% !TeX program = pdflatex
% Automatically machine-translated English edition; the Chinese original is authoritative.
% Original: docs/G_Framework/tex2/G_Framework_Lawspace_Heterogeneous_Governance_ZH.tex
% Original SHA256: 86095955cbf83be59ab96c9a279c7e229da59a7bbf8b183e9a99d934b97d36a1
% Compile from the repository root: .\compile_notebook_tex.ps1 docs\G_Framework\tex4\G_Framework_Lawspace_Heterogeneous_Governance_EN.tex
\documentclass[12pt,a4paper,openany,notitlepage]{report}
\usepackage[utf8]{inputenc}
\usepackage[T1]{fontenc}
\usepackage{CJKutf8}
\usepackage[scaled=0.84]{beramono}
\usepackage{indentfirst,titlesec}
\setlength{\parindent}{0pt}
\renewcommand{\contentsname}{Contents}
\renewcommand{\abstractname}{Abstract}
\renewcommand{\bibname}{References and Source Locations}
\renewcommand{\tablename}{Table}
\renewcommand{\figurename}{Figure}
\usepackage[a4paper,left=20mm,right=20mm,top=25mm,bottom=25mm,headheight=15pt]{geometry}
\usepackage{amsmath,amssymb,amsthm,mathtools,bm,mathrsfs}
\usepackage{booktabs,longtable,array,tabularx,multirow}
\usepackage{enumitem,setspace,fancyhdr,xcolor}
\usepackage{xurl}
\usepackage[unicode=true,colorlinks=true,linkcolor=blue,citecolor=blue,urlcolor=blue,bookmarksnumbered=true,bookmarksdepth=2,linktoc=all]{hyperref}
\Urlmuskip=0mu plus 1mu
\def\UrlFont{\ttfamily\small}
\DeclareRobustCommand{\gframeworkunderscore}{\raisebox{-0.35ex}{\rule{0.55em}{0.35pt}}}
\makeatletter
\g@addto@macro\UrlSpecials{\do\_{\gframeworkunderscore}}
\AddToHook{file/c70gkai.fd/after}{%
  \DeclareFontShape{C70}{gkai}{m}{it}{<->CJK*gkaiu}{\CJKnormal}%
  \DeclareFontShape{C70}{gkai}{bx}{it}{<->CJKb*gkaiu}{\CJKbold}%
}
\let\gframeworkCJKpunctsymbol\CJKpunctsymbol
\renewcommand{\CJKpunctsymbol}[1]{%
  \gframeworkCJKpunctsymbol{#1}%
  \ifnum#1="01\relax
    \ifnum"\CJK@plane="30\relax\allowbreak\fi
  \fi}
\ifdefined\pdfnobuiltintounicode
  \let\gframeworkCJKaddcmap\CJK@addcmap
  \def\CJK@addcmap#1{%
    \ifx\f@encoding\CJK@enc
      \gframeworkCJKaddcmap{#1}%
      \pdfnobuiltintounicode\font@name
    \fi}
\fi
\makeatother

\hypersetup{pdftitle={Heterogeneous Evolutionary Governance in GaoZheng Law-Space},pdfauthor={GaoZheng},pdfsubject={Machine-translated English edition; the Chinese original is authoritative},pdfkeywords={G-Framework,law-space,homotopy repair,task-sufficient quotient,semantic calculus},pdflang={en-US}}
\setstretch{1.10}
\widowpenalty=10000
\clubpenalty=10000
\displaywidowpenalty=10000
\setlength{\parskip}{0.6em}
\setlength{\emergencystretch}{8em}
\setlength{\tabcolsep}{4.5pt}
\renewcommand{\arraystretch}{1.20}
\setlist{nosep,leftmargin=2em,topsep=0.35em}
\allowdisplaybreaks[2]
\titleformat{\section}{\Large\bfseries}{\thesection}{0.8em}{}
\titleformat{\subsection}{\large\bfseries}{\thesubsection}{0.8em}{}
\titlespacing*{\section}{0pt}{2.0ex plus .5ex}{1.0ex plus .2ex}
\titlespacing*{\subsection}{0pt}{1.4ex plus .3ex}{.8ex plus .2ex}
\theoremstyle{plain}
\newtheorem{theorem}{Theorem}[chapter]
\newtheorem{lemma}[theorem]{Lemma}
\newtheorem{proposition}[theorem]{Proposition}
\newtheorem{corollary}[theorem]{Corollary}
\theoremstyle{definition}
\newtheorem{definition}[theorem]{Definition}
\newtheorem{assumption}[theorem]{Assumption}
\newtheorem*{certificate}{Certificate}
\theoremstyle{remark}
\newtheorem{remark}[theorem]{Remark}
\renewcommand{\proofname}{Proof}
\newcolumntype{Y}{>{\raggedright\arraybackslash}X}
\newcolumntype{P}[1]{>{\raggedright\arraybackslash}p{#1}}
\pagestyle{fancy}
\fancyhf{}
\fancyhead[L]{\small Law-Space: Heterogeneous Governance}
\fancyhead[R]{\small GaoZheng}
\fancyfoot[L]{\scriptsize Machine translation; Chinese original prevails}
\fancyfoot[C]{\small\thepage}
\renewcommand{\headrulewidth}{0.3pt}
\renewcommand{\footrulewidth}{0pt}
\title{\bfseries Heterogeneous Evolutionary Governance in GaoZheng Law-Space\\[0.65em]\large Semantic Differential Atoms, Geometric Transport, Homological Repair, and Finite Governance}
\author{GaoZheng}
\date{4 November 2026}

\usepackage{tocloft}
\renewcommand{\cftchapleader}{\cftdotfill{\cftdotsep}}
\renewcommand{\chaptername}{Chapter}
\titleformat{\chapter}[display]{\normalfont\huge\bfseries}{\chaptername~\thechapter}{12pt}{\Huge}
\titlespacing*{\chapter}{0pt}{0pt}{24pt}
\setcounter{secnumdepth}{3}
\setcounter{tocdepth}{2}
\pdfstringdefDisableCommands{%
  \def\mathcal#1{#1}\def\mathsf#1{#1}\def\mathbb#1{#1}%
  \def\infty{infinity}\def\Gamma{Gamma}\def\Phi{Phi}%
}

\begin{document}
\begin{CJK*}
{UTF8} {gkai} \pagenumbering {roman} \begingroup \setstretch {1.14} \setlength {\parskip} {0.20em} \setlength {\parindent} {2em} \maketitle
\begin{center}\small\bfseries Machine-translated English edition\end{center}
\noindent\textbf{Translation notice.} This English text was automatically machine-translated from the Chinese original. The Chinese original is authoritative. In case of any discrepancy, ambiguity, or translation error, the Chinese original prevails. The mathematical notation, labels, and citation keys are retained from the original. Machine translation date: 12 November 2026.\par
\noindent Original Chinese manuscript: \path{docs/G_Framework/tex2/G_Framework_Lawspace_Heterogeneous_Governance_ZH.tex}.\par\medskip
 \thispagestyle {empty} \phantomsection \addcontentsline {toc} {chapter} { Summary }\begin{abstract}
This book investigates how explicit mathematical interfaces preserve task identity, identify obstructions to continuation, construct admissible repairs, and retain verified results as reusable capabilities when object types, representations, admissible actions, and laws change. GaoZheng Law-Space and the rich structure of GZ-GMS provide a common mathematical substrate; each state, function, chain, cochain, operation, and certificate has a specified type, domain, and version. Semantic differentials start from the antecedent, successor, and observation difference of a single admissible action. Boundaries of finite action graphs and dual cochains turn these data into logical inferences and execution records that can be recomputed item by item.

The geometric part starts from the metric and negative curvature of a saddle surface, proving growth of transverse Jacobi fields, lower error bounds for finite observation differences, and compatibility conditions for connections, parallel transport, and heterogeneous bundle maps. Smooth singular chains, integration maps, flat subbundles, and characteristic curvature forms provide explicit interfaces from geometric data to cochain complexes. The algebraic part constructs continuation defects in Hom complexes and provides certificates for homotopy fillings, order-by-order compatibility, and unsolvability of finite matrix systems. Kernel-image quotients, compatible extensions, and executable prefixes separately establish the minimum dimension of additional directions and the admissibility of actual repair.

In a complete parameterized family, hyperbolic geometric inputs, nonzero homology classes, actual unit actions, and full higher-order operations share one registered state. Geometric parameters determine the binary operation, while actual progress determines the ternary filling. An explicit linear map relates the linear residual to the total square of the coderivation, and both vanish along the same finite path. The construction also establishes unreachability under the original action family, reachability under the new family, the minimum number of additional directions, and lower bounds on the number of unit actions. Finally, task normal forms and finite-state refinements provide exact readback, equivalence of future observations, and strategy reuse. Migration of law objects uses target inclusion, commutation on states and chains, transport of atomic truth values, version freezing, and unions of premises to prove that arbitrary finite governance updates and their actual execution preserve the corresponding task conclusions.\end{abstract}


\noindent \textbf { Keywords: } GaoZheng Law-Space; GZ-GMS; semantic differential; negative curvature; fiber bundle; Hom complex; homological obstruction; homotopy algebra; compatible expansion; task normal form; certificate migration \par \endgroup \clearpage \tableofcontents \clearpage \chapter* { Reading and Verification Agreement } \addcontentsline {toc} {chapter} { Reading and Verification Conventions } This book adopts the organization method of "define or clarify the premise, mathematical results, certificate, proof, and remarks". The certificate records the input, source and target, symbols, computation, execution and version information that make the conclusion true; the proof gives the symbolic derivation from these data to the conclusion. Finite execution, partition refinement, order-by-order filling and governance migration respectively use their actual induction variables, and save the base case, induction hypothesis, induction step and all unresolved premises.

The following chapter dependency tables are generated from text tags, and the numbers are automatically updated with the document structure.\begin{center}
\small\begin{tabularx}
{\textwidth}{P{.13\textwidth}P{.32\textwidth}Y}\toprule
Chapter&
core object&
Output interface\\
\midrule
\ref{exp:ch:logic}&
Typed languages and atomic certificates&
Substitution, finite composition, and inductive soundness\\
\ref{exp:ch:law}&
law object and task morphism&
Transport of status, target, observation and cost\\
\ref{exp:ch:higher}&
Strict algebra and higher-order operations&
Non-zero defects, complete compatibility and parametric equations\\
\ref{exp:ch:task}&
Execution path and task chain&
Boundaries, semantic differential and dynamic quotient\\
\ref{exp:ch:geometry}&
Negative curvature and variation&
finite observation difference and error certificate\\
\ref{exp:ch:transport}&
fiber bundle and connection&
Heterogeneous transportation, points and cochain mapping\\
\ref{exp:ch:hom}&
Hom complex and continuation surface&
Defect classes, filling and higher-order continuations\\
\ref{exp:ch:extension}&
Constrained Repair Complex&
Minimum direction, feasibility and practical implementation\\
\ref{exp:ch:tower}&
Parameters and arity filtration&
Stepwise barriers with homotopy algebra closure\\
\ref{exp:ch:joint}&
Complete joint parameter family&
Strictly reachable expansion on the same task\\
\ref{exp:ch:normalform}&
Normal form and finite task quotient&
Decoding, stable partitioning and capability reuse\\
\ref{exp:ch:governance}&
law object and certificate migration&
Target acceptance, version and prerequisite maintenance\\
\bottomrule
\end{tabularx}
\end{center}


The symbol's source and target are fixed at the point where they are first defined. The chain complex adopts the boundary of the descending order $\partial$, and the cochain complex adopts the differential of the ascending order $d$; the symbols of the fractional Hom differential and suspension are completely defined separately. The geometry substrate, task execution chain, and repair change complex have their own object types, and the usage relationship between them is given by the actual chain map, cochain map, or explicitly listed bridging equations. The total cost retains a breakdown of actual usage such as input, search, execution, validation, encoding, migration, and readback.

Finite rational matrices and finite transfer tables use exact arithmetic; geometric observations use differentiability, domain, error bound and corresponding measurement protocols. The conditional conclusions of this book are explained in the task field and admissible actions field registered by the certificate. The basic sources are the attached published manuscripts and reviews, and the listed published mathematical literature \cite{S01,S10,S19,S38,GFrameworkSynthesis,ExtHatcher,ExtWeibel,ExtLadaMarkl,ExtLeeRiemann,ExtKobayashiNomizu,ExtBottTu}. \clearpage \pagenumbering {arabic} \chapter{Typed Formal Language, Semantic Differential Atoms, and Certificate Derivation}
\label{exp:ch:logic}


\label{hgov:chapter}

This book studies the problem of task preservation when objects, representations, patterns, and allowed operations co-vary. The basic objects adopt GaoZheng Law-Space, GZ-GMS rich structure, task-sufficient quotient and semantic calculus in \cite{GFrameworkSynthesis,S01,S38}; chain complex, chain homotopy and higher-order compatibility adopt the typed structure of \cite{ExtHatcher,ExtWeibel,ExtStacks,ExtLadaMarkl}. All conclusions are established within the clear object domain, action domain, observation family and cost specification. The common interface constructed in this article is \begin{equation}
\label{hgov:eq:spine}
\begin{gathered}
\text{\shortstack[c]{\scriptsize Law objects and\\[-1pt]\scriptsize task\\[-1pt]\scriptsize presentations}}\longrightarrow
\text{\shortstack[c]{\scriptsize Typed transport\\[-1pt]\scriptsize and semantic\\[-1pt]\scriptsize differentials}}\longrightarrow
\text{\shortstack[c]{\scriptsize Closed defects\\[-1pt]\scriptsize and their\\[-1pt]\scriptsize homology classes}},\\
\text{\shortstack[c]{\scriptsize Admissible\\[-1pt]\scriptsize fillings or\\[-1pt]\scriptsize structural\\[-1pt]\scriptsize extensions}}\longrightarrow
\text{\shortstack[c]{\scriptsize Admissible\\[-1pt]\scriptsize finite execution}}\longrightarrow
\text{\shortstack[c]{\scriptsize Task normal\\[-1pt]\scriptsize forms and\\[-1pt]\scriptsize subsequent reuse}}.
\end{gathered}
\end{equation}. Geometric variation, operation compatibility, task identity, execution reachability and regular readback are each connected by their own mapping and proof. For the fixed operation family and its authenticated extensions, this book gives reachable set inclusion, strict inclusion witness, minimum number of new directions, and joint constructions of geometric input, non-zero obstacles, unit repair and finite termination in the same parameter family.


\label{hgov:sec:atoms}

\section { Object language and finite certificate }

\begin{definition}
[Object languages, interpretations and finite certificates]\label{hgov:def:language}
Fixed feature zero domain $\mathbb K$ , typed symbol table $\Sigma$ and declaration model $\mathfrak M$ . The finite linear calculation is taken as $\mathbb K=\mathbb Q$; the geometric and analytical conclusions are taken as $\mathbb R$ and differentiability, domain and comparison conditions are retained. Object types include sets, graded spaces, linear maps, partial actions, real observations, chains, cochains, homotopies, and certificates. The judgment is written as \[
\Gamma\vdash e:\tau,\qquad \Gamma\vdash P(e),\qquad
\mathfrak M\models\Gamma\Rightarrow P(e).
\] certificate \[
\mathcal C=(\Sigma,\Gamma,P,w,\mathsf{Parents},\mathsf{Scope},\mathsf{Version})
\] to save the actual input, derivation witness $w$, parent dependency, valid domain and version. The root node is given by a verified input or an explicit premise; all internal rules are accompanied by a proof of their reliability in $\mathfrak M$.\end{definition}


\section { semantic differential atoms and registration rules }

\begin{definition}
[Semantic differential atoms relative to registration primitives]\label{hgov:def:atom}
On a finite directed action graph, an edge $e:x\to y$ registers action $a$, precondition $P_a(x)$, actual successor $y=F_a(x)$ and cost $c(e)$. Observing the scalar $f$, the atomic decomposition of a semantic differential is \[
\begin{aligned}
\alpha_0&:\operatorname{Type}(x)\land\operatorname{Type}(a),\\
\alpha_1&:P_a(x),\qquad
\alpha_2:y=F_a(x),\\
\alpha_3&:u=f(x),\qquad
\alpha_4:v=f(y),\\
\alpha_5&:\delta=v-u,\qquad
\alpha_6:\langle df,[e]\rangle=\delta .
\end{aligned}
\]. Each atom holds the input value, output value, type and parent sequence number. The atomic granularity here is specified by registration primitives; matrix multiplications are further decomposed into single coefficient products with finite sums, and bound sum pairings are further decomposed into generator reads. Vector observations use the same construction coordinate by coordinate.\end{definition}


\begin{definition} [Basic rules of the finite inference kernel]
\label{hgov:def:kernel}
The rule kernel includes typed substitution, finite arithmetic, pairing of matrix coefficients, linear extension from a specified boundary, propositional conjunction and elimination, instantiation of sound implications, and mathematical induction with a base case and induction step. The following input-output relations specify its executable part:
\begin{center} \small
\begin{tabularx} {\textwidth} {P{.12\textwidth}YY}
\toprule
Rule & Input atoms & Output atoms \\
\midrule
R0 & Objects and signatures & Matching domains, codomains, and degrees \\
R1 & $a_{ij},b_{jk}$ and index ranges & Itemwise records of $\sum_j a_{ij}b_{jk}$ \\
R2 & $\partial e=t(e)-s(e)$ and observation values & $\langle df,e\rangle=f(t(e))-f(s(e))$ \\
R3 & Two linear expressions on generators & Equality on the whole space and a proof of linear extension \\
R4 & Bases of the kernel and image, and elimination records & An actual filling or a separating functional \\
R5 & Quotient fibers, observations, and action tables & Factorization and commutation with successors \\
R6 & Preconditions, actions, and successors & One-step execution and resource updates \\
R7 & Parent certificates, rule instances, and open premises & Child conclusions and the union of premises \\
R8 & Base case, induction predicate, and induction step & Conditional conclusions for all finite prefixes \\
R9 & Encoding, decoding, and task observations & Readback equalities and terminal acceptance \\
\bottomrule
\end{tabularx}
\end{center}
Each rule is applied only to registered inputs; undischarged premises remain in the output judgment.
\end{definition}

\section { inductive soundness of finite certificate graph }

\begin{theorem}
[Conditional reliability of finite atomic derivation]\label{hgov:thm:atomsound}
Let the certificate graph be finite, cycle-free, all inputs have types, and every actually used rule holds a truth value in the declarative model. Then each node that passes the verification is established under the combination of its open premises. This property is derived from input and rule construction.\end{theorem}
\begin{certificate}
[Mathematical induction certificate for deduction prefix]Register the topological order $v_1,\ldots,v_N$, node predicate $P_i$, parent set $J_i\subseteq\{1,\ldots,i-1\}$, open premise $\Gamma_i\supseteq\bigcup_{j\in J_i}\Gamma_j$ and the reliability formula of the rule; when implementing hypothesis release, the reliable release rule and its scope are also saved. A rule without a release node is written as \[
\mathfrak M\models
\left(\Gamma_i\land\bigwedge_{j\in J_i}P_j\right)\Rightarrow P_i.
\]. The validator records the rule instance, input type, and prerequisite pass on a node-by-node basis.\end{certificate}
\begin{proof}
Let $I(n)$ be the conditional conclusion of the first $n$ nodes. The input nodes of the base case are established by its verification or retained premises. If $I(n)$ is true, all the parent nodes of the $n+1$ node are before; substituting their proven judgments into reliable rules, we get that the node is true when the parent premise and the local premise are combined. Therefore the concrete predicate chain \[
I(0)\ \land\
\bigl(\forall n<N,\ I(n)\Rightarrow I(n+1)\bigr)
\Rightarrow I(N).
\] is \[
\begin{gathered}
\operatorname{TypedInputs}\land\operatorname{CheckedParents}
\land\operatorname{SoundRule}\\
\Rightarrow\operatorname{ValidSubstitution}
\Rightarrow(\Gamma_i\Rightarrow P_i)
\Rightarrow I(i).
\end{gathered}
\] finite and induction completes the proof. Error premises for numerical readings are propagated with the same dependency graph.\end{proof}


\section { Explanation, replacement and finite composition of atomic judgment }

\begin{definition}
[Signature examples and judging validity]\label{exp:logic:signature}
A registered function symbol $f$ has the signature $f:\tau_1\times\cdots\times\tau_r\rightharpoonup\tau$, the domain predicate $\mathsf D_f$, and the specified interpretation $f^{\mathfrak M}$. The evaluation of the term $f(t_1,\ldots,t_r)$ submits both $t_i:\tau_i$ and $\mathsf D_f(t_1,\ldots,t_r)$. The model validity of an open judgment $\Gamma\vdash P$ is defined as \[
 \mathfrak M\models\Gamma\vdash P
 \quad\Longleftrightarrow\quad
 \forall\nu\in\mathsf{Val}(\Gamma),\quad
 \left(\bigwedge_{A\in\Gamma}\lbrack\!\lbrack A\rbrack\!\rbrack_\nu=1\right)
 \Rightarrow\lbrack\!\lbrack P\rbrack\!\rbrack_\nu=1.
\] where $\nu$ is an assignment that satisfies type constraints; domain antecedents involving some functions are retained in $\Gamma$. Equation atoms are explained according to the equality of evaluation on both sides, linear relationships are explained according to operations in a given domain, and action atoms are explained according to the actual registered preconditions and successors.\end{definition}


\begin{lemma}
[Type substitution and evaluation exchange]\label{exp:logic:substitution}
Let $\sigma$ replace the variables with terms of the same type, and each used partial function domain condition is satisfied. The induced assignment is defined as $(\nu\circ\sigma)(x)=\lbrack\!\lbrack\sigma(x)\rbrack\!\rbrack_\nu$. Then the convention of avoiding variable capture is adopted for the substitution under any finite term $t$ and finite formula $P$, \[
 \lbrack\!\lbrack t\sigma\rbrack\!\rbrack_\nu
 =\lbrack\!\lbrack t\rbrack\!\rbrack_{\nu\circ\sigma},\qquad
 \lbrack\!\lbrack P\sigma\rbrack\!\rbrack_\nu
 =\lbrack\!\lbrack P\rbrack\!\rbrack_{\nu\circ\sigma}.
\] quantifier.\end{lemma}
\begin{certificate}
[Syntax construction tree induction]Saves a finite construction tree of terms and formulas, the source and replacement types of each variable, local domain witnesses, and a record of the renames of bound variables.\end{certificate}
\begin{proof}
The variable case is established by the definition of $\nu\circ\sigma$, and the constant case is established by the fixed interpretation. If $t=f(t_1,\ldots,t_r)$, the induction hypothesis gives consistent evaluation of each parameter, so the \[
 \begin{aligned}
 \lbrack\!\lbrack t\sigma\rbrack\!\rbrack_\nu
 &=f^{\mathfrak M}(\lbrack\!\lbrack t_1\sigma\rbrack\!\rbrack_\nu,
                    \ldots,\lbrack\!\lbrack t_r\sigma\rbrack\!\rbrack_\nu)\\
 &=f^{\mathfrak M}(\lbrack\!\lbrack t_1\rbrack\!\rbrack_{\nu\circ\sigma},
                    \ldots,\lbrack\!\lbrack t_r\rbrack\!\rbrack_{\nu\circ\sigma})
 =\lbrack\!\lbrack t\rbrack\!\rbrack_{\nu\circ\sigma}.
 \end{aligned}
\] domain witness guarantees that both sides are defined at the same time. Equality and relational atoms are established by the conclusion of the terms; conjunctions, negations, and implications are maintained by the same truth operations. For $\forall x:\tau\,Q$, first separate $x$ from the free variables in substitution, and then use the inductive conclusion of $Q$ for each $a$ in the same type domain to obtain the quantifier conclusion; the same is true for existing quantifiers. So \[
 \begin{gathered}
 \operatorname{TypedSubstitution}\land\operatorname{DomainWitness}\\
 \Rightarrow\operatorname{TermEvaluationCommutes}
 \Rightarrow\operatorname{FormulaEvaluationCommutes}.
 \end{gathered}
\]\end{proof}


\begin{theorem}
[Prerequisite union and dependency preservation of certificate compound]\label{exp:logic:certificate-composition}
Suppose finite certificate $\mathcal C_1$ proves $\Gamma_1\vdash P$ and certificate $\mathcal C_2$ proves $\Gamma_2\cup\{P\}\vdash Q$. If the shared objects, types, domains and versions of the two certificates are consistent, then there is a finite composite certificate proving \[
 \Gamma_1\cup\Gamma_2\vdash Q.
\]. Any finite certificates that meet these continuation conditions can be composited according to the same rules.\end{theorem}
\begin{certificate}
[Conclusion node replacement and topological order]Bind the leaf node in $\mathcal C_2$ using the hypothesis $P$ to the conclusion node of $\mathcal C_1$; renumber all nodes, first list $\mathcal C_1$, and then list the rest of the topological order of $\mathcal C_2$. Open premises are recorded as unions, and duplicate premises can share the same record.\end{certificate}
\begin{proof}
Assume that $\Gamma_1\cup\Gamma_2$ is all established. The reliability of $\mathcal C_1$ gives $P$; all inputs of $\mathcal C_2$ are therefore established, and $Q$ is obtained. The above node replacement only adds dependent edges from the first part to the later part, maintaining acyclicity. Each compound continues to satisfy the conditions of theorem ~ \ref{hgov:thm:atomsound}. Let $\Gamma^{(n)}$ be the open premise union of the first $n$ certificates, then \[
 \begin{gathered}
 \Gamma^{(n+1)}=\Gamma^{(n)}\cup\Gamma_{n+1},\\
 \operatorname{ValidComposite}(n)\land\operatorname{MatchedInterface}(n+1)
 \Rightarrow\operatorname{ValidComposite}(n+1).
 \end{gathered}
\] zero certificates correspond to an empty dependency graph; the conclusion of any finite composition can be obtained by induction according to the number of compositions.\end{proof}


\section { Value-wise expansion and recalculable calculation example of semantic differential }

\begin{proposition}
[Atomic differential accumulation on finite paths]\label{exp:logic:telescoping}
For admissible path $x_0\xrightarrow{a_1}x_1\cdots\xrightarrow{a_n}x_n$, assume that observation $f$ is defined in all access states. Register $u_i=f(x_i)$ and $\delta_i=u_i-u_{i-1}$, then \[
 \sum_{i=1}^{n}\delta_i=f(x_n)-f(x_0).
\] vector observation satisfies the same equation according to coordinates.\end{proposition}
\begin{certificate}
[Side-by-side numerical and partial sums]For each step, save the type, precondition, successor, two observation values, difference value and cumulative value $s_0=0$, $s_k=s_{k-1}+\delta_k$.\end{certificate}
\begin{proof}
Base case $s_0=u_0-u_0=0$. If $s_k=u_k-u_0$ , then \[
 s_{k+1}=s_k+(u_{k+1}-u_k)
 =(u_k-u_0)+(u_{k+1}-u_k)=u_{k+1}-u_0.
\] mathematical induction gives all finite prefixes. Applying the same proof to each coordinate of the vector gives us the vector's expression. Its atomic predicate chain is \[
 y=F_a(x)\land u=f(x)\land v=f(y)
 \Rightarrow\delta=v-u
 \Rightarrow s_{k+1}=s_k+\delta
 \Rightarrow s_{k+1}=f(x_{k+1})-f(x_0).
\]\end{proof}


Take $X=\mathbb Q^2$, actions $F_a(x,y)=(x+1,y)$ and $F_b(x,y)=(x,y+1)$, and observation $f(x,y)=x+2y$. Starting at $(0,0)$ and executing the action word $aba$ gives the complete arithmetic record
\begin{center}
\begin{tabular} {ccccc}
\toprule
$k$ & Actual state $x_k$ & $u_k=f(x_k)$ & $\delta_k$ & $s_k$ \\
\midrule
0 & $(0,0)$ & 0 & --- & 0 \\
1 & $(1,0)$ & 1 & 1 & 1 \\
2 & $(1,1)$ & 3 & 2 & 3 \\
3 & $(2,1)$ & 4 & 1 & 4 \\
\bottomrule
\end{tabular}
\end{center}
Each entry can be recomputed using only rational addition, multiplication, and domain lookup. The terminal difference $4-0=4$ equals the accumulated difference $1+2+1=4$. The action word $aba$, visited states, and stepwise costs remain stored with the certificate, giving difference summation and actual execution separate, explicit readback objects.

\begin{remark}
[Proof status of structural conditions]In this book, boundary closure, compatible transport, repair reachability and governance preservation are all given by conditional theorems. Constructive inputs are registered through matrices, operation tables, maps and their domains; each property that needs to be maintained is proven by the corresponding certificate. The explanation of basic logic and the explicit use of domain operations constitute the derivation environment, and the structural validity of the model is established by symbolic calculations and finite induction within this environment. The named predicates in the proof chain are unified as abbreviations of previously defined conditions or proved equations; certificate also records its expansion, actual parameters and derivation rules used.\end{remark}


\chapter{GaoZheng Law-Space, Generalized Mathematical Structure, and Task Substrates}
\label{exp:ch:law}


\label{hgov:sec:law}

\section { Objects, laws and common mathematical substrate }

\begin{definition}
[law object with tasks]\label{hgov:def:law}
Use $\mathsf{Law}$ as a set of typed rules to present. Each $L\in\mathsf{Law}$ registers \[
\mathfrak S_L=(X_L,A_L,F_L,O_L,I_L,G_L,\mathsf{Cost}_L),
\], where $F_L$ is the partial successor with precondition, $O_L$ is the task observation, $I_L$ is the task identity, and $G_L$ is the target predicate. The common mathematical substrate takes the \[
\mathcal E=\coprod_{L\in\mathsf{Law}}\{L\}\times X_L,
\qquad \pi(L,x)=L.
\] structure mapping and simultaneously registers state mapping, action correspondence, observation transportation and identity maintenance. The data organization of GZ-GMS is \[
\mathcal G=(\mathsf{Obj},\mathsf{Law},\mathsf{Flow},
\mathsf{Layer},\mathsf{GZTLC};\mathsf{LawCurv}).
\]. Each data has actual type and compatibility relationship; the smooth sectors of fixed typical fibers are configured with ordinary plexus, and the migration configuration of cross-type sectors has type structure mapping. The above substrate continuation exposes law-space with rich structure definition \cite{S01,S38} .\end{definition}


\section { Fixed regular slicing and recovery of classic tasks }

\begin{proposition}
[Fixed regularity embedding and observation projection]\label{hgov:prop:slice}
Fixed $L_0$, mapping $i_{L_0}:X_{L_0}\to\mathcal E$, $i_{L_0}(x)=(L_0,x)$ as injective; its image on-projection $p_{L_0}(L_0,x)=x$ satisfies $p_{L_0}i_{L_0}=\mathrm{id}$. If the extended action is defined as $\widehat F_a(L_0,x)=(L_0,F_{L_0,a}x)$ on this image, then \[
\widehat F_ai_{L_0}=i_{L_0}F_{L_0,a}.
\] stipulates that when the observation, identity, and target are pulled back according to this projection, the task conclusions of all finite action words are maintained.\end{proposition}
\begin{certificate}
[Embed certificate field by field]Save $L_0$, state domain, part of the action domain, projection and the pullback of each observation; save the true value of the action precondition before and after embedding.\end{certificate}
\begin{proof}
The ordered pairs of equals give $i_{L_0}$ injective , which can be obtained by substituting into the projection and one-step commutation relation . Induction for word length $n$: zero length is identity; one-step exchange and induction hypothesis are combined to obtain \[
\widehat F_{a_{n+1}}\widehat F_{a_n}\cdots\widehat F_{a_1}i_{L_0}
=i_{L_0}F_{a_{n+1}}F_{a_n}\cdots F_{a_1}.
\]. The domain of each prefix is maintained uniformly by precondition, and the observation is maintained by pullback. Therefore \[
\operatorname{Embedding}\land\operatorname{DomainMatch}
\land\operatorname{StepCommute}
\Rightarrow\operatorname{WordCommute}
\Rightarrow\operatorname{TaskPreserved}.\qedhere
\]\end{proof}


\section { law object morphism, finite execution simulation and composition }

\begin{definition}
[Maintain the task's law object morphism]\label{exp:law:morphism}
For law object $\mathfrak S_L,\mathfrak S_{L'}$, the common identity domain $\mathcal I$ is specified, allowing different observation value domains. A task morphism is recorded as $\Phi=(f,\tau,d):\mathfrak S_L\to\mathfrak S_{L'}$, where $f:X_L\to X_{L'}$ is the state mapping, $\tau:A_L\to A_{L'}$ is the label mapping, $d:Y_{L'}\to Y_L$ is the observation decoder, and satisfies \[
 \begin{aligned}
 \operatorname{Pre}_{L,a}(x)
 &\Longleftrightarrow\operatorname{Pre}_{L',\tau(a)}(f x),\\
 fF_{L,a}(x)&=F_{L',\tau(a)}(fx)
 &&\text{\shortstack[c]{\scriptsize on the\\[-1pt]\scriptsize admissible\\[-1pt]\scriptsize domain}},\\
 I_{L'}f&=I_L,& O_L&=dO_{L'}f,\\
 f(G_L)&\subseteq G_{L'},&
 c_{L,a}(x)&=c_{L',\tau(a)}(fx).
 \end{aligned}
\]. The cost adopts the common metering domain. Action families can contain new tags; the above requirements apply to tags registered on the source side one by one.\end{definition}


\begin{theorem}
[Path simulation and cumulative cost preservation of law object morphisms]\label{exp:law:simulation}
If $\Phi$ satisfies the definition ~ \ref{exp:law:morphism} , then each source-side legal action word $w=a_n\cdots a_1$ has a corresponding word $\tau(w)$ on the image state, and the \[
 fF_{L,w}(x)=F_{L',\tau(w)}(fx),\qquad
 \sum_{k=1}^n c_{L,a_k}(x_{k-1})
 =\sum_{k=1}^n c_{L',\tau(a_k)}(f x_{k-1}).
\] identity, prescribed observation, and successful target run along that correspondence are maintained.\end{theorem}
\begin{certificate}
[Synchronization prefixes for status, word and fee]Save each triple $(x_k,fx_k,\tau(a_k))$, action predecessor equivalence, successor exchange, cost equality, and end target containing instance.\end{certificate}
\begin{proof}
When there are zero words, the status exchange is $f\mathrm{id}=\mathrm{id}f$, and the accumulated costs are all zero. If the previous $k$ step correspondence is established, the precondition equivalence guarantees that step $k+1$ is defined at the same time, and \[
 F_{L',\tau(a_{k+1})}(fx_k)=fF_{L,a_{k+1}}(x_k)=fx_{k+1}.
\] adds an equal $k+1$ step fee to the two equal accumulated fees to obtain a new equality. Therefore, induction holds for word length. Then substitute $I_{L'}f=I_L$, $O_L=dO_{L'}f$ and $f(G_L)\subseteq G_{L'}$ to get the task conclusion. \begin{gather*}
 \operatorname{DomainMatch}\land\operatorname{StepCommute}\land\operatorname{StepCostEqual}\\
 \Rightarrow\operatorname{PrefixSimulation}
 \Rightarrow\operatorname{SuccessfulWordPreserved}.\qedhere
\end{gather*}\end{proof}


\begin{proposition}
[Finite composition of task morphisms]\label{exp:law:composition}
Let $\Phi=(f,\tau,d):\mathfrak S_0\to\mathfrak S_1$ and $\Psi=(g,\upsilon,e):\mathfrak S_1\to\mathfrak S_2$ be task morphisms. Then \[
 \Psi\circ\Phi=(gf,\upsilon\tau,de)
\] is still a task morphism. The identity map gives the identity morphism, and the composition satisfies the associative law.\end{proposition}
\begin{certificate}
[Field-by-field compounding]Register the source target, common field, six sets of holding formulas and composite mapping of two morphisms. The order of observation decoders and the order of state transport are saved separately.\end{certificate}
\begin{proof}
precondition by two equivalent continuations. The successor and observation respectively satisfy \[
 \begin{aligned}
 gfF_{0,a}&=gF_{1,\tau(a)}f
 =F_{2,\upsilon\tau(a)}gf,\\
 O_0&=dO_1f=deO_2gf.
 \end{aligned}
\]. The identity is maintained by $I_2gf=I_1f=I_0$; the target is maintained by $gf(G_0)\subseteq g(G_1)\subseteq G_2$; the cost is substituted into the equality equation twice value by value. Identity mapping satisfies the definition item by item, and the associativity of function and label mapping gives the associativity of all fields. Any finite composition is established by induction from the degree of composition.\end{proof}


\section { Encoding and geometric sectors of typing substrate }

\begin{proposition}
[Reversible reading of tagged unions]\label{exp:law:typed-union}
Let each object of $L$ have the injective encoding $e_L:X_L\to Z_L$. Then \[
 E:\coprod_L\{L\}\times X_L\longrightarrow
   \coprod_L\{L\}\times Z_L,
 \qquad E(L,x)=(L,e_L(x))
\] is an injective and has the unique inverse reading $D(L,e_L(x))=(L,x)$ on its image.\end{proposition}
\begin{certificate}
[Type labels and local inverses]Save the label of $L$, the domain of $e_L$, the injective proof, the local inverse reading on the image, and verify that the label of each actual input is consistent with the local code.\end{certificate}
\begin{proof}
If $E(L,x)=E(L',x')$ , the ordering gives $L=L'$ for the first component and $e_L(x)=e_L(x')$ for the second component. From the property of injective, we get $x=x'$, so $E$ is injective. Each item on the image has a unique original image. Define $D$ to get $DE=\mathrm{id}$ and $ED=\mathrm{id}_{\operatorname{im}E}$. \[
 \operatorname{SameTypedCode}
 \Rightarrow L=L'\land e_L(x)=e_L(x')
 \Rightarrow(L,x)=(L',x').\qedhere
\]\end{proof}


\begin{proposition}
[Dimension sectors of local trivial vector bundles]\label{exp:law:rank-sector}
Let $p:E\to B$ be a locally trivial finite-dimensional vector bundle. Then the fiber dimension function $b\mapsto\dim E_b$ is locally constant and therefore constant on each connected branch. For a law object family containing fibers of different dimensions, its local trivial sectors are registered separately according to the dimensions, and cross-sector transportation is specified by an additional typed mapping.\end{proposition}
\begin{certificate}
[Local trivialization and rank]Save coverage $\{U_i\}$ , linear fiber isomorphism $p^{-1}(U_i)\cong U_i\times\mathbb K^{r_i}$ , and reversibility of transition mappings at overlaps.\end{certificate}
\begin{proof}
If $b\in U_i$, local trivialization restricted to the fiber gives $E_b\cong\mathbb K^{r_i}$, so the dimension is always $r_i$ on $U_i$. For any integer $r$, the point set with rank $r$ is the union of these local fields, so it is an open set; its complement consists of other rank fields, which is also an open set. If a connected branch has two different ranks, it can be divided into two non-empty open and closed parts by the corresponding rank domain, which is inconsistent with connectivity. Therefore the rank is constant on each connected branch.\end{proof}


\begin{remark}
[law-space data and its implementation]A common type substrate that labels and provides objects, laws, and implementations. Set encoding, smooth local trivialization, task morphisms, and actual actions each preserve their own structural conditions. A cross-type update can simultaneously include state transport, chain mapping, and geometric implementation; its task conclusion is derived from verified field persistence and finite composition.\end{remark}


\chapter{Strict Algebraic Sectors and Complete Higher-Order Compatibility Structures}
\label{exp:ch:higher}


\section { Zero-extension construction for the strict Lie sector }
\begin{theorem} [Zero-extension construction for the strict Lie sector]
\label{hgov:thm:strictlie}
Suppose $V$ is nonzero only in degree zero, and the binary operation $b:\Lambda^2V\to V$ is antisymmetric. Set $\ell_1=0$, $\ell_2=b$, and $\ell_n=0$ for $n\ge3$. These data define an $L_\infty$ structure if and only if $b$ satisfies the Jacobi identity. Strict Lie homomorphisms preserve this embedding when their higher-arity components are set to zero.
\end{theorem}
\begin{certificate} [Zero extension arity by arity]
Using the convention $\deg\ell_n=2-n$, record the structure constants and antisymmetry of $b$, the ternary Jacobi expression, and the fact that all higher-arity components vanish; for the map, record $fb=b'(f,f)$.
\end{certificate}
\begin{proof}
The arity-one and arity-two relations involve $\ell_1$ and therefore hold; at arity three only the following remains:
\[
b(x,b(y,z))+b(y,b(z,x))+b(z,b(x,y))=0.
\]
When $n\ge4$, the arities in each composite term satisfy $i+j=n+1$. Any term containing a unary operation is zero; in every other case at least one operation has arity at least three, so the term is again zero. Reading the ternary expression in reverse gives the Jacobi identity. The only nonzero part of the map relations is the bracket-preservation equation. Hence
\[
\operatorname{Skew}(b)\land\operatorname{Jacobi}(b)
\Rightarrow\bigwedge_{n\ge1}\mathcal R_n=0,
\qquad
\bigwedge_{n\ge1}\mathcal R_n=0\Rightarrow\operatorname{Jacobi}(b).
\]
The signs in the higher-arity relations and the general definition follow \cite{ExtLadaMarkl}; the restricted construction above is the strict-structure embedding in \cite{S38}.
\end{proof}

\begin{remark}
[Structural Hierarchy and Presentation Interface]Fixed regular slicing, zero-expansion embedding, task quotient and chain homotopy configure their own mappings and holds respectively. All sets, functions, tables, and proof records determine the rich structure used; the compatibility of each sector is determined by specific structural equations.\end{remark}


\section { The complete two-term homotopy structure of the non-zero Jacobiator }\label{hgov:linfty:section}


Zero expansion of strict sectors and presentations carrying higher-order witnesses can be compared in the same system of identities. The two-term model of open document \cite{S38} provides a finite, non-zero, and ternary-verifiable instance: a binary bracket has a non-zero Jacobiator, the ternary operation preserves its filling witness, and all higher-order relations remain closed. The following fixed symbol conventions and complete proofs are given; for the background of general high-order operations, see \cite{ExtLadaMarkl}.

\begin{definition}
[cohomology grading and higher-order compatibility formula]\label{hgov:linfty:convention}
On the field of characteristic zero $\mathbb K$, let $V$ be the graded space, and $\ell_n:\Lambda^n V\to V$ be the graded antisymmetric operation of degree $2-n$. Align the homogeneous input $x_1,\ldots,x_n$, let $\chi(\sigma;x)$ be the product of the permutation symbol and the Koszul symbol produced by exchanging the homogeneous input. Define \[
\begin{aligned}
\mathcal R_n(x_1,\ldots,x_n)
={}&\sum_{i+j=n+1}\ \sum_{\sigma\in\operatorname{Sh}(i,n-i)}
(-1)^{i(j-1)}\chi(\sigma;x)\\
&\qquad\cdot
\ell_j\bigl(\ell_i(x_{\sigma(1)},\ldots,x_{\sigma(i)}),
x_{\sigma(i+1)},\ldots,x_{\sigma(n)}\bigr).
\end{aligned}
\] where $\operatorname{Sh}(i,n-i)$ is a shuffle permutation set that maintains the internal order of the first $i$ positions and the last $n-i$ positions. The complete $L_\infty$ closure predicate is \[
\operatorname{HigherClosed}(V,\ell_\bullet)
\quad\Longleftrightarrow\quad
\forall n\ge1,\quad\mathcal R_n=0.
\] and the zero-degree antisymmetric bracket $b$ stipulates \[
J_b(x,y,z)=b(x,b(y,z))+b(y,b(z,x))+b(z,b(x,y)).
\]. Under this convention, the ternary satisfied by the three $\ell_1$-closed zero-degree inputs is $-J_{\ell_2}+\ell_1\ell_3=0$.\end{definition}


\begin{theorem}
[Complete higher-order rendering with non-zero binary defects]\label{hgov:linfty:model}
Take $E=\mathbb K^3$, the base is $e_1,e_2,e_3$, and define the antisymmetric bilinear operation \[
b(e_1,e_2)=e_2,\qquad b(e_2,e_3)=e_3,\qquad b(e_1,e_3)=0.
\]. Let $\widetilde E$ be a linear copy of $E$, and $x\mapsto\widetilde x$ be the specified isomorphism. Take $V^0=E$ and $V^{-1}=\widetilde E$, and the remaining degrees are zero. Define \[
\begin{aligned}
&\ell_1(\widetilde x)=x,\qquad \ell_1(x)=0,\\
&\ell_2(x,y)=b(x,y),\qquad
\ell_2(x,\widetilde y)=\widetilde{b(x,y)},\qquad
\ell_2(\widetilde x,\widetilde y)=0,\\
&\ell_3(x,y,z)=\widetilde{J_b(x,y,z)},\qquad
\ell_n=0\quad(n\ge4).
\end{aligned}
\]. The remaining binary values are determined by the graded antisymmetry; $\ell_3$ with negative order input is zero. Then these operations have a specified number of times and satisfy $\mathcal R_n=0$ for all $n\ge1$ , and \[
J_b(e_1,e_2,e_3)=-e_3\ne0.
\] so this fixed render has both non-zero low-order defects and full higher-order compatibility witnessed.\end{theorem}


\begin{certificate}
[structure constants, degree and element verification]\label{hgov:linfty:model-certificate}
Save the corresponding basis of $E$ and $\widetilde E$, the matrices $I_3$, $b_{12}^{2}=b_{23}^{3}=1$ of $\ell_1:V^{-1}\to V^0$ and their antisymmetric structure constants. $\ell_3(e_1,e_2,e_3)=-\widetilde e_3$ is saved, and the remaining zero inputs are evaluated alternately. The verification is divided into six groups: $1,2,3,4,5$ and $n\ge6$; each group is divided into negative input numbers. The finite part is expanded basis by basis, and any higher element is uniformly proved by the arithmetic truncation rules.\end{certificate}


\begin{proof}
The bilinear operation $b$ is antisymmetric, so $J_b$ is an alternating trilinear mapping. All definitions of $\ell_n$ satisfy the times $2-n$. For example $\ell_2(\widetilde x,y)=\widetilde{b(x,y)}$ , while the binary output of two negative inputs should be at $V^{-2}=0$ . Check the numbers below.

\begin{enumerate}
\item
Number one. $\ell_1$ feeds $V^{-1}$ into $V^0$ and zeroes in on $V^0$, thus $\mathcal R_1=\ell_1^2=0$ .

\item
Number two. The compatible formula is the fractional derivative law \[
\ell_1\ell_2(a,b)
=\ell_2(\ell_1a,b)+(-1)^{|a|}\ell_2(a,\ell_1b).
\]. Two zero-order inputs make both sides zero; inputting $(x,\widetilde y)$ makes both sides $b(x,y)$. For $(\widetilde x,\widetilde y)$, the left side is zero and the right side is \[
\ell_2(x,\widetilde y)-\ell_2(\widetilde x,y)
=\widetilde{b(x,y)}-\widetilde{b(x,y)}=0.
\]. The graded antisymmetry covers another input order, so $\mathcal R_2=0$.

\item
Number three. Three zero-order inputs give \[
\mathcal R_3(x,y,z)=-J_b(x,y,z)+\ell_1\widetilde{J_b(x,y,z)}=0.
\]. For one negative-order input $(\widetilde x,y,z)$, the binary composite term is $-\widetilde{J_b(x,y,z)}$; the only possible non-zero input differential term is $\ell_3(\ell_1\widetilde x,y,z)=\widetilde{J_b(x,y,z)}$, and the two terms cancel. When there are at least two negative inputs, the output times of each composite term are at most $-2$, so they are all zero. In addition, directly substituting the three basis vectors we get \[
J_b(e_1,e_2,e_3)
=b(e_1,e_3)+b(e_2,0)+b(e_3,e_2)=-e_3.
\]

\item
Number four. The only possible non-zero entries are $\ell_3\ell_2$ and $\ell_2\ell_3$. For zero-order input, its shuffled sum is written as \[
\begin{aligned}
\mathcal R_4(x_1,x_2,x_3,x_4)
={}&\sum_{\sigma\in\operatorname{Sh}(2,2)}
\operatorname{sgn}(\sigma)
\ell_3\bigl(\ell_2(x_{\sigma(1)},x_{\sigma(2)}),
x_{\sigma(3)},x_{\sigma(4)}\bigr)\\
&-\sum_{\sigma\in\operatorname{Sh}(3,1)}
\operatorname{sgn}(\sigma)
\ell_2\bigl(\ell_3(x_{\sigma(1)},x_{\sigma(2)},x_{\sigma(3)}),
x_{\sigma(4)}\bigr).
\end{aligned}
\] This expression is an alternating quadlinear map of $E^4\to\widetilde E$. Since $\dim E=3$ , any four inputs are linearly related, so the expression is zero. When negative input is included, the total number of each compound item is at most $-2$, which is also zero.

\item
Number five. The only possible non-zero combination of operations is $\ell_3\ell_3$ . When the inner layer is non-zero, its output belongs to $V^{-1}$, and the outer layer $\ell_3$ is zero because it contains negative inputs; when the inner layer is zero, the composite is also zero. Hence $\mathcal R_5=0$.

\item
Element number $n\ge6$. Each compound satisfies $i+j=n+1\ge7$, so at least one of $i,j$ is not less than four. The corresponding operation is always zero, so each item of $\mathcal R_n$ is zero.\end{enumerate}
The above situation covers all yuan numbers and all input times. The predicate chain of the proof is \[
\begin{gathered}
\operatorname{GradedTypes}\land
\bigwedge_{n=1}^{5}(\mathcal R_n=0)
\land
\bigl(\forall n\ge6,\ \mathcal R_n=0\bigr)\\
\Rightarrow\operatorname{HigherClosed}(V,\ell_\bullet),\\
J_b(e_1,e_2,e_3)=-e_3\ne0
\Rightarrow\operatorname{NonzeroBinaryDefect}(V,\ell_2).
\end{gathered}
\]. The two conclusions are simultaneously established on the same set of structure constants.\end{proof}


\section { Co-reading of shrinkable witness sectors and task core }

\begin{proposition} [Joint preservation of the contractible witness sector and the task core]
\label{hgov:linfty:task}
Take a task space $T$ concentrated in degree zero, with zero differential and zero higher-arity operations. On $W=T\oplus V$ retain all operations of Theorem \ref{hgov:linfty:model}, and set binary and higher-arity operations involving inputs from $T$ to zero. Let
\[
p:W\to T,\qquad i:T\to W,\qquad D=\ell_1^W,
\]
be the direct-sum projection, inclusion, and total differential, and define a map of degree $-1$ by
\[
h(x)=\widetilde x\ (x\in E),\qquad
h(\widetilde x)=0,\qquad h(T)=0.
\].
Then $W$ is a full $L_\infty$ structure, $p,i$ are strict higher-order structure maps, and
\[
pi=\mathrm{id}_T,\qquad
Dh+hD=\mathrm{id}_W-ip.
\].
Consequently, the underlying complex satisfies $H^\bullet(W)\cong T$, where $T$ lies in degree zero. For any task observation factoring through $p$, denoted by $O=\overline O p$, we have $Oip=O$. If the extended actions additionally satisfy
\[
p\widehat F_a=F_a^Tp,
\]
and action availability is constant on the fibers of $p$, every finite admissible action word preserves this task readback.
\end{proposition}

\begin{certificate}
[Component-wise shrinkage and action swapping]\label{hgov:linfty:task-certificate}
Register $D,h,p,i$ and the contraction equation on $T$ , $E$ , and $\widetilde E$ respectively; each specified observation, identity, or target saves its factorization by $p$ . Finite execution also saves the domain, one-step commutation relation and actual successor of each action. When cost preservation is required, factorization of step-by-step costs is incorporated into the same observation certificate.\end{certificate}


\begin{proof}
The high-order relationship of the pure $V$ input is established by the theorem \ref{hgov:linfty:model}; the high-element composite containing the $T$ input is zero, and the differential is also zero on $T$, so all relationships are maintained. $p$ sends non-zero operations on $V$ to zero, and all higher-order operations in the image of $i$ are zero, so the binary mapping strictly preserves each operation.

On $T$, $Dh+hD$ and $\mathrm{id}_W-ip$ are both zero; on $E$ there is $Dh=\mathrm{id}_E$; on $\widetilde E$ there is $hD=\mathrm{id}_{\widetilde E}$. Adding component-by-component gives the contraction. This formula and $pi=\mathrm{id}_T$ give the homotopy retraction of the underlying complex, so its cohomology is read back by $T$. \[
Oip=\overline O pip=\overline O p=O.
\] is directly obtained from the observation factorization. The action word length is establish by induction, and the one-step commutation relation gives $p\widehat F_w=F_w^Tp$; the availability on the fiber ensures that each prefix is defined at the same time. So \[
\begin{gathered}
\operatorname{Retraction}(p,i,h)\land O=\overline O p
\Rightarrow\operatorname{CoreReadback},\\
\operatorname{DomainMatch}\land\operatorname{StepCommute}
\Rightarrow\operatorname{WordCommute}
\Rightarrow O(\widehat F_w z)=\overline O(F_w^Tpz).
\end{gathered}
\] This gives the component-wise and word-by-word certificate of the public task core direct sum construct \cite{S38}.\end{proof}


\begin{remark}
[Operation rendering and separate reading of underlying homotopy types]\label{hgov:linfty:scope}
This example saves $J_b\ne0$ and its non-zero ternary witness in a fixed arithmetic representation, thereby demonstrating a complete higher-order representation beyond strictly binary sectors. Its underlying complex $V$ is scalable; after adding task core, the underlying homotopy type of $T\oplus V$ is still given by $T$. The addition of presentation fields, the filling of low-order defects, and the maintenance of task homotopy types are expressed by structure constants, full higher-order compatibility, and explicit shrinking certificates, respectively.\end{remark}


\section { Coordinate expansion and parameterization compatibility conditions of structure constants }

\begin{proposition}
[Coordinate expressions of binary operations and ternary defects]\label{exp:higher:coordinates}
For Theorem ~ \ref{hgov:linfty:model}, take $E=\mathbb K^3$ and write $x,y,z$ as a column vector in the specified basis. Then
\[
 b(x,y)=
 \begin{pmatrix}
 0\\x_1y_2-x_2y_1\\x_2y_3-x_3y_2
 \end{pmatrix},
 \qquad
 J_b(x,y,z)=-\det[x\ y\ z]\,e_3.
\].
For $a=x+\widetilde u$, $a'=y+\widetilde v$, and $a''=z+\widetilde w$, all nonzero operations are determined by
\[
 \begin{aligned}
 \ell_1(a)&=u,\\
 \ell_2(a,a')&=b(x,y)+\widetilde{b(x,v)+b(u,y)},\\
 \ell_3(a,a',a'')&=-\det[x\ y\ z]\,\widetilde e_3
 \end{aligned}
\].
These formulas are expanded linearly on each homogeneous component.\end{proposition}
\begin{certificate}
[Primitive product and determinant expansion]Saves six primitives $e_i,\widetilde e_i$, non-zero binary structure constants, antisymmetric reading, homogeneous decomposition of the input, and six calculations of the determinant.\end{certificate}
\begin{proof}
From bilinearity and structure constants, the coefficients of $e_2$ only come from $b(e_1,e_2)=e_2$ and $b(e_2,e_1)=-e_2$, and the coefficients of $e_3$ only come from $b(e_2,e_3)=e_3$ and $b(e_3,e_2)=-e_3$, and the binary coordinate formula is obtained. $J_b$ is an alternating trilinear mapping, so it is uniquely determined by $J_b(e_1,e_2,e_3)$ in the three-dimensional space. The previous basis-by-basis calculation is $-e_3$. Therefore, \[
 \begin{aligned}
 J_b(x,y,z)
 ={}-&\bigl(x_1y_2z_3+x_2y_3z_1+x_3y_1z_2\\
 &\qquad-x_1y_3z_2-x_2y_1z_3-x_3y_2z_1\bigr)e_3.
 \end{aligned}
\] substitutes the zero-order and negative first-order components of the three inputs into the operation definitions. The negative first-order binary operation and the ternary operation containing negative first-order input are all zero, and the latter three formulas are obtained. \[
 \operatorname{StructureConstants}\land\operatorname{GradedExpansion}
 \Rightarrow\operatorname{CoordinateEvaluation}
 \Rightarrow\operatorname{FiniteBasisCertificate}.
\]\end{proof}


\begin{definition}
[Binary coefficients and ternary filling coefficients]\label{exp:higher:parameter-family}
On the same graded space $V$, for $\alpha,\beta\in\mathbb K$ define
\[
 \ell_1^{\alpha,\beta}=\ell_1,
 \qquad \ell_2^{\alpha,\beta}=\alpha\ell_2,
 \qquad \ell_3^{\alpha,\beta}=\beta\ell_3,
 \qquad \ell_n^{\alpha,\beta}=0\quad(n\ge4).
\].
Denote the corresponding arity relation by $\mathcal R_n^{\alpha,\beta}$ and define the ternary operator
$\Omega=\mathcal R_3^{1,0}$.
It retains every ternary defect arising from binary compositions, including homogeneous components with degree-minus-one inputs.\end{definition}


\begin{theorem}
[Parametric equations for complete higher-order compatibility]\label{exp:higher:scaled-closure}
Definition ~ \ref{exp:higher:parameter-family} satisfies \[
 \mathcal R_n^{\alpha,\beta}=0\quad(n\ne3),
 \qquad
 \mathcal R_3^{\alpha,\beta}=(\alpha^2-\beta)\Omega,
 \qquad
 \Omega(e_1,e_2,e_3)=e_3\ne0.
\] Therefore \[
 \operatorname{HigherClosed}(V,\ell_\bullet^{\alpha,\beta})
 \quad\Longleftrightarrow\quad \beta=\alpha^2.
\] This condition holds the specified binary operation coefficient $\alpha$ and determines the required ternary padding coefficient.\end{theorem}
\begin{certificate}
[Collection of coefficients for each element]Save the non-zero operation combinations of ternary numbers one to five, the scalar coefficients, the ary-wise zero relations of the original structure, and the non-zero witness of the ternary operator on $(e_1,e_2,e_3)$.\end{certificate}
\begin{proof}
Element one is $\ell_1^2=0$; each element of element two contains one $\ell_2$, so $\mathcal R_2^{\alpha,\beta}=\alpha\mathcal R_2^{1,1}=0$. The only operation combinations of ternary three are $(2,2)$, $(1,3)$ and $(3,1)$. The sum of the former category is $\Omega$, and the sum of the latter two categories is recorded as $A$. According to the complete ternary relation of theorem ~ \ref{hgov:linfty:model}, $A+\Omega=0$, thus the equation \[
 \mathcal R_3^{\alpha,\beta}
 =\alpha^2\Omega+\beta A=(\alpha^2-\beta)\Omega.
\] holds on all homogeneous inputs. Element four only contains $(2,3)$ and $(3,2)$, so its sum is $\alpha\beta\mathcal R_4^{1,1}=0$. The number five contains only $(3,3)$, so it is $\beta^2\mathcal R_5^{1,1}=0$. Each item with an arity of six or more contains zero-set higher-order operations. For three zeroth-order primitives, $\Omega=-J_b$ , so $\Omega(e_1,e_2,e_3)=e_3$ . So $\alpha^2-\beta=0$ implies that all relations are closed; inversely, $(\alpha^2-\beta)e_3=0$ is obtained by the ternary relation being zero on this primitive group, and $\beta=\alpha^2$ is obtained by $e_3\ne0$.\end{proof}


\begin{corollary}
[Exact residual reading under fixed binary presentation]\label{exp:higher:residual-readout}
Let $e_3^*:E\to\mathbb K$ be the third coordinate functional. Under fixed $\alpha$, define \[
 r_{\rm alg}(\beta)
 =e_3^*\bigl(\mathcal R_3^{\alpha,\beta}(e_1,e_2,e_3)\bigr).
\] then $r_{\rm alg}(\beta)=\alpha^2-\beta$, and have $r_{\rm alg}(\beta)=0$ if and only if the entire set of higher-order relations is closed.\end{corollary}
\begin{certificate}
[Primitive reading and scalar padding]Save the fixed coefficient $\alpha$, the actual coefficient $\beta$, primitive input, third coordinate reading and calculation of the scalar equation $\beta=\alpha^2$.\end{certificate}
\begin{proof}
By the theorem ~ \ref{exp:higher:scaled-closure}, $\mathcal R_3^{\alpha,\beta}(e_1,e_2,e_3)=(\alpha^2-\beta)e_3$. Applying $e_3^*$ yields a scalar residual; zero is equivalent to the complete compatibility condition. This equivalence uses all of the aforementioned checks, while a single read is responsible for extracting the proven residual coordinates.\end{proof}


\chapter{Task Action Graphs, Chain Complexes, and Dynamic Task Quotients}
\label{exp:ch:task}


\label{hgov:sec:taskcomplex}

\section { Task diagram, certification relationship surface and boundary matrix }

\begin{definition}
[Finite task action diagram and continuation surface]\label{hgov:def:taskcomplex}
Fixed finite vertex set $V$ , directed action edge set $E$ , and authenticated continuation relation set $R$ . Each edge $e$ has source $s(e)$ and target $t(e)$. Each relationship $\rho$ gives two admissible paths $\gamma_\rho^+,\gamma_\rho^-$ with the same start and end points, and saves the task compatibility proof required by them. Define \[
C_0=\mathbb K[V],\qquad C_1=\mathbb K[E],\qquad C_2=\mathbb K[R],
\] \[
\partial_1[e]=[t(e)]-[s(e)],\qquad
\partial_2[\rho]=[\gamma_\rho^+]-[\gamma_\rho^-].
\] path chain as the finite sum of directed edges, and other degrees are zero. The relationship surface records the path comparison data that has been given.\end{definition}


\begin{theorem}
[Chain closure and finite generation of task relationships]\label{hgov:thm:complex}
The definition \ref{hgov:def:taskcomplex} gives chain complex \[
C_2\xrightarrow{\partial_2}C_1\xrightarrow{\partial_1}C_0,
\qquad \partial_1\partial_2=0.
\]. When authenticated relationship faces with a common start and end point are added to the graph in sequence, each finite stage is still a chain complex.\end{theorem}
\begin{certificate}
[Boundary entries and face number induction]The source and target are saved edge by edge, and the ordered edge sequences, common start and end points, and compatible certificates of the two paths are saved face by face. Record the boundary matrix $D_1,D_2$ of each stage and each product entry $ (D_1D_2)_{v,\rho}=\sum_e(D_1)_{v,e}(D_2)_{e,\rho}$.\end{certificate}
\begin{proof}
An empty path has zero path chain and its start and end points are the same. For the path $\gamma=(e_1,\ldots,e_n)$ with length $n\ge1$, the adjacent endpoints are equal, so the boundaries of the two paths with the same starting and end points \[
\partial_1[\gamma]
=\sum_{j=1}^n([t(e_j)]-[s(e_j)])
=[t(e_n)]-[s(e_1)].
\] are equal, so each generated surface satisfies $\partial_1\partial_2[\rho]=0$, and linear continuation gives the entire matrix product to be zero. $\partial_2$ of the base case of zero side is zero; when adding a side, only a verified zero product column is added. Therefore, face number induction maintains chaininess. The per-atom predicate chain is \[
\begin{gathered}
\operatorname{ValidPath}(\gamma^\pm)\land
\operatorname{SameEndpoints}(\gamma^+,\gamma^-)
\Rightarrow
\partial_1[\gamma^+]=\partial_1[\gamma^-]\\
\Rightarrow\partial_1\partial_2[\rho]=0
\Rightarrow D_1D_2=0.
\end{gathered}
\]\end{proof}


\section { Semantic cochain, differential and path integral }

\begin{definition} [Semantic cochains, differentials, and path integrals]
\label{hgov:def:semanticdiff}
Set $C^k=\operatorname{Hom}_{\mathbb K}(C_k,\mathbb K)$ and $d^k=(\partial_{k+1})^*$. For an observation $f\in C^0$,
\[
\langle df,[e]\rangle=f(t(e))-f(s(e)).
\].
For $\omega\in C^1$ and a finite path chain $c$, define
$\int_c\omega=\langle\omega,c\rangle$.
This integral includes path orientation, edge coefficients, and the actual observation convention.
\end{definition}

\begin{lemma}
[Squared zeros and path sums of semantic differential]\label{hgov:lem:semanticdiff}
There is $d^{k+1}d^k=0$. For $c=[\gamma]$ , \[
\int_\gamma df=f(t(\gamma))-f(s(\gamma)).
\] if $d\omega=0$ and $c_1-c_0=\partial_2s$ , then $\int_{c_1}\omega=\int_{c_0}\omega$ .\end{lemma}
\begin{certificate}
[Dual and paired atoms]Save the boundary matrix, observation coordinates, each edge difference value, finite sum and $\langle d\omega,s\rangle=\langle\omega,\partial s\rangle$.\end{certificate}
\begin{proof}
For any $c$ , the finite-dimensional linear functional \[
\langle d^2f,c\rangle=\langle f,\partial^2c\rangle=0.
\] zero-by-basis gives $d^2=0$ . The path summation is obtained by the expansion sum of the previous theorem. Finally, \[
\int_{c_1}\omega-\int_{c_0}\omega
=\langle\omega,\partial s\rangle
=\langle d\omega,s\rangle=0.
\] therefore \[
\operatorname{BoundarySquareZero}\Rightarrow\operatorname{DifferentialSquareZero},
\quad
\operatorname{Closed}(\omega)\land(c_1-c_0=\partial s)
\Rightarrow\operatorname{EqualIntegrals}.
\]\end{proof}


\section { Field-by-field retention of dynamic task quotient }

\begin{definition}
[Dynamic task quotient and its chain implementation]\label{hgov:def:quotient}
The following takes the task status set $X=V$. The surjective $q:X\to Q$ is the dynamic sufficient quotient of the specified task, which means that each observation has $O=\bar Oq$, the availability of each action is constant on the quotient fiber, and all legal successors satisfy $qF_a=\bar F_aq$. Goals and identities are also factored separately. If the task requires maintenance fees, register $c(x,a)=\bar c(qx,a)$ at the same time. The chain implementation separately maps $q_k:C_k(X)\to C_k(Q)$ and verifies $q_{k-1}\partial_X=\partial_Qq_k$, among which $q_0[x]=[q(x)]$; the corresponding action side and authentication side are mapped to the quotient side one by one according to the same registered task.\end{definition}


\begin{theorem}
[Co-preservation of semantic differential, task quotient and finite execution]\label{hgov:thm:quotient}
When the conditions defining \ref{hgov:def:quotient} hold, all finite legal action words \[
d_Xq_k^*=q_{k+1}^*d_Q.
\] satisfy $qF_w=\bar F_wq$ and the readback of the specified observation, target, and identity remains.\end{theorem}
\begin{certificate}
[Separate exchange of status and chain]Save the quotient fiber table, observations, available actions table, successor table, chain mapping matrix at all levels, generator by generator boundary exchange and each prefix of the actual word.\end{certificate}
\begin{proof}
For $f\in C^k(Q)$, $c\in C_{k+1}(X)$, \[
\begin{aligned}
\langle d_Xq_k^*f,c\rangle
&=\langle f,q_k\partial_Xc\rangle\\
&=\langle f,\partial_Qq_{k+1}c\rangle
=\langle q_{k+1}^*d_Qf,c\rangle.
\end{aligned}
\] word length zero is the identity; one-step exchange and induction hypothesis give $qF_aF_w=\bar F_aqF_w=\bar F_a\bar F_wq$. The fiber constancy of action availability ensures that the corresponding prefixes are all defined, and then the observation factorization is applied to obtain the conclusion. The predicate causal chain is \[
\begin{gathered}
\operatorname{ChainCommute}(q_\bullet)
\Rightarrow\operatorname{DifferentialCommute}(q^*),\\
\operatorname{DomainMatch}\land\operatorname{StepCommute}
\Rightarrow\operatorname{WordCommute}
\overset{O=\bar Oq}{\Rightarrow}\operatorname{TaskReadback}.
\end{gathered}
\] This gives an atom-by-atom implementation of \cite{S38,GFrameworkSynthesis} that exposes the task quotient and rate--quotient interfaces.\end{proof}


\section { Joint storage of ordered execution records and chain readouts }

\begin{definition}
[Path words and chained mapping]\label{exp:task:word}
Let $\mathsf{Path}(V,E)$ be the path set of finite directed graph, and the empty path is recorded as $1_x$. For $\gamma=e_n\cdots e_1$, the source-target connection condition is $t(e_i)=s(e_{i+1})$. Definition \[
 \kappa(\gamma)=\sum_{i=1}^{n}[e_i],\qquad
 \kappa(1_x)=0.
\] The execution record saves the initial state, edges in order, each action's antecedent and cost; chain readout is calculated from these records.\end{definition}


\begin{proposition}
[Path composition, endpoint boundaries and ordered readback]\label{exp:task:word-chain}
For the continuationable path $\gamma,\eta$ , there is a joint encoding of \[
 \kappa(\eta\gamma)=\kappa(\gamma)+\kappa(\eta),\qquad
 \partial_1\kappa(\gamma)=[t(\gamma)]-[s(\gamma)].
\] and $\mathcal N_{\rm path}(\gamma)=(s(\gamma),(e_1,\ldots,e_n),\kappa(\gamma))$ that is injective on the path collection. When only $\kappa(\gamma)$ is saved, its recoverable observation is exactly the constant path observation on the chained fiber.\end{proposition}
\begin{certificate}
[Ordered edge lists and endpoint elimination]Save the source-target equations, edge sequences, chain coefficients, terminals and initial states of adjacent edges, and register the cancellation of boundary terms edge by edge.\end{certificate}
\begin{proof}
The edge table of a continuation word is the connection of two edge tables, and the sum is the first formula. The boundary expression is expanded to \[
 \partial_1\kappa(\gamma)
 =\sum_{i=1}^{n}([t(e_i)]-[s(e_i)]),
\] and the adjacent terms are eliminated by $t(e_i)=s(e_{i+1})$, leaving the terminal minus initial state; both sides of the empty word are zero. The ordered edge list and initial state in joint coding directly restore the path, so injective. If the path observation $O$ has read back $D\kappa=O$, the two paths of the same chain will get the same observation through $D$; conversely, when $O$ is constant on the fiber, the definition $D(\kappa\gamma)=O(\gamma)$ is well-defined. For example, the two named self-loops $a,b$ on a single vertex satisfy $\kappa(ab)=\kappa(ba)=[a]+[b]$, but $ab,ba$ are different ordered words. Joint encoding preserves the two words separately, and chain reading preserves their common edge count.\end{proof}


\section { Complete matrix calculation example of four-state continuation surface }

\begin{definition}
[Square execution graph for 2D translation task]\label{exp:task:square}
Take the vertex $v_{00},v_{10},v_{01},v_{11}$ and represent the corresponding binary coordinates respectively. The action $a$ adds the first coordinate, and the action $b$ adds the second coordinate; only legal actions remaining within these four vertices are registered, and the cost is one. Register edge \[
 e_1:v_{00}\to v_{10},\quad e_2:v_{10}\to v_{11},\quad
 e_3:v_{00}\to v_{01},\quad e_4:v_{01}\to v_{11}.
\] in order. Two paths $e_2e_1$ and $e_4e_3$ have the same end and the same total cost. Register the specified terminal task compatibility relationship as a face $\rho$, $\partial_2\rho=e_1+e_2-e_3-e_4$.\end{definition}


\begin{theorem}
[Chain closure, homology and semantic integration of the square task]\label{exp:task:square-homology}
The boundary matrix of the above complex is \[
 D_1=\begin{pmatrix}
 -1&0&-1&0\\1&-1&0&0\\0&0&1&-1\\0&1&0&1
 \end{pmatrix},\qquad
 D_2=\begin{pmatrix}1\\1\\-1\\-1\end{pmatrix}.
\] which satisfies $D_1D_2=0$ and has \[
 H_2=0,\qquad H_1=0,\qquad H_0\cong\mathbb K.
\] pairs of observations $f(x,y)=x+2y$, the vertex value is $(0,1,2,3)^{\mathsf T}$ and the edge difference is $(1,2,2,1)^{\mathsf T}$. The differential integral on both paths is three.\end{theorem}
\begin{certificate}
[Matrix, Rank, Goal and Cost]Save four vertices, four action domains, two actual paths and their target $v_{11}$. Save $D_1,D_2$, a third-order non-zero subformula, kernel generator, vertex observations and edge-by-edge differences. The terminal is compatible and the cost of saving two additional paths is two.\end{certificate}
\begin{proof}
Reading the end point minus the starting point column by column gives $D_1$; the plane direction gives $D_2$. Multiply row by row to get \[
 D_1D_2=\begin{pmatrix}-1+1\\1-1\\-1+1\\1-1\end{pmatrix}=0.
\] $D_1$. The sum of the components in each column is zero, so the rank is at most three. The subformulas of the first three rows and the first three columns are one, so the rank is exactly three. So $\dim\ker D_1=1$ and the non-zero vector $D_2$ form this core. The rank of $D_2$ is one and the domain is one-dimensional, so $H_2=0$, $H_1=0$; $\dim H_0=4-3=1$. The sum of vertex coefficients drops to the isomorphism of $H_0\to\mathbb K$. The cochain is calculated as \[
 D_1^{\mathsf T}\begin{pmatrix}0\\1\\2\\3\end{pmatrix}
 =\begin{pmatrix}1\\2\\2\\1\end{pmatrix},\qquad
 D_2^{\mathsf T}D_1^{\mathsf T}f=1+2-2-1=0.
\], so the two path integrals are $1+2=3$ and $2+1=3$ respectively, which are equal to the common terminal value minus the initial value.\end{proof}


\section { Homology and cochain reading of compound chain mapping }

\begin{theorem}
[Induced mapping and finite composition of chain maps]\label{exp:task:homology-functor}
If $f:C_\bullet\to D_\bullet$ is a chain mapping, then \[
 H_k(f):H_k(C)\to H_k(D),\qquad [c]\mapsto[f_kc]
\] is well-defined. For chain mapping $g:D_\bullet\to E_\bullet$, \[
 H_k(gf)=H_k(g)H_k(f),\qquad H_k(\mathrm{id})=\mathrm{id}.
\] corresponds to the cochain pullback that satisfies $(gf)^*=f^*g^*$, and is commutation with cochain with differentials.\end{theorem}
\begin{certificate}
[Closed links, boundaries and composite matrices]Save each boundary exchange $\partial_Df_k=f_{k-1}\partial_C$, closed chain representation, boundary original image representing element difference, and composite mapping matrix.\end{certificate}
\begin{proof}
If $\partial_Cc=0$, then $\partial_Df_kc=f_{k-1}\partial_Cc=0$, so the closed chain is reflected as a closed chain. If $c'=c+\partial_Cb$ , then \[
 f_kc'=f_kc+f_k\partial_Cb
 =f_kc+\partial_Df_{k+1}b,
\] therefore the two images represent the same homology class. The conclusions of composition and identity are established by substituting the representative elements step by step. For cochain $\omega$, there is $((gf)^*\omega)(c)=\omega(gfc)=(f^*g^*\omega)(c)$; and commutation with differentials are obtained by duality of the commutation relation with the same boundary. The conclusions of finite multiple mappings are summarized by the number of mappings. \[
 \operatorname{ChainCommute}(f)
 \Rightarrow\operatorname{CyclesPreserved}\land\operatorname{BoundariesPreserved}
 \Rightarrow\operatorname{WellDefined}(H_kf).
\]\end{proof}


\chapter{Negative Curvature Geometry, Variational Growth, and Task Observability}
\label{exp:ch:geometry}


\label{hgov:geo:section}

The public G-Framework uses GZ-GMS to register objects, law rules, admissible evolution, property layering and connection, and uses triple-layer connection to organize the joint transportation of geometry, parameters and law rule directions \cite{S01,GFrameworkSynthesis}. This chapter constructs geometric variations and task observations on structural sectors with smooth realizations, and the next chapter establishes inter-bundle transport and cochain continuations. Each construct saves its substrate type, valid domain, and exchange relationships; the resulting geometric certificate can be used as input for subsequent obstacle determination and repair. The geometric calculation goes through the metric, connection, variational field, observation derivative and finite error budget in sequence, and finally outputs the differentiability at the specified accuracy.

\section { Parting geometry substrate and curvature convention }\label{hgov:geo:typed}


\begin{definition}
[Fractal geometry implementation of law-space]\label{hgov:geo:def:sector}
Let $I$ be the index set of the registered structure sector. The geometric implementation of sector $\alpha\in I$ contains \[
\mathfrak G_\alpha=
\bigl(M_\alpha,g_\alpha,
P_\alpha\xrightarrow{\pi_\alpha}M_\alpha,
G_\alpha,\rho_\alpha,E_\alpha,\nabla^\alpha\bigr).
\] where $M_\alpha$ is a smooth manifold, $g_\alpha$ is a Riemannian metric, $P_\alpha$ is a smooth principal bundle with the structural group $G_\alpha$, $\rho_\alpha:G_\alpha\to\operatorname{GL}(V_\alpha)$ is a finite-dimensional real or complex representation, $E_\alpha=P_\alpha\times_{\rho_\alpha}V_\alpha$ is an adjoint vector bundle, and $\nabla^\alpha$ is a vector bundle induced by a registered main connection connection . The rank of each vector bundle is fixed on each connected component; different sectors store their basis spaces, structural groups, fiber types and allowed mappings respectively. The Levi--Civita connection measuring $g_\alpha$ is recorded as $\nabla^{g_\alpha}$, and the connection of the task coefficient bundle is recorded as $\nabla^\alpha$. The principal bundle and connection meanings of these data adopt the usual smooth geometry definition \cite{ExtKobayashiNomizu}.\end{definition}


\begin{definition}
[Negative Section Curvature and Transverse Variation]\label{hgov:geo:def:jacobi}
The abbreviated Riemann implementation in the fixed sector is $(M,g)$, and the curvature convention \[
R^g(X,Y)Z=
\nabla^g_X\nabla^g_YZ-
\nabla^g_Y\nabla^g_XZ-
\nabla^g_{[X,Y]}Z,
\qquad
\operatorname{Sec}_g(u,v)=
\frac{g(R^g(u,v)v,u)}{
g(u,u)g(v,v)-g(u,v)^2}.
\] is used. The cross-sectional curvature is defined when $u,v$ is linearly independent. Let $\gamma:[0,T]\to M$ be the unit speed geodesic, $T>0$; let $J$ be the Jacobi field along $\gamma$, satisfying \begin{equation}
\label{hgov:geo:eq:jacobi}
D_t^2J+R^g(J,\dot\gamma)\dot\gamma=0,
\qquad J(0)=0,\qquad D_tJ(0)=v,
\qquad 0\ne v\perp\dot\gamma(0),
\end{equation} where $D_t=\nabla^g_{\dot\gamma}$. Given $\kappa>0$ , the condition for negative curvature along this curve is: for every $t\in[0,T]$ and every nonzero $w\perp\dot\gamma(t)$ , there is \[
\operatorname{Sec}_g(w,\dot\gamma(t))\le-\kappa.
\] The Jacobi field gives the infinitesimal variation of the geodesic family with respect to the initial conditions; the connections, curvature signs and initial values used are jointly registered with the above formula \cite{ExtLeeRiemann}.\end{definition}


\begin{remark}
[Local coordinates and tensor types]Write $g=g_{ij}\,dx^i\otimes dx^j$ in the coordinates $x^1,\ldots,x^n$, and record the inverse matrix as $(g^{ij})$. Levi--Civita connection is uniquely determined by metric compatibility and inflexibility; its coefficient is \[
\Gamma^k_{ij}=\frac12\sum_\ell g^{k\ell}
\bigl(\partial_i g_{j\ell}+\partial_jg_{i\ell}-\partial_\ell g_{ij}\bigr).
\] The vector field $J=J^k\partial_k$ along the curve satisfies \[
(D_tJ)^k=\frac{dJ^k}{dt}
+\sum_{i,j}\Gamma^k_{ij}(\gamma(t))\dot\gamma^iJ^j,
\qquad
\ddot\gamma^k+\sum_{i,j}\Gamma^k_{ij}\dot\gamma^i\dot\gamma^j=0.
\] so the coordinate certificate can be generated item by item from the first and second order partial derivatives of $g_{ij}$. If there is another task function $L:M\to\mathbb R$, its covariant Hessian is \[
(\operatorname{Hess}_gL)_{ij}
=\partial_i\partial_jL-\sum_k\Gamma^k_{ij}\partial_kL
\in\Gamma(\operatorname{Sym}^2T^*M).
\], the Riemann curvature $R^g$ belongs to $\Omega^2(M,\operatorname{End}(TM))$, and the curvature of the task coefficient connection $F_\nabla$ belongs to $\Omega^2(M,\operatorname{End}(E))$. When registering $E=TM$ and $\nabla=\nabla^g$, they satisfy $F_\nabla(X,Y)Z=R^g(X,Y)Z$. These type identifications are saved with metric selections for comparison between different geometric readouts.\end{remark}


\section { Item-by-item calculation of saddle surface and variation norm }\label{hgov:geo:observation}


\begin{proposition}
[Explicit negative curvature implementation of saddle surfaces]\label{hgov:geo:prop:saddle}
Take the embeddings $\iota:\mathbb R^2\to\mathbb R^3$ , $\iota(x,y)=(x,y,xy)$ , and use the Euclidean inner product-induced metric on the surface. Its Gaussian curvature is \[
K(x,y)=-\frac{1}{(1+x^2+y^2)^2}.
\]. On the area of $x^2+y^2\le R_0^2$, there is $K\le-(1+R_0^2)^{-2}<0$.\end{proposition}
\begin{certificate}
[First and second basic forms and determinants]Let $D=1+x^2+y^2$, save \[
g=\begin{pmatrix}1+y^2&xy\\xy&1+x^2\end{pmatrix},
\quad
n=\frac{(-y,-x,1)}{\sqrt D},
\quad
\mathrm{II}=\begin{pmatrix}0&D^{-1/2}\\D^{-1/2}&0\end{pmatrix}.
\] and save the embedded first and second order partial derivatives, unit normals and determinants of the two matrices item by item.\end{certificate}
\begin{proof}
The partial derivatives are $\iota_x=(1,0,y)$ , $\iota_y=(0,1,x)$ , so the first fundamental form is listed as $g$ and $\det g=D$ . The second-order derivative satisfies $\iota_{xx}=\iota_{yy}=0$ and $\iota_{xy}=(0,0,1)$. When paired with $n$, $\mathrm{II}$ is obtained, and its determinant is $-D^{-1}$. The matrix of the shape operator is $g^{-1}\mathrm{II}$ , the Gaussian equation gives \[
K=\det(g^{-1}\mathrm{II})
=\frac{-D^{-1}}{D}=-D^{-2}.
\] and the zone condition gives $1\le D\le1+R_0^2$ , thus obtaining the listed negative curvature bounds. Therefore \[
\operatorname{InducedMetric}(\iota)
\Rightarrow(\det g=D,\ \det\mathrm{II}=-D^{-1})
\Rightarrow K=-D^{-2}
\Rightarrow\operatorname{NegCurv}_{(1+R_0^2)^{-2}} .
\]\end{proof}


\begin{remark}
[Valid time of area curvature bound]In the saddle surface example, if the parametric coordinates of the geodesic satisfy $x(t)^2+y(t)^2\le R_0^2$ for all $t\in[0,T]$, then $\kappa=(1+R_0^2)^{-2}$ can be taken along the interval. This time interval enters the negative curvature certificate along with the proof that the curve remains in the region. The second derivative of the coordinate function $xy$ is used to calculate the second fundamental form, and the final section curvature is determined by both the induced metric and this fundamental form.\end{remark}


\begin{lemma} [Preservation of transversality and a differential inequality for the norm]
\label{hgov:geo:lem:norm}
Under the conditions of Definition ~ \ref{hgov:geo:def:jacobi}, we have $J(t)\perp\dot\gamma(t)$. On an interval where $J(t)\ne0$, set $f(t)=\lVert J(t)\rVert_g$. Then
\begin{equation}
\label{hgov:geo:eq:norm-inequality}
f(t)f''(t)+(f'(t))^2
=\lVert D_tJ(t)\rVert_g^2
-g(R^g(J(t),\dot\gamma(t))\dot\gamma(t),J(t)),
\qquad f''(t)\ge\kappa f(t).
\end{equation}
\end{lemma}

\begin{certificate}
[Transversality and term-by-term differentiation certificate]Save $g$, $\nabla^g$, $\gamma$, $J$, $v$, definition interval and upper bound of curvature. The certificate itemizes the geodesic equations, the formula ~ \eqref{hgov:geo:eq:jacobi} , initial orthogonality, metric symmetry relations for curvature tensors, and applicable inner products of Cauchy--Schwarz inequalities. Curvature bounds can be provided by analytical proofs or by verification bounds covering the entire claimed interval.\end{certificate}


\begin{proof}
Order $a(t)=g(J(t),\dot\gamma(t))$. Due to the antisymmetry of $D_t\dot\gamma=0$ and the curvature tensor in the last two indicators, the initial value of \[
a''(t)=-g(R^g(J,\dot\gamma)\dot\gamma,\dot\gamma)=0.
\] is $a(0)=0$ and $a'(0)=g(v,\dot\gamma(0))=0$, so $a(t)=0$. $q=g(J,J)=f^2$ is placed first in the interval of $f>0$. From $\nabla^g g=0$, differentiate term by term to obtain \[
\begin{aligned}
q'&=g(D_tJ,J)+g(J,D_tJ)=2g(D_tJ,J),\\
q''&=2g(D_t^2J,J)+2g(D_tJ,D_tJ),\\
f'&=\frac{q'}{2\sqrt q}=\frac{g(D_tJ,J)}{f},\\
f''&=\frac{q''}{2f}-\frac{(q')^2}{4f^3}
=\frac{\lVert D_tJ\rVert_g^2+g(D_t^2J,J)}{f}
-\frac{g(D_tJ,J)^2}{f^3}.
\end{aligned}
\]. Multiply $f$ and substitute it into the Jacobi equation to obtain the equation of ~ \eqref{hgov:geo:eq:norm-inequality}. To record the origin of each item, further let \[
n=J/f,\qquad W=D_tJ-f'n.
\] due to $g(n,n)=1$ and $f'=g(D_tJ,n)$, and $g(W,n)=0$ and $\lVert D_tJ\rVert_g^2=(f')^2+\lVert W\rVert_g^2$. $J\perp\dot\gamma$ and $\lVert\dot\gamma\rVert=1$ give \[
g(R^g(J,\dot\gamma)\dot\gamma,J)
=\operatorname{Sec}_g(J,\dot\gamma)f^2.
\] and thus obtain the exact decomposition \begin{equation}
\label{exp:geo:eq:norm-remainder}
f''-\kappa f
=\frac{\lVert W\rVert_g^2}{f}
+\bigl(-\operatorname{Sec}_g(J,\dot\gamma)-\kappa\bigr)f.
\end{equation}. The first term is non-negative by the orthogonal decomposition and the second term is non-negative by the registered section curvature bound, so \[
\bigl(W\perp n\bigr)\land
\bigl(\operatorname{Sec}_g\le-\kappa\bigr)\land(f>0)
\Longrightarrow f''-\kappa f\ge0.
\] This also gives a term-by-term verifiable expression of the remainder of the comparison equation.\end{proof}


\section { Scalar comparison, zero point continuation and constant curvature examples }\label{exp:geo:section:comparison}


\begin{theorem}
[Transverse variational growth under negative curvature]\label{hgov:geo:thm:growth}
Under the conditions of Definition ~ \ref{hgov:geo:def:jacobi}, for all $t\in[0,T]$,
\begin{equation}
\label{hgov:geo:eq:growth}
\lVert J(t)\rVert_g
\ge\frac{\lVert v\rVert_g}{\sqrt\kappa}
\sinh(\sqrt\kappa\,t).
\end{equation}.
In particular, $J(t)\ne0$ holds for every $t\in(0,T]$. If the geodesic, initial-value problem, and curvature conditions hold on all finite intervals, the same estimate extends to all $t\ge0$.\end{theorem}


\begin{certificate}
[Scalar comparisons and zero-point continuation certificate]Order $b=\sqrt\kappa$. Save the initial right derivatives $f(0)=0$, $f'(0^+)=\lVert v\rVert_g$, and the remainder $h=f''-b^2f\ge0$ on the positive interval. Compare the certificate to \begin{equation}
\label{hgov:geo:eq:comparison}
f(t)=\frac{\lVert v\rVert_g}{b}\sinh(bt)
+\int_0^t\frac{\sinh(b(t-s))}{b}h(s)\,\mathrm ds.
\end{equation} and save the exclusion arguments for the initial positive neighborhood and the first positive zero point.\end{certificate}


\begin{proof}
Take a parallel orthogonal frame along $\gamma$, and write $J$ as the Euclidean vector function $u(t)$. The Jacobi equation becomes $u''(t)+C(t)u(t)=0$ , where $C$ is smooth, $u(0)=0$ , and $u'(0)=v$ . So $u''(0)=-C(0)u(0)=0$, Taylor expansion gives $u(t)=tv+O(t^3)$. More specifically, set \[
a(t)=\int_0^1u'(st)\,ds,
\qquad u(t)=t a(t),\qquad a(0)=v,\qquad a'(0)=0.
\] $a$ to be smooth and non-zero around zero. So in the right neighborhood $f(t)=t\lVert a(t)\rVert$ is smoothed to the endpoint, satisfying \[
f(0)=0,\quad f'(0^+)=\lVert v\rVert,
\quad f(t)=t\lVert v\rVert+O(t^3),\quad f''(t)=O(t).
\] This gives the initial positive neighborhood and makes $h=f''-b^2f$ a continuous continuation to $0$ with a value of zero. Lemma ~ \ref{hgov:geo:lem:norm} gives $h\ge0$ on any positive interval continued by this neighborhood. Let \[
G(t)=\int_0^t\frac{\sinh(b(t-s))}{b}h(s)\,ds.
\] be the upper limit of differential integration and the integral kernel, respectively, we get \[
G'(t)=\int_0^t\cosh(b(t-s))h(s)\,ds,
\qquad
G''(t)=h(t)+b^2G(t).
\] and thus $G(0)=G'(0)=0$, and $\lVert v\rVert\sinh(bt)/b+G(t)$ satisfies the same scalar equation and initial value as $f$. The uniqueness of the scalar linear initial value problem is given by the formula ~ \eqref{hgov:geo:eq:comparison}. Both the integration kernel and $h$ are non-negative, so the formula ~ \eqref{hgov:geo:eq:growth} is obtained.

If $f$ has a zero point at $(0,T]$, the first positive zero point $t_*$ can be obtained based on the initial positive value and continuity. The aforementioned lower bound is established at $0<t<t_*$; this causes a contradiction between $t\uparrow t_*$ and \[
0=f(t_*)\ge\frac{\lVert v\rVert_g}{b}\sinh(bt_*)>0,
\]. Therefore, there is no positive zero point, and the comparison estimate holds over the entire $[0,T]$. When all finite intervals satisfy the same condition, the last conclusion can be obtained by selecting the finite interval containing it at any given moment.\end{proof}


\begin{proposition}
[Exact Jacobi solution under constant radial curvature]\label{exp:geo:prop:constant-jacobi}
Along the unit speed geodesic $\gamma$, it is assumed that there is a real constant $\rho$ such that all transverse vectors satisfy $R^g(w,\dot\gamma)\dot\gamma=\rho w$. The initial value of reserved ~ \eqref{hgov:geo:eq:jacobi}, let $P_t$ be the Levi--Civita parallel transportation from $\gamma(0)$ to $\gamma(t)$, then \[
J(t)=S_\rho(t)P_tv,
\qquad
S_\rho(t)=
\begin{cases}
\sin(\sqrt\rho\,t)/\sqrt\rho,&\rho>0,\\
t,&\rho=0,\\
\sinh(\sqrt{-\rho}\,t)/\sqrt{-\rho},&\rho<0.
\end{cases}
\] therefore $\lVert J(t)\rVert=|S_\rho(t)|\lVert v\rVert$. When taking $\rho=-\kappa$, the formula ~ \eqref{hgov:geo:eq:growth} takes equal signs everywhere.\end{proposition}
\begin{certificate}
[Parallel frame and scalar initial value solution]Save the metric-preserving, radial curvature effects of $P_t$, $\rho\operatorname{id}$, $S_\rho''+\rho S_\rho=0$, $S_\rho(0)=0$, $S_\rho'(0)=1$. The declared time interval is taken as the interval in which the geodesic and the curvature effect are both effective.\end{certificate}
\begin{proof}
From $D_t(P_tv)=0$ , there is \[
D_t(S_\rho P_tv)=S_\rho'P_tv,
\qquad
D_t^2(S_\rho P_tv)=S_\rho''P_tv.
\] parallel transport keeping the inner product, hence $P_tv\perp\dot\gamma(t)$ . Substituting the curvature assumption and the scalar equation, the initial value of \[
D_t^2(S_\rho P_tv)+R^g(S_\rho P_tv,\dot\gamma)\dot\gamma
=(S_\rho''+\rho S_\rho)P_tv=0.
\] is the same as $(0,v)$. Therefore, the uniqueness of Jacobi equation gives the listed $J$. The norm formula is obtained from $\lVert P_tv\rVert=\lVert v\rVert$. The equations and initial values of the three $S_\rho$ can be verified by term-by-term differentiation; when $\rho=-\kappa$, $S_\rho(t)>0$ is established for $t>0$, and is consistent with the right end of the formula ~ \eqref{hgov:geo:eq:growth}. Therefore the certificate chain is \[
\text{\shortstack[c]{\scriptsize Constant radial\\[-1pt]\scriptsize curvature}}\land\text{\shortstack[c]{\scriptsize Parallel frame}}
\Longrightarrow\text{\shortstack[c]{\scriptsize Scalar\\[-1pt]\scriptsize constant-coefficient\\[-1pt]\scriptsize equation}}
\Longrightarrow\text{\shortstack[c]{\scriptsize Exact variation\\[-1pt]\scriptsize and its norm}}.
\]\end{proof}


\begin{proposition}
[Exact separation of finite angles in the hyperbolic plane]\label{exp:geo:prop:hyperbolic-separation}
Let $a>0$ , take the Lorentz bilinear $\langle X,Y\rangle_L=-X_0Y_0+X_1Y_1+X_2Y_2$ on $\mathbb R^3$ , and let \[
\mathbb H_a^2=\{X:\langle X,X\rangle_L=-a^{-2},\ X_0>0\}
\] be the restricted Riemannian metric. The cross-sectional curvature of this measure is $-a^2$ . Two unit velocity rays starting from $o=(a^{-1},0,0)$ and with an included angle of $\theta\in[0,\pi]$ can be written as \[
\gamma_\theta(t)=\frac1a
\bigl(\cosh(at),\sinh(at)\cos\theta,\sinh(at)\sin\theta\bigr).
\], denoted by $d_\theta(t)=d(\gamma_0(t),\gamma_\theta(t))$, then \begin{equation}
\label{exp:geo:eq:finite-angle}
d_\theta(t)=\frac2a\operatorname{arsinh}
\left(\sinh(at)\sin\frac\theta2\right).
\end{equation} under fixed $t\ge0$, $\lim_{\theta\downarrow0}d_\theta(t)/\theta=\sinh(at)/a$ . Under fixed $\theta\in(0,\pi]$, \[
d_\theta(t)=2t+\frac2a\log\sin(\theta/2)+O(e^{-2at})
\qquad(t\to\infty).
\]\end{proposition}
\begin{certificate}
[Embedding, Radial Metrics, and Hyperbolic Distances]Saves Lorentz types, upper lobe conditions for hyperboloids, positive definite limits on tangent spaces, ray velocities and accelerations. Distance certificate save \[
\cosh\bigl(a\,d(X,Y)\bigr)=-a^2\langle X,Y\rangle_L,
\qquad
1+2\sinh^2z=\cosh(2z).
\] The two limits respectively register the quantifiers of fixed time and fixed non-zero angle and their valid domains. This model and its Riemannian calculations use the usual hyperbolic space convention \cite{ExtLeeRiemann}.\end{certificate}
\begin{proof}
In the tangent space of $o$, the Lorentz form is restricted to the Euclidean inner product over the last two coordinates. For any point $X$, first perform Euclidean rotation of the last two coordinates to make it $(X_0,r,0)$, where $r\ge0$. Let $c=aX_0$ , $s=ar$ , then $c^2-s^2=1$ , matrix \[
L=\begin{pmatrix}c&-s&0\\-s&c&0\\0&0&1\end{pmatrix}
\] maintains Lorentz type and upper lobe, and $L(X_0,r,0)^{\mathsf T}=(a^{-1},0,0)^{\mathsf T}$ . Therefore, there is an equidistant property change that moves any point to $o$, and the limiting metric is positive definite everywhere. Let the environment straight derivative be $D$. Tangent vector field $X,Y$, differentiate $\langle Y,p\rangle_L=0$, get $\langle D_XY,p\rangle_L=-g(X,Y)$, so the environmental curvature of \[
D_XY=\nabla^g_XY+a^2g(X,Y)p.
\] is zero, take the tangential component of $D_XD_YZ-D_YD_XZ-D_{[X,Y]}Z$, we can get \[
R^g(X,Y)Z=-a^2\bigl(g(Y,Z)X-g(X,Z)Y\bigr),
\], so the cross-sectional curvature is $-a^2$. The listed rays satisfy $\lVert\dot\gamma_\theta\rVert_g=1$ and $\ddot\gamma_\theta=a^2\gamma_\theta$; the acceleration is normal, so it is a unit velocity geodesic.

In the radial coordinate centered on $o$, directly differentiate the above parametric formula to obtain \[
g=dr^2+\frac{\sinh^2(ar)}{a^2}\,d\theta^2.
\]. From $o$ to the point with the radial coordinate $r$, the length of any path is at least the total change of the radial coordinate $r$, and the radial curve reaches this value. Therefore $d(o,Y)=r$ , and $\cosh(ar)=aY_0=-a^2\langle o,Y\rangle_L$ . Then use the isometric transformation that maintains the Lorentz type to move the first point, and obtain the general distance formula in certificate. Now directly calculate \[
\begin{aligned}
-a^2\langle\gamma_0(t),\gamma_\theta(t)\rangle_L
&=\cosh^2(at)-\sinh^2(at)\cos\theta\\
&=1+2\sinh^2(at)\sin^2(\theta/2).
\end{aligned}
\]. Since the distance is non-negative, apply $\cosh(2\operatorname{arsinh}z)=1+2z^2$ to get the formula ~ \eqref{exp:geo:eq:finite-angle}.

With $t$ fixed, $\operatorname{arsinh}z=z+O(z^3)$ and $\sin(\theta/2)=\theta/2+O(\theta^3)$ give the first limit. The corresponding angular variational field is $J(t)=\partial_\theta\gamma_\theta(t)|_{\theta=0}
=a^{-1}(0,0,\sinh(at))$, and its norm is exactly this limit. When fixing non-zero $\theta$, use $\operatorname{arsinh}z=\log(2z)+O(z^{-2})$ and $\sinh(at)=\tfrac12e^{at}(1-e^{-2at})$ to get the last expansion. These calculations save infinitesimal variations, finite angular distances, and their limits respectively, making it easy to select observation parameters on the same time scale.\end{proof}


\section { finite observation difference, Taylor remainder and measurement accuracy }\label{exp:geo:section:finite-observation}


\begin{definition}
[Task observation with remainder and measurement accuracy]\label{hgov:geo:def:observation}
Assume $\Upsilon:[-\varepsilon_0,\varepsilon_0]\times[0,T]\to M$ is a smooth variation that satisfies $\Upsilon(0,t)=\gamma(t)$ and $\partial_\varepsilon\Upsilon(0,t)=J(t)$. The task observation is that $C^2$ maps $\mathcal O:U\to\mathbb R^m$ , where $U\subseteq M$ is an open set, and $\Upsilon([-\varepsilon_0,\varepsilon_0]\times[0,T])\subseteq U$ . Note that \[
B_t=D\mathcal O_{\gamma(t)},\qquad
y_\varepsilon(t)=\mathcal O(\Upsilon(\varepsilon,t)).
\] is given the registered subspace $V_t\subseteq T_{\gamma(t)}M$ containing $J(t)$, assuming that there is a constant $c>0$ such that \[
\lVert B_tw\rVert\ge c\lVert w\rVert_g
\quad(w\in V_t),\qquad
\sup_{|\sigma|\le\varepsilon_0}
\left\lVert\partial_\varepsilon^2
(\mathcal O\circ\Upsilon)(\sigma,t)\right\rVert\le L_{\mathcal O}(t).
\] actually reads $\widehat y_0(t)$, $\widehat y_\varepsilon(t)$ The error bounds are $e_0(t)$ and $e_\varepsilon(t)$ respectively; the distinction threshold of the observation protocol is $\eta(t)\ge0$.\end{definition}


\begin{proposition}
[Sufficient conditions for geometric variation to finite observational distinction]\label{hgov:geo:prop:distinguish}
In the common valid domain of theorem ~ \ref{hgov:geo:thm:growth} and definition ~ \ref{hgov:geo:def:observation}, if \begin{equation}
\label{hgov:geo:eq:distinguish}
\frac{c|\varepsilon|\lVert v\rVert_g}{\sqrt\kappa}
\sinh(\sqrt\kappa\,t)
-\frac{L_{\mathcal O}(t)}2\varepsilon^2
-e_0(t)-e_\varepsilon(t)>\eta(t),
\end{equation} then $\lVert\widehat y_\varepsilon(t)-\widehat y_0(t)\rVert>\eta(t)$.\end{proposition}


\begin{certificate}
[Observation mapping, non-degenerate lower bounds and error budgets]Retain $\Upsilon$, $\mathcal O$, $V_t$, $B_t$, $c$, $L_{\mathcal O}(t)$, $\varepsilon_0$, the chosen $\varepsilon,t$, both measurement error bounds, and the distinction threshold. In an orthonormal frame, use $\bigl(B_t|_{V_t}\bigr)^*\bigl(B_t|_{V_t}\bigr)\succeq c^2\operatorname{id}_{V_t}$ to verify the lower bound. Verify the second-derivative bound on the entire parameter segment. A numerical or symbolic check of equation ~ \eqref{hgov:geo:eq:distinguish} serves as the final distinction witness.\end{certificate}


\begin{proof}
Fixed $t$, let $Y_t(\sigma)=\mathcal O(\Upsilon(\sigma,t))$. The chain rule gives $Y_t'(0)=B_tJ(t)$ . In order to handle positive and negative $\varepsilon$ at the same time, write the integral remainder in the fixed interval $[0,1]$: \[
\begin{aligned}
y_\varepsilon-y_0
&=\varepsilon B_tJ(t)+r_\varepsilon(t),\\
r_\varepsilon(t)
&=\varepsilon^2\int_0^1(1-s)Y_t''(s\varepsilon)\,ds,\\
\lVert r_\varepsilon(t)\rVert
&\le\varepsilon^2L_{\mathcal O}(t)\int_0^1(1-s)\,ds
=\frac{L_{\mathcal O}(t)}2\varepsilon^2.
\end{aligned}
\] If the variation falls in the coordinate diagram, record its coordinates as $x^i(\sigma,t)$, then each observation component satisfies \[
Y_{t,\ell}''(\sigma)
=\sum_{i,j}(\partial_i\partial_j\mathcal O_\ell)
\,\partial_\sigma x^i\partial_\sigma x^j
+\sum_i(\partial_i\mathcal O_\ell)\,\partial_\sigma^2x^i.
\] and the right end is all in Evaluate at $\Upsilon(\sigma,t)$. Therefore, the verification of $L_{\mathcal O}(t)$ includes the second-order change of the observation map and the second-order change of the variation curve; when multiple coordinate diagrams are covered, it can be calculated piece by piece and a common upper bound can be taken. Applying the inverse triangle inequality, observing the lower bound, theorem ~ \ref{hgov:geo:thm:growth} and measuring error bound in sequence, we get \[
\begin{aligned}
\lVert\widehat y_\varepsilon-\widehat y_0\rVert
&\ge\lVert y_\varepsilon-y_0\rVert-e_0-e_\varepsilon\\
&\ge c|\varepsilon|\lVert J(t)\rVert_g
-\frac{L_{\mathcal O}(t)}2\varepsilon^2-e_0-e_\varepsilon\\
&\ge\frac{c|\varepsilon|\lVert v\rVert_g}{\sqrt\kappa}
\sinh(\sqrt\kappa\,t)
-\frac{L_{\mathcal O}(t)}2\varepsilon^2-e_0-e_\varepsilon.
\end{aligned}
\]. Finally, we get the conclusion from the formula ~ \eqref{hgov:geo:eq:distinguish}. Its predicate continuation is \[
\begin{gathered}
\text{\shortstack[c]{\scriptsize Curvature bounds\\[-1pt]\scriptsize and variation\\[-1pt]\scriptsize initial data}}
\Longrightarrow\text{\shortstack[c]{\scriptsize Variation norm\\[-1pt]\scriptsize lower bound}},\\
\text{\shortstack[c]{\scriptsize Variation lower\\[-1pt]\scriptsize bound}}\land\text{\shortstack[c]{\scriptsize Observation\\[-1pt]\scriptsize lower bound}}\land\text{\shortstack[c]{\scriptsize Remainder and\\[-1pt]\scriptsize accuracy bounds}}
\Longrightarrow\text{\shortstack[c]{\scriptsize Distinguishable\\[-1pt]\scriptsize task\\[-1pt]\scriptsize observations}}.
\end{gathered}
\]\end{proof}


\begin{proposition}
[Observe differentiation under derivative perturbations certificate stability]\label{exp:geo:prop:observation-robustness}
Fixed $t$ , let $B_t|_{V_t}$ have a lower bound ~ defined in \ref{hgov:geo:def:observation} $c>0$ . Another $C^2$ observation is that the derivative of $\widetilde{\mathcal O}$ is $\widetilde B_t$ , and \[
\bigl\|(\widetilde B_t-B_t)|_{V_t}\bigr\|\le\delta<c.
\] then $\lVert\widetilde B_tw\rVert\ge(c-\delta)\lVert w\rVert_g$ holds for $w\in V_t$ . If the alternative observation has the second-order residual bound $\widetilde L(t)$ and error bound $\widetilde e_0(t),\widetilde e_\varepsilon(t)$, then the strict inequalities obtained by substituting $c-\delta,\widetilde L,\widetilde e_0,\widetilde e_\varepsilon$ in the equations ~ \eqref{hgov:geo:eq:distinguish} give the same distinguishable conclusion for the alternative readings.\end{proposition}
\begin{certificate}
[Restricted operator norm and remainder versions]Save the common subspace $V_t$, its inner product, the effective domain of the two observations, and the operator norm bound of $(\widetilde B_t-B_t)|_{V_t}$. The Taylor remainder of the alternative observation is saved separately from the measurement accuracy, and $c-\delta>0$ is checked against the final threshold inequality.\end{certificate}
\begin{proof}
For any $w\in V_t$, the inverse triangle inequality gives \[
\lVert\widetilde B_tw\rVert
\ge\lVert B_tw\rVert-\lVert(\widetilde B_t-B_t)w\rVert
\ge(c-\delta)\lVert w\rVert_g.
\] whose predicate chain is \[
\text{\shortstack[c]{\scriptsize Original\\[-1pt]\scriptsize observation\\[-1pt]\scriptsize lower bound} }c\land\text{\shortstack[c]{\scriptsize Derivative-difference\\[-1pt]\scriptsize bound} }\delta<c
\Longrightarrow\text{\shortstack[c]{\scriptsize Alternative\\[-1pt]\scriptsize observation\\[-1pt]\scriptsize lower bound} }c-\delta.
\]. Substituting this lower bound and the integral remainder bound of the alternative observation itself into the proof of the proposition ~ \ref{hgov:geo:prop:distinguish}, we can get the finite reading conclusion. The entire derivation uses operator differences only on the specified $V_t$; therefore changes in unobserved directions are not mixed into this lower bound.\end{proof}


\begin{remark}
[Matrix verification of one-dimensional observable directions]Take the Euclidean tangent spaces $\mathbb R^2$, $V=\operatorname{span}(e_1)$, $B=(1\ \ 0)$ and $\widetilde B=(1+u\ \ v)$. On $V$ there are $c=1$ and $\delta=|u|$; if $|u|<1$, then \[
\lVert\widetilde B(te_1)\rVert=|1+u|\,|t|
\ge(1-|u|)|t|.
\] $v$ does not enter this restricted operator norm. If the mission additionally requires the restoration of the $e_2$ direction, the lower bound needs to be rechecked for the expanded direction space. For affine observation and linear parameter variation $z(\varepsilon)=z_0+\varepsilon w$, $Y''(\varepsilon)=0$, the Taylor remainder bound of this example can be zero.\end{remark}


\section { Disturbance amplitude selection and distinguishable threshold }\label{exp:geo:section:amplitude}


\begin{proposition}
[Optimal certificate budget for finite perturbation amplitude]\label{exp:geo:prop:amplitude-budget}
Fixed the time $t>0$ in the common valid domain, recorded as \[
a=\frac{c\lVert v\rVert_g}{\sqrt\kappa}\sinh(\sqrt\kappa t)>0,
\qquad L=L_{\mathcal O}(t)\ge0.
\]. Suppose that for all $0<|\varepsilon|\le\varepsilon_0$, there is a unified error budget $E\ge0$, so that $e_0(t)+e_\varepsilon(t)+\eta(t)\le E$. Set $Q(r)=ar-Lr^2/2$. The existence of amplitude $r\in(0,\varepsilon_0]$ makes the necessary and sufficient condition of $Q(r)>E$ \[
Q(r_*)>E,\qquad
r_*=
\begin{cases}
\varepsilon_0,&L=0,\\
\min\{\varepsilon_0,a/L\},&L>0.
\end{cases}
\]. The amplitude that satisfies this condition gives the proposition ~ \ref{hgov:geo:prop:distinguish}'s distinction certificate. When $L>0$, the set of feasible amplitudes can also be written as \[
(0,\varepsilon_0]\cap
\left(\frac{a-\sqrt{a^2-2LE}}L,
      \frac{a+\sqrt{a^2-2LE}}L\right),
\] where this expression requires $a^2-2LE>0$; the discriminant non-positive set is empty.\end{proposition}
\begin{certificate}
[Unified budget, concave quadratic and interval constraints]Save the source of $a,L,E,\varepsilon_0$, the uniformity of the error budget over the entire amplitude interval, $Q'(r)=a-Lr$, $Q''(r)=-L$, and the common valid domain where the selected $r$ is located. This optimization is for the proven lower bound $Q$; actual readings still use the same observation version and accuracy protocol.\end{certificate}
\begin{proof}
If $L=0$, $Q(r)=ar$ is strictly increasing, so its maximum value on $[0,\varepsilon_0]$ is $Q(\varepsilon_0)$. If $L>0$, the derivative $a-Lr$ is positive on the left side of $a/L$ and negative on the right side, so the maximum value of the interval is reached at $r_* =\min\{\varepsilon_0,a/L\}$. $a>0$ and $\varepsilon_0>0$ guarantee $r_*>0$. Thus \[
\exists r\in(0,\varepsilon_0]:Q(r)>E
\Longleftrightarrow\max_{[0,\varepsilon_0]}Q>E
\Longleftrightarrow Q(r_*)>E.
\] if $Q(r)>E\ge e_0+e_\varepsilon+\eta$ , then $Q(r)-e_0-e_\varepsilon>\eta$ , is exactly the required strict threshold condition. When $L>0$, complete the square to get \[
Q(r)-E=\frac{a^2-2LE}{2L}
-\frac L2\left(r-\frac aL\right)^2.
\]. This formula is valid and only when the discriminant is positive and $r$ strictly falls between the two roots; then intersect with the allowable range interval to obtain the set listed.\end{proof}


\begin{remark}
[a numeric certificate instance with a non-zero remainder]Take $\kappa=1$ and $c\lVert v\rVert_g=1$, and observe $t=\operatorname{arsinh}3$ at the verified and valid time, then $a=3$. Assume $L=2$, $\varepsilon_0=1$, and $\eta=1/2$, and the sum of the two measurement errors does not exceed $1/2$ in the range, so $E=1$ can be taken. At this time, $r_*=1$ and $Q(r_*)=3-1=2>1$, the selected amplitude satisfies the sufficient condition of \[
\lVert\widehat y_\varepsilon-\widehat y_0\rVert
\ge2-\frac12=\frac32>\eta.
\] and is $((3-\sqrt5)/2,1]$. This example uses a negative curvature variational lower bound, a second-order remainder, and a nonzero measurement error simultaneously, giving parameter choices that can be directly reviewed.\end{remark}


\begin{remark}
[Observe the task field of certificate]The above certificate preserves the distinguishability of geometric perturbations under specified observations and precision. The candidate's task legality, obstacle fillability, and goal establishment after execution are registered by the corresponding task predicate, complex, and execution certificate respectively. When entering these interfaces from observations, maintain common object identifiers, parameter ranges, and observation versions.\end{remark}


\chapter{Fiber-Bundle Connections, Heterogeneous Transport, and Geometric Cochain Interfaces}
\label{exp:ch:transport}


This chapter converts the connection data continuation of geometric sectors into composite heterogeneous transport, and constructs cochain data that can enter homology calculations. The required exchange conditions are given at the local coordinate, path, subbundle and integral levels respectively. Each cochain construction registers a specific space, degree and differential; the curvature form resulting from the connection change enters the scalar cochain through the trace and Stokes formula.

\section { connection matrix, covariant external derivatives and curvature calculation }\label{hgov:geo:transport}


\begin{definition}
[Covariant external derivatives and curvature effects]\label{hgov:geo:def:exterior}
Let $E\to M$ be a real or complex vector bundle, and $\nabla$ be a smooth connection. Denote $\Omega^k(M,E)$ as the $E$ value $k$ sub-differential form space. The covariant external derivative $d_\nabla:\Omega^k(M,E)\to\Omega^{k+1}(M,E)$ is written as \[
d_\nabla\xi=d\xi+A\wedge\xi,
\qquad F_\nabla=dA+A\wedge A.
\] in the local frame, where $A$ is a matrix-valued form, and the wedge product contains the matrix composite at the coefficient end. Specifically, under the local frame $e=(e_1,\ldots,e_r)$ with rank $r$, the section is written as $s=e u$, $u$ is a column vector, then \[
\nabla(eu)=e(du+Au),\qquad
A=\sum_iA_i\,dx^i,
\qquad
F_\nabla=\frac12\sum_{i,j}(F_\nabla)_{ij}\,dx^i\wedge dx^j.
\] has a value of $E$ $k$ form, the effect of matrix one form is \[
(A\wedge\xi)(X_0,\ldots,X_k)
=\sum_{i=0}^k(-1)^iA(X_i)
\xi(X_0,\ldots,\widehat X_i,\ldots,X_k).
\]. The vector field of the coordinate frame satisfies $[\partial_i,\partial_j]=0$; the curvature effect of the general vector field is registered as $F_\nabla(X,Y)=[\nabla_X,\nabla_Y]-\nabla_{[X,Y]}$. If the new frame is $e'=e g$, the coefficient vector is $u'=g^{-1}u$. $A'=g^{-1}Ag+g^{-1}dg$ can be read from $\nabla(e g u')=e(dg\,u'+g\,du'+Ag\,u')$. For the $\operatorname{End}(E)$ value $k$ of the minor form $\Theta$, the induced covariant external derivative is \[
d_{\nabla^{\operatorname{End}}}\Theta
=d\Theta+A\wedge\Theta-(-1)^k\Theta\wedge A.
\]\end{definition}


\begin{theorem}
[Local commutation relation and covariant square of curvature]\label{hgov:geo:thm:curvature}
Write $\nabla_i=\partial_i+A_i$ in local coordinates. For the smooth section $s$ and any $E$ differential form $\xi$ , there are \begin{align}
\label{hgov:geo:eq:commutator}
[\nabla_i,\nabla_j]s
&=(\partial_iA_j-\partial_jA_i+[A_i,A_j])s
=(F_\nabla)_{ij}s,\\
\label{hgov:geo:eq:covariant-square}
d_\nabla^2\xi&=F_\nabla\wedge\xi,\\
\label{hgov:geo:eq:bianchi}
d_{\nabla^{\operatorname{End}}}F_\nabla&=0.
\end{align} These equations are consistent between all local frames \cite{ExtKobayashiNomizu,ExtBottTu} .\end{theorem}


\begin{certificate}
[The local multiplication law and degree are compatible with the frame]Save the coordinate plot, local frame, connection matrix $A_i$, formal degree and $d(\alpha\wedge\beta)=d\alpha\wedge\beta+(-1)^{\deg\alpha}\alpha\wedge d\beta$. When changing the frame, the transformation rules of $A'=g^{-1}Ag+g^{-1}dg$, $F'=g^{-1}Fg$ and coefficient vectors are saved. The certificate consists of the following term-by-term expansions and consistent identities.\end{certificate}


\begin{proof}
Expanding $\nabla_i\nabla_js$ gives
\[
\partial_i\partial_js+(\partial_iA_j)s
+A_j\partial_is+A_i\partial_js+A_iA_js.
\].
Interchange $i,j$ and subtract. The second mixed derivatives and first derivatives of the section cancel, leaving precisely equation ~ \eqref{hgov:geo:eq:commutator}. Next, using the degree of the one-form $A$, compute
\[
\begin{aligned}
d_\nabla^2\xi
&=d(d\xi+A\wedge\xi)+A\wedge(d\xi+A\wedge\xi)\\
&=dA\wedge\xi-A\wedge d\xi
+A\wedge d\xi+A\wedge A\wedge\xi\\
&=(dA+A\wedge A)\wedge\xi.
\end{aligned}
\].
For the curvature,
\[
dF=dA\wedge A-A\wedge dA,
\qquad
A\wedge F-F\wedge A=A\wedge dA-dA\wedge A,
\].
Adding gives equation ~ \eqref{hgov:geo:eq:bianchi}. To verify frame compatibility, write $a=g^{-1}dg$ and $\widetilde A=g^{-1}Ag$. From $d(g^{-1})=-g^{-1}(dg)g^{-1}$ and the multiplication rule for one-forms, we obtain
\[
\begin{aligned}
da&=-a\wedge a,\\
d\widetilde A&=-a\wedge\widetilde A+g^{-1}(dA)g-\widetilde A\wedge a,\\
d(\widetilde A+a)+(\widetilde A+a)\wedge(\widetilde A+a)
&=g^{-1}(dA+A\wedge A)g.
\end{aligned}
\].
The first equation eliminates $da+a\wedge a$. In the second, the two cross terms cancel those in the square expansion. Thus $F'=g^{-1}Fg$; applying it to $u'=g^{-1}u$ gives $F'u'=g^{-1}Fu$. Consequently, the covariant square and Bianchi identity agree between local frames. The local calculations therefore define global bundle-valued equalities, with continuation
\[
\text{\shortstack[c]{\scriptsize Connection\\[-1pt]\scriptsize matrix and\\[-1pt]\scriptsize degree}}
\Longrightarrow d_\nabla^2=F_\nabla\wedge(-)
\Longrightarrow\text{\shortstack[c]{\scriptsize Frame-compatible\\[-1pt]\scriptsize curvature action}}.
\].
\end{proof}

\section { Heterogeneous pullback, parallel transportation and path recording }\label{exp:geo:section:pullback}


\begin{definition}
[Keep connection heterogeneous transport]\label{hgov:geo:def:heterogeneous}
The transport between the two sectors $\alpha,\beta$ is given by the smooth basis map $f_{\alpha\beta}:M_\alpha\to M_\beta$ and the smooth linear bundle map \[
T_{\alpha\beta}:E_\alpha\longrightarrow f_{\alpha\beta}^*E_\beta
\]. For each $s\in\Gamma(E_\alpha)$, it is required that the fibers in each sector of \begin{equation}
\label{hgov:geo:eq:connection-intertwining}
(f_{\alpha\beta}^*\nabla^\beta)(T_{\alpha\beta}s)
=(\operatorname{id}_{T^*M_\alpha}\otimes T_{\alpha\beta})(\nabla^\alpha s).
\end{equation} can have different dimensions; the transport record maintains the conditions for its nuclei, images and mission observations. If the task requires reversible transport, add the bundle inverse and the bilateral composite identity to the record.\end{definition}


\begin{proposition}
[Pull back the connection with the coordinates of the heterogeneous transport certificate]\label{exp:geo:prop:pullback-coordinates}
In the source coordinate $x^1,\ldots,x^n$ and the target coordinate $y^1,\ldots,y^m$, let $f=(f^1,\ldots,f^m)$ and the target connection be written as $\nabla^\beta=d+\sum_aA_a^\beta(y)\,dy^a$. In the pull-back frame, the coefficients and curvature of $f^*\nabla^\beta=d+\widetilde A$ satisfy \begin{align}
\label{exp:geo:eq:pullback-connection}
\widetilde A_i(x)&=\sum_aA_a^\beta(f(x))\partial_i f^a(x),\\
\label{exp:geo:eq:pullback-curvature}
\widetilde F_{ij}(x)&=\sum_{a,b}F^\beta_{ab}(f(x))
\partial_i f^a(x)\partial_j f^b(x).
\end{align}. If the source connection is $d+A^\alpha$ and the local matrix transporting $T$ is $H(x)$, then the necessary and sufficient conditions for maintaining connection are \begin{equation}
\label{exp:geo:eq:matrix-intertwining}
\partial_iH+\widetilde A_iH-HA_i^\alpha=0
\qquad(i=1,\ldots,n).
\end{equation}. Among them, $H$ can be a rectangular matrix corresponding to different source and target ranks.\end{proposition}
\begin{certificate}
[Chain rule, mixed partial derivatives and matrix exchange]Save $f$ , the source and target frames, $A_i^\alpha,A_a^\beta$ , $H$ , and the row and column dimensions of each matrix. Save $\partial_i\partial_jf^a=\partial_j\partial_if^a$ for each $i,j$, and verify the derivative terms and commutative subterms of the formula ~ \eqref{exp:geo:eq:pullback-curvature} respectively. To maintain connection, the identity of the formula ~ \eqref{exp:geo:eq:matrix-intertwining} in the common coordinate domain is established as certificate.\end{certificate}
\begin{proof}
The pullback-form satisfies $f^*(dy^a)=df^a=\sum_i\partial_if^a\,dx^i$, so the pullback of the target connection-form directly gives the formula ~ \eqref{exp:geo:eq:pullback-connection}. Differentiating term by term gives \[
\partial_i\widetilde A_j
=\sum_{a,b}(\partial_bA_a^\beta)(f(x))
\partial_if^b\partial_jf^a
+\sum_aA_a^\beta(f(x))\partial_i\partial_jf^a.
\] exchange $i,j$ After subtraction, the second term disappears by the mixed partial derivative exchange. Rename the summation index, and the first term becomes \[
\sum_{a,b}(\partial_aA_b^\beta-\partial_bA_a^\beta)(f(x))
\partial_if^a\partial_jf^b.
\]. The matrix commutator is also given as \[
[\widetilde A_i,\widetilde A_j]
=\sum_{a,b}[A_a^\beta,A_b^\beta](f(x))
\partial_if^a\partial_jf^b.
\]. The addition of the two formulas is the formula ~ \eqref{exp:geo:eq:pullback-curvature}. The global expression is $F_{f^*\nabla^\beta}=f^*F_{\nabla^\beta}$.

For any source section coefficient $u$, the two equations \[
\begin{aligned}
(d+\widetilde A)(Hu)&=(dH)u+H\,du+\widetilde A H u,\\
H(d+A^\alpha)u&=H\,du+H A^\alpha u.
\end{aligned}
\] are equal to all $u$, if and only if $(dH+\widetilde A H-H A^\alpha)u=0$ is true for all $u$, which is in turn if and only if the equation ~ \eqref{exp:geo:eq:matrix-intertwining} was established. The entire calculation uses only matrix multiplications that are type-compatible and therefore suitable for heterogeneous transports of different ranks.\end{proof}


\begin{theorem}
[Curvature,paths for heterogeneous transport are compatible with finite,composites.]\label{hgov:geo:thm:transport}
For transportation in the definition ~ \ref{hgov:geo:def:heterogeneous}, let $\gamma:[0,1]\to M_\alpha$ be a piecewise smooth path. The corresponding parallel transportation is recorded as $\operatorname{PT}$, then the transportation of \begin{align}
\label{hgov:geo:eq:curvature-intertwining}
F_{f_{\alpha\beta}^*\nabla^\beta}\,T_{\alpha\beta}
&=(\operatorname{id}_{\Lambda^2T^*M_\alpha}\otimes T_{\alpha\beta})F_{\nabla^\alpha},\\
\label{hgov:geo:eq:parallel-intertwining}
\operatorname{PT}^{\beta}_{f_{\alpha\beta}\circ\gamma}
\circ(T_{\alpha\beta})_{\gamma(0)}
&=(T_{\alpha\beta})_{\gamma(1)}
\circ\operatorname{PT}^{\alpha}_{\gamma}.
\end{align} that maintains connection can be made finite composite; the composite transportation still satisfies the above two equations. If each step also satisfies the observation exchange $\mathcal O_\beta(f_{\alpha\beta}(x),(T_{\alpha\beta})_xv)=\mathcal O_\alpha(x,v)$ in the same task range, then any finite composition maintains the task observation.\end{theorem}


\begin{certificate}
[Complex mapping, covariant exchange and path certificate]Save the type identification of the base map, bundle map, pullback bundle and the formula ~ \eqref{hgov:geo:eq:connection-intertwining}. The path certificate stores the starting and ending points, the connection coefficients on each smooth segment, and the linear initial value problem. The composite certificate holds each intermediate sector with the natural identification $f_{\alpha\beta}^*f_{\beta\gamma}^*E_\gamma
\cong(f_{\beta\gamma}\circ f_{\alpha\beta})^*E_\gamma$ . The observation holds the certificate as well as the input type and common range of the observation.\end{certificate}


\begin{proof}
Expand $T_{\alpha\beta}$ to the bundle value form according to the coefficient function. The multiplication method of the connection commutation relation and the external differential gives \[
d_{f_{\alpha\beta}^*\nabla^\beta}(T_{\alpha\beta}\xi)
=T_{\alpha\beta}(d_{\nabla^\alpha}\xi).
\]. Act on the section again and use the formula ~ \eqref{hgov:geo:eq:covariant-square} to get the curvature commutation relation. Given the initial value $v$, let $u(t)$ be the parallel section along $\gamma$. In a coordinate segment, the parallel equation and the basic solution matrix are \[
u'(t)+A^\alpha_{\gamma(t)}(\dot\gamma(t))u(t)=0,
\qquad
U'(t)=-A^\alpha_{\gamma(t)}(\dot\gamma(t))U(t),\quad U(0)=I.
\] respectively. Let $\widetilde u(t)=H(\gamma(t))u(t)$ . Use the formula ~ \eqref{exp:geo:eq:matrix-intertwining} to condense along $\gamma$, and get \[
\begin{aligned}
\widetilde u'+\widetilde A(\dot\gamma)\widetilde u
&=\bigl(dH(\dot\gamma)+\widetilde A(\dot\gamma)H
       -HA^\alpha(\dot\gamma)\bigr)u
 +H\bigl(u'+A^\alpha(\dot\gamma)u\bigr)\\
&=0.
\end{aligned}
\]. By pulling back the connection formula, $\widetilde A(\dot\gamma)=A^\beta(df\dot\gamma)$, so this equation is the target parallel equation along $f\circ\gamma$, and the initial value is $H(\gamma(0))v$. The uniqueness of the solution to the linear ordinary differential equation is given by the formula ~ \eqref{hgov:geo:eq:parallel-intertwining}. When the path crosses the frame, the coefficients are transformed according to $g^{-1}$ and the fiber vectors remain the same; the endpoints of each smooth segment are compounded sequentially to obtain the commutation relation of the entire segmented smooth path.

The composite of the two-step transport is \[
T_{\alpha\gamma}
=(f_{\alpha\beta}^*T_{\beta\gamma})\circ T_{\alpha\beta},
\qquad
f_{\alpha\gamma}=f_{\beta\gamma}\circ f_{\alpha\beta}.
\]. First use the covariant commutation relation in the first step, and then use the pull-back commutation relation in the second step to obtain the composite covariant exchange. Taking the identity transport as a basic example and generalizing the number of steps, all finite compositions maintain connection and inherit curvature and path exchange. The same successive substitutions are made for the observation exchange, and \[
\mathcal O_{\alpha_n}\bigl(f_{\alpha_0\alpha_n}(x),
(T_{\alpha_0\alpha_n})_xv\bigr)
=\mathcal O_{\alpha_0}(x,v),
\] is obtained. The base mapping of each compound is understood according to the above formula, so the task observation is maintained.\end{proof}


\begin{proposition}
[Overall transport and coefficient cohomology of a flat circumferential connection]\label{exp:geo:prop:circle-holonomy}
Let $\eta=x\,dy-y\,dx$ be defined on the unit circle $S^1\subset\mathbb R^2$, and take the trivial polylinear bundle $E=S^1\times\mathbb C$. Given $\lambda\in\mathbb R$ , connection \[
\nabla_\lambda=d+i\lambda\eta
\] satisfies $F_{\nabla_\lambda}=0$ . A path with an angle lift of $\theta(t)$ has a parallel transport of \[
z(1)=\exp\bigl(-i\lambda(\theta(1)-\theta(0))\bigr)z(0).
\] and a closed path with a winding number of $n$ has an overall transport of $e^{-2\pi i\lambda n}$. Furthermore, the cohomology of a flat cochain complex is \[
H^k\bigl(\Omega^\bullet(S^1,E),d_{\nabla_\lambda}\bigr)
\cong
\begin{cases}
\mathbb C,&\lambda\in\mathbb Z\text{ and }k=0,1,\\
0,&\lambda\notin\mathbb Z\text{ or }k\ge2.
\end{cases}
\] where $k$ takes a non-negative integer.\end{proposition}
\begin{certificate}
[Angle lifting, periodic conditions and first-order equations]Save the global form $\eta$, the local identity $\eta=d\theta$, the path angle lift and its end point difference. The periodic function representation uses $\theta\sim\theta+2\pi$. The cohomology calculation saves the differential operator $D_\lambda u=u'+i\lambda u$ and the integral $\int_0^{2\pi}e^{i\lambda s}f(s)\,ds$, and verifies the zero-order kernel and the first-order quotient space.\end{certificate}
\begin{proof}
$S^1$ is a one-dimensional manifold, so $d\eta=0$ has a form $\eta\wedge\eta=0$. So $F_{\nabla_\lambda}=i\lambda d\eta+(i\lambda\eta)\wedge(i\lambda\eta)=0$. The parallel equation along the path is $z'+i\lambda\theta'z=0$ and when multiplied by the integration factor $e^{i\lambda\theta(t)}$ the derivative is zero, resulting in the parallel transport listed. The angular endpoint difference of the closed path is $2\pi n$, so its overall transportation is determined only by $n$ and $\lambda$. When $\lambda\notin\mathbb Z$, the overall transport around number one is not identical.

To calculate cohomology, write the zero form as the smooth periodic function $u(\theta)$ and the one form as $f(\theta)d\theta$. Its differential is \[
d_{\nabla_\lambda}u=(u'+i\lambda u)d\theta,
\qquad\Omega^k(S^1,E)=0\quad(k\ge2).
\] and the general solution of $D_\lambda u=0$ is $u=c e^{-i\lambda\theta}$. Its periodic condition is $(e^{-2\pi i\lambda}-1)c=0$, so the zero-order kernel is one-dimensional when it is an integer $\lambda$, and it is zero when it is a non-integer.

For any smooth period $f$, the general solution of the equation $D_\lambda u=f$ is \[
u(\theta)=e^{-i\lambda\theta}
\left(c+\int_0^\theta e^{i\lambda s}f(s)\,ds\right).
\]. The periodic condition is equivalent to \[
(e^{2\pi i\lambda}-1)c
=\int_0^{2\pi}e^{i\lambda s}f(s)\,ds.
\]. If $\lambda\notin\mathbb Z$, the left-hand coefficient is non-zero, so each $f$ has a unique periodic solution, differential surjective , once the cohomology is zero. If $\lambda\in\mathbb Z$, the solvable condition is that the right-hand integral is equal to zero. This integral is a nonzero linear functional because it is equal to $2\pi$ when acting on the periodic function $f=e^{-i\lambda\theta}$. Therefore, the differential image is the kernel of this functional, and its quotient is one dimension. One time representation is desirable $e^{-i\lambda\theta}d\theta$. The computation of these kernels, images, and quotients gives the full cohomology and preserves the effect of the global path data on the flat coefficient complex.\end{proof}


\section { Flat subbundle and coefficient cochain complex }\label{hgov:geo:complexes}


\begin{proposition}
[cochain construction of flat coefficient bundles]\label{hgov:geo:prop:flat-complex}
Let $K\subseteq E$ be a smooth subbundle that satisfies $\nabla\Gamma(K)\subseteq\Omega^1(M,K)$ and $F_\nabla|_K=0$. Then the restriction connection $\nabla^K$ gives the cochain complex \[
\bigl(\Omega^\bullet(M,K),d_{\nabla^K}\bigr),
\qquad d_{\nabla^K}^2=0.
\]. On the entire $E$, the necessary and sufficient conditions of $d_\nabla^2=0$ are $F_\nabla=0$. If the heterogeneous transport maintains the selected sub-plexes and their connections, it induces cochain mapping between these cochain complexes after pulling back.\end{proposition}


\begin{certificate}
[Subbundle stability, curvature elimination and commutation with differentials]Save the local subframe, connection stable form, zero effect of curvature on $K$ and all formal degrees of $K$. The transportation part saves $T(K)\subseteq f^*K'$ and the covariant commutation relation. The mapping between zero square and cochain is verified successively by the formula ~ \eqref{hgov:geo:eq:covariant-square}.\end{certificate}


\begin{proof}
connection stability guarantee $d_{\nabla^K}$ is defined on the listed space. For any $K$ value of the form $\xi$, the theorem ~ \ref{hgov:geo:thm:curvature} gives $d_{\nabla^K}^2\xi=F_\nabla\wedge\xi=0$, so we get the cochain complex. Over the entire plexus, $F_\nabla=0$ also follows that the square is zero. On the contrary, if the square is zero on all sections, then for each local smooth section $s$ there is $F_\nabla s=0$; randomly selecting a local section with a given fiber value at each point, we get that $F_\nabla$ has zero effect on all fiber vectors, so $F_\nabla=0$. The transport that keeps connection satisfies $d_{f^*\nabla^{K'}}T=T d_{\nabla^K}$ , which is exactly the commutation relation required for cochain mapping.\end{proof}


\begin{remark} [An instance of flat coefficients within nonzero total curvature]
In the trivial rank-two real vector bundle over $M=\mathbb R^2$, take
\[
A=\begin{pmatrix}0&0\\0&x\,dy\end{pmatrix},
\qquad
F_\nabla=\begin{pmatrix}0&0\\0&dx\wedge dy\end{pmatrix}.
\].
Let $K$ be the subbundle generated by the first coordinate vector $e_1$. For every function $u$, we have $\nabla(u e_1)=du\,e_1$. Thus $K$ is connection-stable, and the curvature vanishes on $K$. The corresponding $K$-valued form is $(\alpha,0)$, with differential and square
\[
d_\nabla(\alpha,0)=(d\alpha,0),\qquad
 d_\nabla^2(\alpha,0)=0.
\].
The inclusion $T:u\mapsto(u,0)$ satisfies $d_\nabla T=T d$, providing a coefficient choice that can enter cochain computations directly. On the entire bundle, $d_\nabla^2e_2=(dx\wedge dy)e_2$. Closedness in degree zero can also be checked explicitly: if $s=(u,v)$ and $\nabla s=0$, then $du=0$ and $dv+xv\,dy=0$. Applying the covariant differential again gives $v\,dx\wedge dy=0$; hence $v=0$, while $u$ is constant on the connected $\mathbb R^2$. These equalities jointly verify subbundle stability, zero curvature action, and explicit closed elements.
\end{remark}

\section { Directed chain, Stokes formula and finite cochain interface }\label{exp:geo:section:stokes}


\begin{definition}
[Smooth singular chain and integral cochain]\label{hgov:geo:def:singular}
For the smooth manifold $M$ and the coefficient domain $\mathbb F\in\{\mathbb R,\mathbb C\}$, for $k\ge0$, let $C_k^{\mathrm{sm}}(M;\mathbb F)$ be the free vector space generated by the smooth singular $k$ simplex; the smooth singular simplex adopts a mapping that can be smoothly extended to the standard simplex neighborhood. Set $C_k^{\mathrm{sm}}=0$ to $k<0$ and set the boundary of the zero chain to zero. For the simplex $\sigma:\Delta^k\to M$ of $k\ge1$, define the boundary; at the same time, define the dual cochain in all non-negative degrees: \[
\partial\sigma=\sum_{i=0}^k(-1)^i\sigma\circ\delta_i,
\qquad
C_{\mathrm{sm}}^k(M;\mathbb F)
=\operatorname{Hom}_{\mathbb F}(C_k^{\mathrm{sm}}(M;\mathbb F),\mathbb F),
\qquad
(\delta c)(\sigma)=c(\partial\sigma),
\] where $\delta_i$ is the $i$th directed surface embedding. The integration map in scalar form is \[
I_k:\Omega^k(M;\mathbb F)\longrightarrow C_{\mathrm{sm}}^k(M;\mathbb F),
\qquad I_k(\theta)(\sigma)=\int_{\Delta^k}\sigma^*\theta.
\]\end{definition}


\begin{theorem}
[Exchange of scalar differential forms to singular cochains]\label{hgov:geo:thm:integration}
Definition ~ \ref{hgov:geo:def:singular} gives chain complex and its dual cochain complex, and \[
\partial^2=0,\qquad\delta^2=0,\qquad
\delta I_k=I_{k+1}d.
\]. Therefore every closed form $\theta$ yields closed cochain $I_k\theta$; every full differential $d\zeta$ yields coboundary $\delta I_{k-1}\zeta$ \cite{ExtBottTu,ExtHatcher}.\end{theorem}


\begin{certificate}
[Directed offset and Stokes certificate]Preserves the vertex order, face embedding, formal smoothness, and integral orientation of the simplex. Each codimension-two surface retains two occurrences of the symbol; integral exchange preserves surface constraints and Stokes' equations. A certificate on a finite chain consists of simplex certificates linearly superposed.\end{certificate}


\begin{proof}
For $k\ge2$, each face with two deleted vertices $i<j$ appears exactly twice in $\partial^2\sigma$, and the coefficients are $(-1)^{i+j-1}$ and $(-1)^{i+j}$ respectively. The two terms cancel. The low-degree square is defined by the zero boundary under the zero chain, so $\partial^2=0$. So for each cochain $c$ and chain $a$, there is $(\delta^2c)(a)=c(\partial^2a)=0$. For the smooth $(k+1)$ simplex $\sigma$, the Stokes formula gives \[
\begin{aligned}
(\delta I_k\theta)(\sigma)
&=\sum_{i=0}^{k+1}(-1)^i
\int_{\Delta^k}(\sigma\circ\delta_i)^*\theta\\
&=\int_{\partial\Delta^{k+1}}\sigma^*\theta
=\int_{\Delta^{k+1}}d(\sigma^*\theta)
=\int_{\Delta^{k+1}}\sigma^*(d\theta)
=(I_{k+1}d\theta)(\sigma).
\end{aligned}
\], which proves the commutation relation; the closure and coboundary conclusions are obtained by substituting $d\theta=0$ or $\theta=d\zeta$.\end{proof}


\begin{proposition} [Naturality of integration and finite chain realizations]
\label{exp:geo:prop:integration-naturality}
For a smooth map $f:M\to N$, define
$f_\#\sigma=f\circ\sigma$
and extend it linearly to smooth singular chains. Then
\[
\partial f_\#=f_\#\partial,
\qquad
I_M(f^*\theta)=f_\#^*(I_N\theta),
\],
where $f_\#^*c=c\circ f_\#$. Furthermore, suppose $(K_\bullet,b)$ is a chain complex of finite-dimensional vector spaces, and $j_k:K_k\to C_k^{\mathrm{sm}}(M;\mathbb F)$ satisfies
$\partial j_k=j_{k-1}b_k$.
On $K^k=\operatorname{Hom}_{\mathbb F}(K_k,\mathbb F)$ set
$\delta_K=b^*$
and define
\[
I_{K,k}(\theta)(c)=\int_{j_k(c)}\theta.
\].
Then
\[
\delta_K I_{K,k}=I_{K,k+1}d.
\].
In particular, closed forms yield finite closed cochains, and exact forms integrate to zero on every finite cycle $c\in\ker b$.
\end{proposition}
\begin{certificate} [Composition of simplices, boundary matrices, and commutation of realizations]
Retain $f$, the vertex ordering of simplices, the chain bases, and the boundary matrices $b_k$. For each basis vector $c$, the finite-realization certificate checks $\partial j_k(c)=j_{k-1}(b_kc)$ and additionally checks $b_{k-1}b_k=0$. Dual cochains use the same bases, so the differential matrix from $K^{k-1}$ to $K^k$ is $b_k^{\mathsf T}$. The orientation of integration agrees with the Stokes certificate of Theorem ~ \ref{hgov:geo:thm:integration}.
\end{certificate}
\begin{proof}
For any simplex $\sigma$, associativity of map composition gives
\[
\partial(f\circ\sigma)
=\sum_i(-1)^if\circ\sigma\circ\delta_i
=f_\#(\partial\sigma).
\].
On $N$, for a $k$-form $\theta$, we have
\[
\begin{aligned}
I_M(f^*\theta)(\sigma)
&=\int_{\Delta^k}\sigma^*(f^*\theta)
=\int_{\Delta^k}(f\circ\sigma)^*\theta\\
&=(I_N\theta)(f_\#\sigma)
=(f_\#^*I_N\theta)(\sigma).
\end{aligned}
\].
Both equalities hold for arbitrary chains after linear extension.

For $c\in K_{k+1}$, the exchange relation realized by finite and the Stokes formula give \[
\begin{aligned}
(\delta_K I_{K,k}\theta)(c)
&=I_{K,k}\theta(b_{k+1}c)
=\int_{j_k(b_{k+1}c)}\theta\\
&=\int_{\partial j_{k+1}(c)}\theta
=\int_{j_{k+1}(c)}d\theta
=(I_{K,k+1}d\theta)(c).
\end{aligned}
\]. If $d\theta=0$, this formula directly gives the closure property. If $\theta=d\zeta$ and $b c=0$ , then $I_K(\theta)(c)=(\delta_KI_K\zeta)(c)=I_K\zeta(bc)=0$ . Its verifiable continuation is \[
\text{\shortstack[c]{\scriptsize Directed\\[-1pt]\scriptsize geometric\\[-1pt]\scriptsize realization}}\land(\partial j=jb)
\Longrightarrow\text{\shortstack[c]{\scriptsize Integration\\[-1pt]\scriptsize commutes with\\[-1pt]\scriptsize differentiation}}
\Longrightarrow\text{\shortstack[c]{\scriptsize Finite closed\\[-1pt]\scriptsize cochains and\\[-1pt]\scriptsize period pairings}}.
\]\end{proof}


\begin{remark}
[Complete matrix calculation for a single triangle]Take the vertices $v_0=(0,0)$, $v_1=(1,0)$, and $v_2=(0,1)$, and take the basis in the order of the edge $[01],[02],[12]$ and the face $[012]$. The boundary matrix is \[
b_1=\begin{pmatrix}-1&-1&0\\1&0&-1\\0&1&1\end{pmatrix},
\qquad b_2=\begin{pmatrix}1\\-1\\1\end{pmatrix},
\qquad b_1b_2=0.
\]. Let $j$ be the affine implementation of the standard simplex to the triangle and its edges and vertices. For $\alpha=x\,dy$, the integrals on the three sides are $0,0,1/2$: The last side is $x=1-s,y=s$, so its integral is $\int_0^1(1-s)ds=1/2$. Therefore \[
I_{K,1}\alpha=\begin{pmatrix}0\\0\\1/2\end{pmatrix},
\qquad
b_2^{\mathsf T}I_{K,1}\alpha=\frac12
=\int_{[012]}dx\wedge dy=I_{K,2}(d\alpha).
\] boundary cancellation, differential matrix and area integral correspond item by item in this finite instance.\end{remark}


\section { Curvature scalar closed form and connection deformation }\label{exp:geo:section:curvature-classes}


\begin{proposition}
[Closed cochain continuation of curvature characteristic form]\label{hgov:geo:prop:curvature-cocycle}
For any real or complex vector bundle connection $\nabla$ and the integer $m\ge1$, let \[
\chi_m(\nabla)=\operatorname{tr}
\bigl(F_\nabla^{\wedge m}\bigr)\in\Omega^{2m}(M;\mathbb F),
\qquad
\omega_m=I_{2m}(\chi_m(\nabla)).
\] then $\chi_m(\nabla)$ is globally defined and satisfies \[
d\chi_m(\nabla)=0,\qquad\delta\omega_m=0.
\]. This leads to an unambiguous construction from geometric data with registered connections to scalar closed cochains.\end{proposition}


\begin{certificate}
[Curvature polynomials, traces and closure properties]Saves the coefficient field of the bundle, connection, curvature, integer $m$, frame transformation, Bianchi identity, graded cyclic properties of the trace, and integration map. The degree $2m$ is registered together with the degree of subsequent cochain differentiation.\end{certificate}


\begin{proof}
$F'=g^{-1}Fg$ under frame transformation, so $(F')^{\wedge m}=g^{-1}F^{\wedge m}g$; unchanged after tracing, so $\chi_m$ is in global scalar form. The induced covariant external derivative is a graded derivative. Bianchi's identity gives \[
d_{\nabla^{\operatorname{End}}}(F^{\wedge m})
=\sum_{j=0}^{m-1}
F^{\wedge j}\wedge
(d_{\nabla^{\operatorname{End}}}F)
\wedge F^{\wedge(m-1-j)}=0.
\] to $\operatorname{End}(E)$ value $q$ in the form $\Theta$. The graded cycle property makes \[
\operatorname{tr}\bigl(A\wedge\Theta-(-1)^q\Theta\wedge A\bigr)=0,
\qquad
\operatorname{tr}(d_{\nabla^{\operatorname{End}}}\Theta)
=d\operatorname{tr}(\Theta).
\]. Substituting $\Theta=F^{\wedge m}$, we get $d\chi_m=0$. Finally, the theorem ~ \ref{hgov:geo:thm:integration}, $\delta\omega_m=\delta I_{2m}\chi_m=I_{2m+1}d\chi_m=0$.\end{proof}


\begin{proposition}
[connection variant of explicit coboundary certificate]\label{exp:geo:prop:transgression}
Given two terms connection $\nabla^0,\nabla^1$ on the same real or complex vector bundle $E\to M$, let $a=\nabla^1-\nabla^0\in\Omega^1(M,\operatorname{End}(E))$, $\nabla^t=\nabla^0+t a$, $0\le t\le1$, and note its curvature as $F_t$. For the integer $m\ge1$ , define the global scalar $(2m-1)$ of the form \[
\mathsf T_{2m-1}(\nabla^0,\nabla^1)
=m\int_0^1\operatorname{tr}\bigl(a\wedge F_t^{\wedge(m-1)}\bigr)\,dt,
\] where $F_t^{\wedge0}=\operatorname{id}_E$ . Then \begin{align}
\label{exp:geo:eq:transgression-form}
\chi_m(\nabla^1)-\chi_m(\nabla^0)&=d\mathsf T_{2m-1},\\
\label{exp:geo:eq:transgression-cochain}
I_{2m}\chi_m(\nabla^1)-I_{2m}\chi_m(\nabla^0)
&=\delta I_{2m-1}\mathsf T_{2m-1}.
\end{align} so the scalar curvature cochains given by the two terms connection represent the same cohomology class; their pairings are equal on every $2m$ smooth singular period.\end{proposition}
\begin{certificate}
[Same bundle identification, parameter derivatives and Bianchi continuation]Save the common bundle of two connections $E$, difference form $a$, parameter interval, degree $2m-1$ and $2m$, and \[
F_t=F_0+t\,d_{(\nabla^0)^{\operatorname{End}}}a+t^2a\wedge a,
\qquad
\partial_tF_t=d_{(\nabla^t)^{\operatorname{End}}}a.
\] save the fractional trace cycle formula, Bianchi identity and integral cochain commutation relation. All operations are performed on this common bundle; when using this certificate between different bundles, the corresponding bundle isomorphism and connection pullback need to be registered first. This proof gives the specific form of the curvature polynomial deformation formula \cite{ExtBottTu,ExtKobayashiNomizu}.\end{certificate}
\begin{proof}
In the common frame $A_t=A_0+t a$. Directly expand $dA_t+A_t\wedge A_t$ to get the first form of certificate. Then differentiate $t$ to get \[
\partial_tF_t
=da+A_t\wedge a+a\wedge A_t
=d_{(\nabla^t)^{\operatorname{End}}}a,
\], where the two product terms with the same sign come from $\deg a=1$. The difference connection $a$ is transformed according to the conjugate under the frame transformation, and the same is true for $F_t$, so $\operatorname{tr}(a\wedge F_t^{m-1})$ is in global scalar form.

Let $\mathcal D_t=d_{(\nabla^t)^{\operatorname{End}}}$ be the covariant external derivative on the End bundle. Since $\deg F_t=2$ is an even number, the product derivation and trace loop give \[
\begin{aligned}
\partial_t\operatorname{tr}(F_t^{\wedge m})
&=\sum_{j=0}^{m-1}
\operatorname{tr}\bigl(F_t^{\wedge j}\wedge(\partial_tF_t)
\wedge F_t^{\wedge(m-1-j)}\bigr)\\
&=m\operatorname{tr}\bigl((\mathcal D_ta)\wedge F_t^{\wedge(m-1)}\bigr).
\end{aligned}
\] and the Bianchi identity gives $\mathcal D_tF_t=0$ , so $\mathcal D_t(F_t^{m-1})=0$ . Using the fractional multiplication method and $\deg a=1$, we have \[
\begin{aligned}
d\operatorname{tr}(a\wedge F_t^{m-1})
&=\operatorname{tr}\bigl(\mathcal D_t(a\wedge F_t^{m-1})\bigr)\\
&=\operatorname{tr}\bigl((\mathcal D_ta)\wedge F_t^{m-1}
-a\wedge\mathcal D_t(F_t^{m-1})\bigr)\\
&=\operatorname{tr}\bigl((\mathcal D_ta)\wedge F_t^{m-1}\bigr).
\end{aligned}
\] so the parameter derivative is equal to $m$ times the above total differential. Integrate coordinate by coordinate on the tight parameter interval $[0,1]$, exchange external differentials with parameter integrals, and obtain \[
\int_0^1\partial_t\chi_m(\nabla^t)\,dt
=d\left(m\int_0^1\operatorname{tr}(a\wedge F_t^{m-1})dt\right),
\], which is the formula ~ \eqref{exp:geo:eq:transgression-form}. Then apply the theorem ~ \ref{hgov:geo:thm:integration} to get the formula ~ \eqref{exp:geo:eq:transgression-cochain}. If $c$ is a period, the value of the difference between the two cochains on $c$ is $(\delta I\mathsf T)(c)=(I\mathsf T)(\partial c)=0$. Therefore \[
\text{\shortstack[c]{\scriptsize Connection\\[-1pt]\scriptsize deformation on\\[-1pt]\scriptsize the same bundle}}\Longrightarrow\text{\shortstack[c]{\scriptsize Difference of\\[-1pt]\scriptsize curvature forms\\[-1pt]\scriptsize is exact}}
\Longrightarrow\text{\shortstack[c]{\scriptsize Difference of\\[-1pt]\scriptsize integration\\[-1pt]\scriptsize cochains is a\\[-1pt]\scriptsize coboundary}}
\Longrightarrow\text{\shortstack[c]{\scriptsize Period pairings\\[-1pt]\scriptsize are preserved}}.
\]\end{proof}


\begin{remark}
[Locally computable examples of first-order curvature polynomials]Taking $A_0=0$ and $A_1=x\,dy$ on the real trivial line bundle of $\mathbb R^2$, then $a=x\,dy$ and $F_t=t\,dx\wedge dy$. $m=1$ sometimes has $\mathsf T_1=x\,dy$ and $\chi_1(\nabla^1)-\chi_1(\nabla^0)=dx\wedge dy$. In the single-triangle implementation described above, the integral of this difference form is $1/2$ , which is exactly the integral of $\mathsf T_1$ on its directed boundary. For periods with zero bounds, the same commutation relation gives zero integral difference.\end{remark}


\begin{remark}
[Continuation conventions for geometry and obstacle complexes]Flat coefficient bundles, scalar differential forms and smooth singular chains respectively give the specific complexes of this chapter. The square of the general covariant external derivative is fully documented by the equation ~ \eqref{hgov:geo:eq:covariant-square}. If the alternative obstacle complex $(C_{\mathrm{obs}}^\bullet,d_{\mathrm{obs}})$ is used, \[
\Psi:C_{\mathrm{sm}}^\bullet(M;\mathbb F)
\longrightarrow C_{\mathrm{obs}}^\bullet,
\qquad d_{\mathrm{obs}}\Psi=\Psi\delta.
\] is further registered, so the closed cochain of this chapter enters the complex along $\Psi$ and satisfies $d_{\mathrm{obs}}\Psi(\omega_m)=0$. The adequacy of the selected curvature readout and mapping $\Psi$ for the task is verified by the observation factorization certificate; the existence of specific repairs is determined by the corresponding obstacle class and allowed filling elements.\end{remark}


\section { Differential factorization of geometrically observed policy changes }\label{hgov:geo:response}


\begin{theorem}
[The kernel contains the local response operator given]\label{hgov:geo:thm:factor}
Let $\mathcal Z,\mathcal W,\mathcal U$ be a finite-dimensional smooth manifold, $\mathcal K:\mathcal Z\to\mathcal W$ be a registered geometric observation, and $\Theta:\mathcal Z\to\mathcal U$ be a strategic observation. Fix $z\in\mathcal Z$, let \[
A=D\mathcal K_z:T_z\mathcal Z\to T_{\mathcal K(z)}\mathcal W,
\qquad
B=D\Theta_z:T_z\mathcal Z\to T_{\Theta(z)}\mathcal U.
\] exist a unique linear operator $\mathcal R_z:\operatorname{im}A\to T_{\Theta(z)}\mathcal U$ such that $B=\mathcal R_zA$, if and only if \begin{equation}
\label{hgov:geo:eq:kernel}
\ker A\subseteq\ker B.
\end{equation} specifies the inner product for these tangent spaces, there exists a right inverse $S:\operatorname{im}A\to T_z\mathcal Z$, satisfying \[
AS=\operatorname{id}_{\operatorname{im}A},\qquad
\mathcal R_z=BS,\qquad
\lVert\mathcal R_z\rVert\le\lVert B\rVert\lVert S\rVert.
\]\end{theorem}


\begin{certificate}
[Kernel, Image, Right Inverse and Response Upper Bound]Save the fixed point $z$, the domain of the two observations, the derivative matrix, the basis of $\ker A$ and the zero value of $B$ on the basis. Select the complementary space $V$ of $\ker A$ in $T_z\mathcal Z$, and save the reversibility and inverse matrix $S$ of $A|_V:V\to\operatorname{im}A$. The operator norm is calculated using the same set of registered inner products.\end{certificate}


\begin{proof}
If $B=\mathcal R_zA$, then $Av=0\Rightarrow Bv=\mathcal R_z0=0$, the core is included. On the contrary, for $w\in\operatorname{im}A$, select $v$ to make $Av=w$ and define $\mathcal R_z(w)=Bv$. If $Av=Av'$ , then $v-v'\in\ker A\subseteq\ker B$ , and thus $Bv=Bv'$ ; so this definition is independent of the chosen representative, and linear operators are given by linearity. $A$ is like surjective , so the equation $\mathcal R_zA=B$ uniquely determines $\mathcal R_z$ .

There is a complementary space $V$ in finite dimensional space, and $A|_V$ is injective and surjective to $\operatorname{im}A$. Taking its inverse as $S$, we have $AS=\operatorname{id}$. For each $w\in\operatorname{im}A$ , $\mathcal R_z(w)=B(Sw)$ , then $\mathcal R_z=BS$ , and the upper bound is obtained by the multiplicative inequality of the operator norm. When $\operatorname{im}A=\{0\}$, the kernel contains $B=0$, and the zero operator satisfies the same conclusion.\end{proof}


\begin{proposition}
[Exact factorization and minimal supplementary readout of affine observations]\label{exp:geo:prop:affine-factorization}
Suppose \[
\mathcal K(z)=Az+a_0\in\mathbb R^p,
\qquad
\Theta(z)=Bz+b_0\in\mathbb R^q,
\qquad z\in\mathbb R^n.
\] has a unique affine mapping $R:a_0+\operatorname{im}A\to\mathbb R^q$ such that $\Theta=R\circ\mathcal K$ iff $\ker A\subseteq\ker B$ . If $r$ linear readouts $Cz$ are allowed to be added, then the necessary and sufficient condition for $(\mathcal K(z),Cz)$ to accurately determine $\Theta(z)$ is \[
\ker A\cap\ker C\subseteq\ker B.
\] and the minimum number of reads for which this condition is true is \[
r_{\min}=\operatorname{rank}(B|_{\ker A})
=\dim\frac{\ker A}{\ker A\cap\ker B}.
\]\end{proposition}
\begin{certificate}
[Affine fibers, nuclear bases and supplementary matrices]Save the matrix $A,B$, offset $a_0,b_0$, kernel and image basis. If kernel inclusion is satisfied, save the linear factor $R_0:\operatorname{im}A\to\mathbb R^q$ of theorem ~ \ref{hgov:geo:thm:factor} and concatenate $R(a_0+w)=b_0+R_0w$. The certificate read out additionally preserves the rank of $B|_{\ker A}$, the isomorphism of its image space to $\mathbb R^{r_{\min}}$, and the complement space used to extend the corresponding linear mapping from $\ker A$ to $\mathbb R^n$.\end{certificate}
\begin{proof}
If $\Theta=R\mathcal K$ , for any $v\in\ker A$ and $z\in\mathbb R^n$ , there is $\mathcal K(z+v)=\mathcal K(z)$ , and thus $Bv=\Theta(z+v)-\Theta(z)=0$ . This gives kernel inclusion. On the contrary, according to the theorem ~ \ref{hgov:geo:thm:factor} there is a unique $R_0$ that satisfies $B=R_0A$. Let $R(a_0+w)=b_0+R_0w$ , then \[
R(\mathcal K(z))=R(a_0+Az)=b_0+R_0Az=b_0+Bz=\Theta(z).
\] $\mathcal K$ is surjective to the affine image $a_0+\operatorname{im}A$ , so $R$ is unique on this image. Replace $A$ with the stacked matrix $\binom AC$, which is verified to be $\ker A\cap\ker C$, to obtain the supplementary conditions.

Remember $N=\ker A$, $r_0=\operatorname{rank}(B|_N)$. If a certain $C:\mathbb R^n\to\mathbb R^r$ satisfies the supplementary condition, then $\ker(C|_N)\subseteq\ker(B|_N)$ , so there is $L$ such that $B|_N=L(C|_N)$ exists by linear factorization. Therefore \[
r_0=\operatorname{rank}(B|_N)
\le\operatorname{rank}(C|_N)\le r,
\] gives a lower bound on the minimum number of reads. On the other hand, select the isomorphism $Q:\operatorname{im}(B|_N)\to\mathbb R^{r_0}$ and define $C_N=Q(B|_N)$ on $N$. Take any complement space $\mathbb R^n=N\oplus V$ and concatenate $C(n+v)=C_Nn$. Therefore, for $\ker(C|_N)=\ker(B|_N)$, the required core inclusion is established, so $r=r_0$ can be achieved. The rank-zero degree formula finally gives $r_0=\dim\bigl(N/\ker(B|_N)\bigr)$, which is the listed quotient space dimension.\end{proof}


\begin{remark}
[Add a precise example of readout]Let $z=(u,v,w)$, the original geometric observation be $\mathcal K(z)=(u,v)$, and the strategy observation be $\Theta(z)=(u+w,v-w)$. The original observed core is $\operatorname{span}(e_3)$ and $Be_3=(1,-1)$, so the number of supplementary reads required is one. After selecting $Cz=w$, the exact response matrix is \[
R=\begin{pmatrix}1&0&1\\0&1&-1\end{pmatrix},
\qquad R(u,v,w)^{\mathsf T}=(u+w,v-w)^{\mathsf T}.
\]. If the Euclidean error vector of the three readings is $e$, the response error is $Re$. Since the eigenvalue of $RR^{\mathsf T}=\begin{pmatrix}2&-1\\-1&2\end{pmatrix}$ is $1,3$ , there is $\lVert Re\rVert\le\sqrt3\lVert e\rVert$ . This example also gives the minimum dimension of the supplementary observations, the complete response matrix, and the error amplification bound.\end{remark}


\begin{remark}
[Respond to continuations with executable fixes]Theorem ~ \ref{hgov:geo:thm:factor} gives a well-defined response on the first variation of a fixed point. For the allowed geometric variation $\dot k\in\operatorname{im}A$ , the policy direction is $u=\mathcal R_z\dot k$ . The subsequent execution interface saves the feasible region of $u$, the actual step size, the task observation hold, and the acceptance result of the target or residual; when these data are jointly established, the geometric observation, heterogeneous transportation, obstacle representation and executable repair are exchanged by the same group of certificate continuations.\end{remark}


% The system sources of this chapter can be found in the Pure Mathematics Volume (S01) and the Unified Overview (GFrameworkSynthesis).

% The external basic literature follows ExtKobayashiNomizu, ExtBottTu, and ExtHatcher in the main file.


\chapter{Hom Complexes and Compatible Fillings for Heterogeneous Continuation}
\label{exp:ch:hom}


\label{hgov:homology}

The local description of the heterogeneous system is first typed and mapped into a common comparable substrate, and then the continuation defect generates repair obligations. The structural expansion in law-space thus obtains three connected mathematical objects: the image of the original operation, the core where the defect is located, and the quotient effect of the new operation on the two. This section expands the barrier, homotopy padding, and finite certificate interfaces of the document ~ \cite{S01,S37,S38} into a structure that can be verified step by step. The cochain complex, chain homotopy and kernel-image quotient adopt the hierarchical language of the document ~ \cite{ExtWeibel,ExtStacks}.

\section { Typed Hom complex with cross-substrate continuation defect }\label{hgov:typed-hom}


\begin{definition}
[substrate, graded mapping and Hom differentiation]\label{hgov:def-hom}
Taking the domain $\mathbb k$, let the $i$th local substrate be the cochain complex \[
 C_i=(C_i^\bullet,d_i),\qquad
 d_i^m:C_i^m\longrightarrow C_i^{m+1},\qquad
 d_i^{m+1}d_i^m=0.
\]. The mapping space of $r$ and its differential are defined as \[
 \operatorname{Hom}^r(C_i,C_j)
 =\prod_m\operatorname{Hom}_{\mathbb k}(C_i^m,C_j^{m+r}),
 \qquad
 \delta_{ij}f=d_jf-(-1)^rfd_i.
\]. Therefore, the degree of $\delta_{ij}f$ is $r+1$. For $f:C_i\to C_j$ and $g:C_j\to C_k$, the composite $gf$ is only defined when the public substrate $C_j$ and its degree are consistent. A continuation mapping $S_{ij}$ is a mapping of degree zero that satisfies \[
 d_jS_{ij}=S_{ij}d_i
\]. Each map simultaneously registers local objects, versions, bases, and valid domains.\end{definition}


\begin{lemma}
[Square zeros and signed composite laws of Hom differentials]\label{hgov:lem-hom}
If $|f|=r$, $|g|=s$, then \[
 \delta_{ij}^{\,2}f=0,\qquad
 \delta_{ik}(gf)=(\delta_{jk}g)f+(-1)^s g(\delta_{ij}f).
\]\end{lemma}
\begin{certificate}
[Fractionation, source-target and item-by-item elimination]Save $d_i^2=d_j^2=d_k^2=0$, $r,s$ and the source and target of each composite item. Verify the sum of the coefficients of the two intermediate terms $(-1)^r+(-1)^{r+1}=0$ item by item, and verify that $-(-1)^s gd_jf$ and $(-1)^s gd_jf$ cancel in the compound law.\end{certificate}
\begin{proof}
By definition, \[
 \begin{aligned}
 \delta^2f
 &=d_j^2f-(-1)^rd_jfd_i
   -(-1)^{r+1}d_jfd_i+(-1)^{2r+1}fd_i^2\\
 &=0.
 \end{aligned}
\] and \[
 \begin{aligned}
 (\delta g)f+(-1)^sg(\delta f)
 &=d_kgf-(-1)^sgd_jf
   +(-1)^sgd_jf-(-1)^{r+s}gfd_i\\
 &=\delta(gf).
 \end{aligned}
\], so the predicate deduction chain is \[
 \bigl(d_i^2=d_j^2=d_k^2=0\bigr)\wedge
 \operatorname{Typed}(g,f)
 \ \Longrightarrow\
 \delta^2=0\ \wedge\ \operatorname{Leibniz}_{r,s}.
\]\end{proof}


\begin{definition}
[Triple continuation flaw]\label{hgov:def-defect}
Given the continuation mapping $S_{12},S_{23},S_{13}$, define \[
 J_{123}=S_{23}S_{12}-S_{13}
 \in\operatorname{Hom}^0(C_1,C_3).
\] and its two items respectively read the composite conversion and direct conversion of $C_2$. The filling witness of the number $-1$ is \[
 h_{123}\in\operatorname{Hom}^{-1}(C_1,C_3),\qquad
 d_3h_{123}+h_{123}d_1=J_{123}.
\]. This witness organizes the two transformations into a common chain homotopy.\end{definition}


\begin{theorem} [Closedness, fillability, and homological readback of continuation]
\label{hgov:thm-filler}
When the three continuation maps satisfy Definition ~ \ref{hgov:def-hom},
\[
 \delta_{13}J_{123}=0,
 \qquad
 \exists h_{123}\ \delta_{13}h_{123}=J_{123}
 \ \Longleftrightarrow\
 [J_{123}]=0\in H^0\operatorname{Hom}(C_1,C_3).
\].
When a filling witness exists, for every degree $m$,
\[
 H^m(S_{23})H^m(S_{12})=H^m(S_{13}).
\].
\end{theorem}
\begin{certificate} [Closedness and primitives]
Retain the three families of commutation relations $d_jS_{ij}=S_{ij}d_i$. In the successful branch, retain all graded components of $h_{123}$ and $d_3h_{123}+h_{123}d_1-J_{123}=0$. Given a closed element $z\in C_1^m$, also retain its preimage $h_{123}z\in C_3^{m-1}$.
\end{certificate}
\begin{proof}
Lemma ~ \ref{hgov:lem-hom} gives
\[
 \delta J_{123}
 =(\delta S_{23})S_{12}+S_{23}(\delta S_{12})
   -\delta S_{13}=0.
\].
Degree $-1$ implies
$\delta h=d_3h+h d_1$.
Consequently,
\[
 \begin{aligned}
 \exists h\ (\delta h=J_{123})
 &\Longleftrightarrow
 J_{123}\in\operatorname{im}
 \bigl(\delta:\operatorname{Hom}^{-1}\to\operatorname{Hom}^{0}\bigr)\\
 &\Longleftrightarrow [J_{123}]=0.
 \end{aligned}
\].
When $d_1z=0$,
\[
 S_{23}S_{12}z-S_{13}z
 =d_3(hz)+h(d_1z)=d_3(hz),
\],
so the two outputs differ by a coboundary. The complete derivation forms
\[
 \begin{gathered}
 \operatorname{Chain}(S_{12},S_{23},S_{13})
 \Longrightarrow\operatorname{Closed}(J_{123}),\\
 \operatorname{Closed}(J_{123})
 \Longrightarrow
 \bigl(\operatorname{Fill}(J_{123})\Longleftrightarrow[J_{123}]=0\bigr),\\
 \operatorname{Chain}(S_{12},S_{23},S_{13})
 \wedge\operatorname{Fill}(J_{123})
 \Longrightarrow\operatorname{EqualOnCohomology}.
 \end{gathered}
\].
\end{proof}

\begin{lemma}
[finite composite of continuation homotopy]\label{hgov:lem-compose-homotopy}
Assume that $f_0,f_1:C\to D$ and $g_0,g_1:D\to E$ are co-chain mappings, and \[
 f_1-f_0=d_DH+Hd_C,\qquad
 g_1-g_0=d_EK+Kd_D,
 \qquad |H|=|K|=-1.
\] then $L=g_1H+Kf_0$ satisfies \[
 g_1f_1-g_0f_0=d_EL+Ld_C.
\]. Any finite number of such certificates can be combined in continuation order.\end{lemma}
\begin{certificate}
[Composite witness and length induction]Save the source-target consistency, $L=g_1H+Kf_0$, and two sets of original homotopy formulas. Zero segments use identity mapping and zero homotopy; each additional segment uses $L$ listed.\end{certificate}
\begin{proof}
Using the same intermediate term $g_1f_0$ to decompose, the third formula of \[
 \begin{aligned}
 g_1f_1-g_0f_0
 &=g_1(f_1-f_0)+(g_1-g_0)f_0\\
 &=g_1(d_DH+Hd_C)+(d_EK+Kd_D)f_0\\
 &=d_E(g_1H+Kf_0)+(g_1H+Kf_0)d_C.
 \end{aligned}
\] uses $g_1d_D=d_Eg_1$ and $d_Df_0=f_0d_C$ respectively. If $P(n)$ indicates that the previous $n$ segment has synthesized homotopy, then $P(0)$ is established, and the listed structure gives $P(n)\Rightarrow P(n+1)$; mathematical induction gives $\forall n\in\mathbb N,\ P(n)$.\end{proof}


\section { Two-dimensional filling and its semantic implementation }\label{hgov:two-cell}


\begin{definition}
[Labeled triangular continuation rendering]\label{hgov:def-triangle}
Take the vertex $v_1,v_2,v_3$ and the directed edge $e_{12},e_{23},e_{13}$, and let the semantic tags of the \[
 \partial_1e_{ij}=v_j-v_i,\qquad
 z=e_{12}+e_{23}-e_{13}.
\] edges be $S_{12},S_{23},S_{13}$ respectively. The geometric boundary of the two-dimensional cell $f$ is specified as $\partial_2f=z$; its semantic implementation is specified as $h_f$ of degree $-1$, and satisfies \[
 \delta h_f=S_{23}S_{12}-S_{13}.
\]. The geometric field records the directional boundary, and the semantic field records morphism composition and homotopy. Witness that the two are bound by the same cellular identifier.\end{definition}


\begin{proposition}
[Filling of triangular holes and implementation of continuation]\label{hgov:prop-triangle}
When only three edges are retained, the obtained complex $K^-$ satisfies \[
 H_1(K^-;\mathbb k)=\mathbb k[z],\qquad
 H_0(K^-;\mathbb k)\cong\mathbb k.
\]. After appending $f$ and $\partial_2f=z$, the complex $K^+$ satisfies \[
 H_1(K^+;\mathbb k)=0,\qquad
 H_0(K^+;\mathbb k)\cong\mathbb k.
\]. If the cell has a definition The semantic implementation of ~ \ref{hgov:def-triangle} is the same in cohomology as direct continuation via intermediate substrate.\end{proposition}
\begin{certificate}
[Boundary matrices and realization equations]Sort by $(v_1,v_2,v_3)$ and $(e_{12},e_{23},e_{13})$, save \[
 [\partial_1]=
 \begin{pmatrix}-1&0&-1\\1&-1&0\\0&1&1\end{pmatrix},
 \qquad
 [\partial_2]=\begin{pmatrix}1\\1\\-1\end{pmatrix}.
\], verify $\partial_1\partial_2=0$, $\operatorname{rank}\partial_1=2$, $\operatorname{rank}\partial_2=1$, and save $\delta h_f=S_{23}S_{12}-S_{13}$ separately.\end{certificate}
\begin{proof}
The first two columns are linearly independent, and the third column is equal to the sum of the first two columns, so $\operatorname{rank}\partial_1=2$, the core is composed of $z$. In $K^-$, $C_2=0$, so the first-level boundary image is zero; in $K^+$, the image is equal to $\mathbb k[z]=\ker\partial_1$. $\partial_1$ is the same for both, and the remaining kernels are one-dimensional. Therefore the last formula \[
 \begin{gathered}
 \partial_1z=0
 \Longrightarrow z\in\ker\partial_1,\\
 \operatorname{im}\partial_2^-=\{0\}
 \Longrightarrow[z]_{K^-}\ne0,\\
 \partial_2^+f=z
 \Longrightarrow[z]_{K^+}=0,\\
 \delta h_f=J_{123}
 \Longrightarrow H(S_{23})H(S_{12})=H(S_{13}).
 \end{gathered}
\] is given by the theorem ~ \ref{hgov:thm-filler}.\end{proof}


\section { Filling the solution space, second-order homotopy and its quotient }\label{exp:hom:fill-space}


Filling of the same continuation defect often has multiple implementations. To itemize this choice, set \[
 \begin{aligned}
 Z^r_{ij}
 &=\ker\bigl(\delta_{ij}:
       \operatorname{Hom}^r(C_i,C_j)\to
       \operatorname{Hom}^{r+1}(C_i,C_j)\bigr),\\
 B^r_{ij}
 &=\operatorname{im}\bigl(\delta_{ij}:
       \operatorname{Hom}^{r-1}(C_i,C_j)\to
       \operatorname{Hom}^{r}(C_i,C_j)\bigr),\\
 H^r_{ij}&=Z^r_{ij}/B^r_{ij}.
 \end{aligned}
\] Lemma ~ \ref{hgov:lem-hom} gives $B^r_{ij}\subseteq Z^r_{ij}$ . The following results separately register the existence of solutions, the differences between solutions, and the homotopy between solutions.

\begin{theorem}
[Filled free transfer effects and cohomology parameters]\label{exp:hom:filler-torsor}
Fixed $J\in Z^0_{13}$ and obtained a witness for $h_0\in\operatorname{Hom}^{-1}(C_1,C_3)$, $\delta h_0=J$. Let \[
 \mathcal F(J)
 =\{h\in\operatorname{Hom}^{-1}(C_1,C_3):\delta h=J\}.
\] then the additive effect of \[
 \mathcal F(J)=h_0+Z^{-1}_{13}.
\] $(z,h)\mapsto h+z$ is the free transfer effect of $Z^{-1}_{13}$ on $\mathcal F(J)$. Define \[
 h\sim h'
 \quad\Longleftrightarrow\quad
 \exists b\in\operatorname{Hom}^{-2}(C_1,C_3),\
 h'-h=\delta b.
\] and the resulting quotient set $\mathcal F(J)/{\sim}$ has the free transitive effect of $H^{-1}_{13}$. After selecting $h_0$, there is bijective \[
 \mathcal F(J)/{\sim}\longrightarrow H^{-1}_{13},
 \qquad [h]_{\sim}\longmapsto[h-h_0].
\]\end{theorem}
\begin{certificate}
[One primitive and two layers of kernel images]Save the basis of $\delta h_0=J$, $Z^{-1}_{13}$, and the basis of $B^{-1}_{13}$. For any candidate $h$, save $z=h-h_0$ and $\delta z=0$; for the equivalent candidate $(h,h')$, save $b$ the number of times $-2$, and verify $h'-h=d_3b-bd_1$ one by one. . When these spaces have finite dimensions and the base is on $\mathbb Q$, the base and quotient coordinates can be obtained by finite elimination.\end{certificate}
\begin{proof}
If $\delta h=J$ , linearity gives $\delta(h-h_0)=J-J=0$ ; if $z\in Z^{-1}_{13}$ , then $\delta(h_0+z)=J$ . This proves the solution set formula. The freedom of action is obtained from $h+z=h\Rightarrow z=0$; for any $h,h'$, the only closed element that sends $h$ to $h'$ is $h'-h$, so the action is transferred. Since $B^{-1}_{13}$ is a subspace, the relation $\sim$ is reflexive, symmetric, and transitive. Both sets of candidates give the same quotient coordinates which are exactly equivalent to \[
 [h-h_0]=[h'-h_0]
 \Longleftrightarrow h'-h\in B^{-1}_{13}
 \Longleftrightarrow h\sim h'.
\]. Each closed element $z$ is realized by candidate $h_0+z$, so the coordinates map surjective. Action and coordinate mapping thus have the properties described. When the division is $-2$, it is $(-1)^{-2}=1$, so the differential of the second-level witness is $\delta b=d_3b-bd_1$. The deduction chain is \[
 \begin{gathered}
 \delta h_0=J\ \land\ \delta h=J
 \Rightarrow\delta(h-h_0)=0,\\
 \delta b=h'-h
 \Rightarrow[h'-h_0]=[h-h_0],\\
 \operatorname{Chosen}(h_0)\ \land\
 \operatorname{KernelImage}(Z^{-1}_{13},B^{-1}_{13})
 \Rightarrow
 \mathcal F(J)/{\sim}\cong H^{-1}_{13}.
 \end{gathered}
\]\end{proof}


\section { Secondary continuation defect of four interfaces }\label{exp:hom:four-interfaces}


Given $C_1,C_2,C_3,C_4$ and all six zero-degree continuation mappings $S_{ij}:C_i\to C_j$, $1\le i<j\le4$. For each triple $i<j<k$, specify an implemented padding \[
 h_{ijk}:C_i^\bullet\to C_k^{\bullet-1},
 \qquad
 \delta h_{ijk}=S_{jk}S_{ij}-S_{ik}.
\] where the target of $h_{123}$ is $C_3$ and the source of $h_{234}$ is $C_2$; composite Both $S_{34}h_{123}$ and $h_{234}S_{12}$ have source $C_1$ and target $C_4$. In this way, homotopy comparisons of two triple transformations can be performed in the same Hom complex.

\begin{theorem}
[Secondary obstacles and times of four interfaces $-2$ filled]\label{exp:hom:secondary-obstruction}
Under the above conditions, define \[
 \begin{aligned}
 H_{\mathrm L}&=S_{34}h_{123}+h_{134},\\
 H_{\mathrm R}&=h_{234}S_{12}+h_{124},\\
 \Theta_{1234}
 &=H_{\mathrm L}-H_{\mathrm R}
 \in\operatorname{Hom}^{-1}(C_1,C_4).
 \end{aligned}
\], then \[
 \begin{aligned}
 \delta H_{\mathrm L}
 &=S_{34}S_{23}S_{12}-S_{14}
 =\delta H_{\mathrm R},\\
 \delta\Theta_{1234}&=0.
 \end{aligned}
\] The necessary and sufficient condition for the existence of the second-level compatible witness $k_{1234}$ is \[
 \begin{gathered}
 \exists k_{1234}\in\operatorname{Hom}^{-2}(C_1,C_4),\
 d_4k_{1234}-k_{1234}d_1=\Theta_{1234}\\
 \Longleftrightarrow\
 [\Theta_{1234}]=0\in H^{-1}_{14}.
 \end{gathered}
\]. Given this witness and the closed element $z\in C_1^m$, two homotopies are in $z$ The difference on is the boundary $\Theta_{1234}z=d_4(k_{1234}z)$.\end{theorem}
\begin{certificate}
[Four faces and one secondary primitive]Save six groups of $\delta S_{ij}=0$, four groups of $\delta h_{ijk}=S_{jk}S_{ij}-S_{ik}$, and two homotopy typed expansions. The secondary acceptance branches hold $|k_{1234}|=-2$ and $d_4k_{1234}-k_{1234}d_1-\Theta_{1234}=0$. When $z$ is input, $d_1z=0$ and the boundary original image $k_{1234}z\in C_4^{m-2}$ are recorded.\end{certificate}
\begin{proof}
The signed compound law gives \[
 \begin{aligned}
 \delta(S_{34}h_{123})
 &=S_{34}\delta h_{123},\\
 \delta(h_{234}S_{12})
 &=(\delta h_{234})S_{12}
   -h_{234}\delta S_{12}
 =(\delta h_{234})S_{12}.
 \end{aligned}
\]. Substituting four triangular filling formulas in sequence, \[
 \begin{aligned}
 \delta H_{\mathrm L}
 &=S_{34}(S_{23}S_{12}-S_{13})
   +(S_{34}S_{13}-S_{14})\\
 &=S_{34}S_{23}S_{12}-S_{14},\\
 \delta H_{\mathrm R}
 &=(S_{34}S_{23}-S_{24})S_{12}
   +(S_{24}S_{12}-S_{14})\\
 &=S_{34}S_{23}S_{12}-S_{14}.
 \end{aligned}
\] is subtracted to obtain the closed property. Defined by the kernel-image quotient of cohomology, the closed degree $-1$ element $\Theta_{1234}$ is the boundary exactly equivalent to the existence degree $-2$ primitive element. Calculate \[
 \Theta_{1234}z
 =(d_4k_{1234}-k_{1234}d_1)z
 =d_4(k_{1234}z).
\] on the closed input. This also verifies that the output times are in sequence $m-2,m-1$. Its predicate chain is \[
 \begin{gathered}
 \operatorname{SixChainMaps}\land\operatorname{FourFillers}
 \Rightarrow\delta H_{\mathrm L}=\delta H_{\mathrm R}
 \Rightarrow\Theta_{1234}\in Z^{-1}_{14},\\
 \Theta_{1234}\in Z^{-1}_{14}
 \Rightarrow
 \bigl(\exists k:\delta k=\Theta_{1234}
       \Longleftrightarrow[\Theta_{1234}]=0\bigr),\\
 \delta k=\Theta_{1234}\land d_1z=0
 \Rightarrow\Theta_{1234}z=d_4(kz).
 \end{gathered}
\]\end{proof}


\begin{remark}
[A secondary barrier that can be verified on a base-by-base basis]\label{exp:hom:secondary-example}
Take four substrates, all of which are $C^0=\mathbb Q e$ and $C^{-1}=\mathbb Q f$, and the remaining times are zero and $d=0$. The six $S_{ij}$ are all identity mappings. Let $h_{123}(e)=f$ , all remaining components and the remaining three $h_{ijk}$ be zero. Each triangle defect is zero, the four surface filling formulas are true, and $\Theta_{1234}(e)=f\ne0$. Since $\delta=0$ , this element gives a non-zero $H^{-1}_{14}$ class.

The extended substrate is \[
 C'^{-2}=\mathbb Q g,\qquad
 C'^{-1}=\mathbb Q f,\qquad C'^0=\mathbb Q e,
 \qquad d'g=f,\quad d'f=d'e=0.
\], keeping the six identity maps and the components of the above $h_{ijk}$, adding $k_{1234}(e)=g$, and the remaining components are zero. The base-by-base calculation yields \[
 \begin{array}{c|ccc}
 x&e&f&g\\ \hline
 (d'k_{1234}-k_{1234}d')(x)&f&0&0\\
 \Theta_{1234}(x)&f&0&0
 \end{array}
\] to obtain secondary filling. Containing $C\hookrightarrow C'$ is a co-chain mapping and induces the identity $H^0(C)=\mathbb Q[e]\to H^0(C')=\mathbb Q[e]$ on degree zero; here the task read is specified to be of degree zero class. The extension also makes the original degree $-1$ of class $[f]$ the boundary $d'g$. The number of times the task is retained and the number of times it is resolved are registered by two groups of independent witnesses.\end{remark}


\begin{proposition}
[Surface-filling secondary changes and obstacle class maintenance]\label{exp:hom:secondary-invariance}
Fix six $S_{ij}$ and give the degree $-2$ of $b_{ijk}$ to the four faces. If \[
 h'_{ijk}=h_{ijk}+\delta b_{ijk},\qquad
 \Xi=S_{34}b_{123}+b_{134}-b_{234}S_{12}-b_{124}.
\] is set, the new surface filling still satisfies the original triangle formula, and if \[
 \Theta'_{1234}-\Theta_{1234}=\delta\Xi,
 \qquad [\Theta'_{1234}]=[\Theta_{1234}]\in H^{-1}_{14}.
\] already exists $\delta k_{1234}=\Theta_{1234}$, then $k'_{1234}=k_{1234}+\Xi$ is the new secondary filling.\end{proposition}
\begin{certificate}
[Transportation of surface and body modifications]Save the source target and times of four $b_{ijk}$, four groups of $h'_{ijk}-h_{ijk}=\delta b_{ijk}$, and $\Xi\in\operatorname{Hom}^{-2}(C_1,C_4)$. Verification $\Theta'-\Theta-\delta\Xi=0$; additional verification of successful branch $\delta(k+\Xi)=\Theta'$.\end{certificate}
\begin{proof}
By $\delta^2=0$, $\delta h'_{ijk}=\delta h_{ijk}$. Since $S_{ij}$ is closed and has degree zero, $b_{ijk}$ has degree $-2$, \[
 \begin{aligned}
 \delta(S_{34}b_{123})&=S_{34}\delta b_{123},\\
 \delta(b_{234}S_{12})&=(\delta b_{234})S_{12}
                       +b_{234}\delta S_{12}
                      =(\delta b_{234})S_{12}.
 \end{aligned}
\]. Substitute four $h'_{ijk}-h_{ijk}$ into $\Theta'-\Theta$ to get $\delta\Xi$. The class of coboundary is zero, so the barrier class remains; linearity gives a new filler formula. The formula-by-form deduction is \[
 \begin{gathered}
 h'-h=\delta b\land\delta^2=0
 \Rightarrow\operatorname{SameFaceDefects},\\
 \operatorname{Chain}(S)\land|b|=-2
 \Rightarrow\Theta'-\Theta=\delta\Xi
 \Rightarrow[\Theta']=[\Theta],\\
 \delta k=\Theta\land\Theta'=\Theta+\delta\Xi
 \Rightarrow\delta(k+\Xi)=\Theta'.
 \end{gathered}
\]\end{proof}


\section { Allows changing the secondary obstacle elimination in the subspace }\label{exp:hom:allowed-face-variations}


When a face fill has been saved, subsequent changes to it can be restricted by version, resources, or committed interfaces. In the following, modifications are allowed by explicit linear subspace registration, and a restricted existential quantifier is obtained. Assume that \[
 \begin{aligned}
 W_{123}&\subseteq Z^{-1}_{13},&
 W_{134}&\subseteq Z^{-1}_{14},\\
 W_{234}&\subseteq Z^{-1}_{24},&
 W_{124}&\subseteq Z^{-1}_{14},\\
 W_{\mathrm{face}}
 &=W_{123}\oplus W_{134}\oplus W_{234}\oplus W_{124}.
 \end{aligned}
\] stipulates that the allowed surface filling is exactly $h_{ijk}+\eta_{ijk}$, $\eta_{ijk}\in W_{ijk}$. The fixed face corresponds to $W_{ijk}=0$. The target defining the linear mapping \[
 \begin{aligned}
 F:W_{\mathrm{face}}&\longrightarrow Z^{-1}_{14},\\
 F(\eta)&=S_{34}\eta_{123}+\eta_{134}
          -\eta_{234}S_{12}-\eta_{124},\\
 \Psi:W_{\mathrm{face}}&\longrightarrow H^{-1}_{14},
 \qquad \Psi(\eta)=[F(\eta)].
 \end{aligned}
\] $F$ is guaranteed by closure and Hom's composite law.

\begin{theorem}
[Precise repairability conditions for restricted surface modifications]\label{exp:hom:restricted-secondary}
Under the above fixed data, the following conditions are equivalent: \[
 \begin{gathered}
 \exists\eta\in W_{\mathrm{face}},\
 \exists k\in\operatorname{Hom}^{-2}(C_1,C_4),\
 \delta k=\Theta_{1234}+F(\eta),\\
 -[\Theta_{1234}]\in\operatorname{im}\Psi.
 \end{gathered}
\] If each space has a finite rational basis, this condition can be solved by a linear finite linear equation: \[
 \begin{pmatrix}\delta_{-2}&-F\end{pmatrix}
 \begin{pmatrix}k\\ \eta\end{pmatrix}
 =\Theta_{1234}.
\] The failed branch can submit a rational row vector $\lambda^{\mathsf T}$ satisfies \[
 \lambda^{\mathsf T}\delta_{-2}=0,\qquad
 \lambda^{\mathsf T}F=0,\qquad
 \lambda^{\mathsf T}\Theta_{1234}\ne0.
\] where the row coordinates of all matrices are taken from the same $\operatorname{Hom}^{-1}(C_1,C_4)$ base.\end{theorem}
\begin{certificate}
[Allow domains, combined preimages or joint vetoes]Save the basis and closure of each $W_{ijk}$, the matrices of $F$, $\Theta_{1234}$ and $\delta_{-2}$. Accepting the branch commit allows the coordinates $\eta$ and $k$; rejecting the branch commit $\lambda$ is the product of two sets of zeros and a non-zero reading. All face transformations are played back as $h'_{ijk}=h_{ijk}+\eta_{ijk}$.\end{certificate}
\begin{proof}
For the allowed $\eta$, $\delta\eta_{ijk}=0$ maintains four surface defects, and the new secondary defect is exactly $\Theta_{1234}+F(\eta)$. Theorem ~ \ref{exp:hom:secondary-obstruction} Therefore, given \[
 \begin{aligned}
 \exists k:\delta k=\Theta_{1234}+F(\eta)
 &\Longleftrightarrow
 [\Theta_{1234}]+\Psi(\eta)=0.
 \end{aligned}
\], and then taking the existential quantifier for the allowed $\eta$, we can get the image criterion. In finite basis, $\delta k-F\eta=\Theta_{1234}$ is the coordinate representation of the same equation. Perform rational row elimination on its augmented matrix with finite number of pivots; the back substitution gives $(k,\eta)$ and the contradiction row gives the listed $\lambda$ . If the two branches are established at the same time, a conflict occurs in \[
 0=\lambda^{\mathsf T}\delta k-\lambda^{\mathsf T}F\eta
   =\lambda^{\mathsf T}\Theta_{1234}\ne0,
\]. All quantifiers here are restricted to the allowed space of registration. Its predicate chain is \[
 \begin{gathered}
 \eta\in W_{\mathrm{face}}
 \Rightarrow\operatorname{SameFaceDefects}
 \land\Theta'=\Theta+F\eta,\\
 \operatorname{SameFaceDefects}
 \Rightarrow
 \bigl(\exists\eta\in W_{\mathrm{face}},k:\delta k=\Theta+F\eta
 \Longleftrightarrow-[\Theta]\in\operatorname{im}\Psi\bigr),\\
 \operatorname{FiniteRationalMatrices}
 \Rightarrow
 \operatorname{JointWitness}(k,\eta)
 \ \vee\ \operatorname{JointObstruction}(\lambda).
 \end{gathered}
\]\end{proof}


\chapter{Finite Repair Complexes, Admissible Directions, and Minimal Compatible Extensions}
\label{exp:ch:extension}


\section { Orthogonal residuals, dual vetoes and structural solutions of finite matrices }\label{hgov:linear-certificate}


\begin{definition}
[finite linear repair rendering]\label{hgov:def-linear}
Taking the finite dimensional space and mapping \[
 U\xrightarrow{L}V\xrightarrow{D}W,\qquad DL=0,
 \qquad \omega\in V,\quad D\omega=0.
\], the repair equation is $Lu=\omega$. For Hom filling, $U=\operatorname{Hom}^{-1}$, $V=\operatorname{Hom}^{0}$, $W=\operatorname{Hom}^{1}$, and $L,D$ can be taken as the corresponding Hom differentials respectively. In the real implementation, $G$ is the positive definite Gram matrix on $V$, $\|v\|_G^2=v^{\mathsf T}Gv$.\end{definition}


\begin{theorem}
[Equivalent certificate of orthogonal residuals and homology solvability]\label{hgov:thm-projection}
In the real implementation defining ~ \ref{hgov:def-linear}, let $P_G$ be the orthogonal projection of $G$ to $\operatorname{im}L$, $r=(I-P_G)\omega$. Then \[
 \min_{u\in U}\|Lu-\omega\|_G=\|r\|_G,
 \qquad
 \|r\|_G=0
 \Longleftrightarrow
 [\omega]=0\in\ker D/\operatorname{im}L.
\] when $r\ne0$, $\lambda=Gr$ gives \[
 \lambda^{\mathsf T}L=0,\qquad
 \lambda^{\mathsf T}\omega=\|r\|_G^2>0.
\]. When the equation is solvable and $U$ is equipped with a positive definite inner product, all solutions are $u_0+\ker L$, where $u_0\in(\ker L)^\perp$ is the only minimum norm solution.\end{theorem}
\begin{certificate}
[Image kernel decomposition and dual readout]Save $L,D,G,\omega$, $DL=0$, $D\omega=0$, like space basis, $\omega=P_G\omega+r$, $L^{\mathsf T}Gr=0$ and some $u_\parallel$ that satisfies $Lu_\parallel=P_G\omega$. The solvable branches save $Lu_0=\omega$ and $u_0\perp\ker L$; the other branches save $\lambda=Gr$ and the strictly positive value $\lambda^{\mathsf T}\omega$.\end{certificate}
\begin{proof}
Finite-dimensional subspaces have orthogonal decompositions. For any $u$, \[
 \begin{aligned}
 Lu-\omega&=(Lu-P_G\omega)-r,\\
 \|Lu-\omega\|_G^2
 &=\|Lu-P_G\omega\|_G^2+\|r\|_G^2.
 \end{aligned}
\] takes $u=u_\parallel$ to reach the lower bound. The closed property $D\omega=0$ gives \[
 \|r\|_G=0
 \Longleftrightarrow \omega\in\operatorname{im}L
 \Longleftrightarrow[\omega]=0.
\] and because $r\perp_G\operatorname{im}L$ , \[
 \lambda^{\mathsf T}L=r^{\mathsf T}GL=0,\qquad
 \lambda^{\mathsf T}\omega
 =r^{\mathsf T}G(P_G\omega+r)=r^{\mathsf T}Gr.
\] if $Lu=\omega$ , then the difference between the two solutions belongs to $\ker L$ . Decompose a solution along $(\ker L)^\perp\oplus\ker L$ to get the unique $u_0$ and $\|u_0+k\|^2=\|u_0\|^2+\|k\|^2$. The predicate causal chain is \[
 \operatorname{Typed}(L,D)\wedge DL=0\wedge D\omega=0
 \Longrightarrow
 \begin{cases}
 r=0\ \Longrightarrow\ Lu_0=\omega,\\
 r\ne0\ \Longrightarrow\
 \lambda^{\mathsf T}L=0,\ \lambda^{\mathsf T}\omega>0.
 \end{cases}
\]. The second branch obtains the unsolvable certificate by applying $\lambda^{\mathsf T}$ to any candidate solution.\end{proof}


\begin{proposition}
[Rationally presented finite elimination certificate]\label{hgov:prop-elimination}
When the coordinates of $L,\omega$ belong to $\mathbb Q$, you can construct the rational vector $Lu=\omega$ in finite steps, or construct the rational row vector $\lambda^{\mathsf T}$ such that \[
 \lambda^{\mathsf T}L=0,\qquad
 \lambda^{\mathsf T}\omega\ne0.
\]\end{proposition}
\begin{certificate}
[Row Actions and Denial Rows]Perform line wrapping, non-zero rational scaling, and row addition on the augmented matrix $[L\mid\omega]$; simultaneously record the invertible matrix $E$. The successful branch records the free variable selection and back-substitution solution. The row of $E$ corresponding to the contradictory rows $[0\mid c]$ and $c\ne0$ is $\lambda^{\mathsf T}$.\end{certificate}
\begin{proof}
Each time a pivot is created, the number of processed pivot columns increases by one; the rows and columns are finite, so the elimination of finite elements terminates. Reversibility gives \[
 Lu=\omega\Longleftrightarrow ELu=E\omega.
\] when there are contradictory rows, $\lambda^{\mathsf T}Lu=0\ne\lambda^{\mathsf T}\omega$ . In other cases, the free variable is set to zero and back-substituted, and all pivot rows and zero rows are established. Therefore \[
 \operatorname{RationalFinite}(L,\omega)
 \Longrightarrow
 \operatorname{FiniteElimination}
 \Longrightarrow
 \operatorname{Witness}(u)\ \vee\ \operatorname{Witness}(\lambda).
\]\end{proof}


\section { Compatible expansion and at least new independent direction }\label{hgov:compatible-extension}


\begin{definition}
[Structural expansion under conservation and identity constraints]\label{hgov:def-extension}
Taking the finite dimension chain complex \[
 C_2\xrightarrow{A}C_1\xrightarrow{B}C_0,\qquad BA=0,
\] and the linear identity specified to be preserved read $T:C_1\to Q$ , satisfying $TA=0$ . Note $U=\operatorname{im}A$. The target repair subspace satisfies \[
 Z\subseteq\ker B\cap\ker T.
\]. The compatible extension gives a new operation space $E$ and $R:E\to C_1$. It satisfies \[
 BR=0,\qquad TR=0,\qquad
 A^+:C_2\oplus E\longrightarrow C_1,\quad A^+(u,v)=Au+Rv.
\]. The mapping from the original complex to the extended complex $\iota$ is in degree two. $u\mapsto(u,0)$, the degrees one and zero are identical. The actual implementation of the new operation is separately bound by the actual action, precondition and resource fields.\end{definition}


\begin{theorem}
[Expanded homology quotient and growth in rigorous repair capabilities]\label{hgov:thm-extension}
Define the data of ~ \ref{hgov:def-extension} to satisfy \[
 BA^+=0,\qquad TA^+=0.
\] Let $\pi:\ker B\to H_1(C)=\ker B/U$ be a quotient map, then \[
 H_1(C^+)\cong
 H_1(C)\big/\pi(\operatorname{im}R),
 \qquad
 \ker\iota_*=\pi(\operatorname{im}R).
\] for any $\omega\in\ker B$, \[
 [\omega]_{C^+}=0
 \Longleftrightarrow
 \exists u\in C_2,\ v\in E:\ \omega=Au+Rv.
\] if otherwise $[\omega]_C\ne0$, then $\omega$ is a strict witness of the new repair capability.\end{theorem}
\begin{certificate}
[Contains, image and new original image]Save $A,B,T,R$, two sets of zero products and include $\iota$; save the basis of $U$ and $U+\operatorname{im}R$. The strict growth branch simultaneously submits the dual veto vector of the original system and the new solution $(u,v)$.\end{certificate}
\begin{proof}
The matrix is divided into blocks to give $BA^+=[BA\ BR]=0$, $TA^+=[TA\ TR]=0$. The map $\iota$ is swapped with the boundary, and the induced $\iota_*$ is \[
 \ker B/U\longrightarrow\ker B/(U+\operatorname{im}R),
 \quad [x]_U\longmapsto[x]_{U+\operatorname{im}R}.
\] This map is surjective and its core is $(U+\operatorname{im}R)/U=\pi(\operatorname{im}R)$ . By the definition of quotient, \[
 [\omega]_{C^+}=0
 \Longleftrightarrow
 \omega\in U+\operatorname{im}R
 \Longleftrightarrow
 \exists(u,v)\ A^+(u,v)=\omega.
\] then \[
 \operatorname{Compatible}(R)
 \wedge[\omega]_C\ne0
 \wedge A^+(u,v)=\omega
 \Longrightarrow
 \operatorname{Repair}_{C^+}(\omega)
 \wedge\neg\operatorname{Repair}_{C}(\omega).
\] each conclusion is verified by the zero product, the old veto witness and the new preimage respectively.\end{proof}


\begin{theorem}
[The minimum number of new independent directions required for the target subspace]\label{hgov:thm-minimal-extension}
Under the definition ~ \ref{hgov:def-extension}, set \[
 r_Z=\dim\pi(Z)
 =\dim\frac{U+Z}{U}.
\]. If the extension can repair all $Z$, then \[
 r_Z\le\operatorname{rank}(\pi R)\le\dim E.
\] allows any base representative to be taken from $Z$ as a new operation. $E=\mathbb k^{r_Z}$ and compatible $R$ reach equal sign. This minimum value measures the number of independent directions added.\end{theorem}
\begin{certificate}
[Lower bound, quotient base and reach witness]The lower bound certificate saves $Z\subseteq U+\operatorname{im}R$. Reach certificate to save $z_1,\ldots,z_{r_Z}\in Z$, make $\pi(z_1),\ldots,\pi(z_{r_Z})$ form the base of $\pi(Z)$, and define $Re_j=z_j$. Each candidate $z_j$ saves the allowed operation domain to which it belongs.\end{certificate}
\begin{proof}
All can be repaired and equivalent to $Z\subseteq U+\operatorname{im}R$. Applying $\pi$ results in $\pi(Z)\subseteq\operatorname{im}(\pi R)$, so the lower bound of dimensionality is obtained. Take the quotient basis listed in reverse. For any $z\in Z$ there is a unique coefficient $a_1,\ldots,a_{r_Z}$ such that \[
 \pi(z)=\sum_{j=1}^{r_Z}a_j\pi(z_j).
\] such that $z-\sum_j a_jz_j\in U$ exists, $u$ satisfies $z=Au+Ra$. Also composed of $z_j\in Z$ and $Bz_j=Tz_j=0$, so $R$ is compatible. Taking $E=0$ when $r_Z=0$, the same derivation gives $Z\subseteq U$. The predicate causal chain is \[
 \begin{gathered}
 \forall z\in Z\ \operatorname{Repair}_{C^+}(z)
 \Longrightarrow
 \pi(Z)\subseteq\operatorname{im}(\pi R)
 \Longrightarrow \dim E\ge r_Z,\\
 \operatorname{Basis}\{\pi(z_j)\}_{j=1}^{r_Z}
 \wedge z_j\in Z
 \Longrightarrow BR=TR=0
 \Longrightarrow
 \forall z\in Z\ \exists(u,a)\ z=Au+Ra.
 \end{gathered}
\]. The lower bound and reaching construction jointly prove minimality.\end{proof}


\begin{corollary}
[Specifies the minimum cost expansion within the menu]\label{hgov:cor-finite-cost}
Given a finite candidate extension menu $\mathcal M$ , each member registers a compatibility check, a repair check for all $Z$ , and determines the real cost of $c(M)$ . If the set of qualified members $\mathcal M_Z$ is non-empty, then there exists $M_*\in\mathcal M_Z$ that satisfies \[
 c(M_*)=\min_{M\in\mathcal M_Z}c(M).
\]. If the cost comparison is decidable, item-by-item comparison gives a finite selection algorithm.\end{corollary}
\begin{certificate}
[Menu overlay and item-by-item comparison]Save all menu members, qualified predicates, the same cost specification, actual cost values and comparison order; the candidate enumeration range is used as the certificate valid domain.\end{certificate}
\begin{proof}
The set of finite non-empty real numbers has a minimum element. The algorithm starts with the first qualified member and keeps the lower cost member of the two for each new member. Let $P(n)$ be the cost of the current member not greater than the cost of the previous $n$ qualified members, then $P(1)$ is established and the comparison operation gives $P(n)\Rightarrow P(n+1)$. The $M_*$ is obtained when the traversal terminates.\end{proof}


\section { Accurate counting of relative repair complexes and new relationships }\label{exp:hom:relative-extension}


In this section, the three-term complex definition ~ \ref{hgov:def-extension} is denoted as $C_n=0$ ($n\notin\{0,1,2\}$). So $H_2(C)=\ker A$ , $H_0(C)=C_0/\operatorname{im}B$ . The quotient chain complex $C^+/C$ only has $E$ in degree two, and its boundary is zero: \[
 (C^+/C)_2=E,\qquad
 (C^+/C)_1=(C^+/C)_0=0.
\] The following directly constructs the connection map of this relatively complex, and proves the exactness point by point; for the usual expression of the chain complex short exact sequence, please refer to the literature ~ \cite{ExtWeibel,ExtStacks}.

\begin{theorem}
[Relative exact column and rank count of new operations]\label{exp:hom:relative-exact}
For the compatible extension $R:E\to\ker B\cap\ker T$, if \[
 \begin{aligned}
 j:\ker A&\longrightarrow\ker A^+,&j(u)&=(u,0),\\
 p:\ker A^+&\longrightarrow E,&p(u,e)&=e,\\
 \beta:E&\longrightarrow H_1(C),&\beta(e)&=[Re].
 \end{aligned}
\] is defined, then there is an exact sequence \[
 0\longrightarrow H_2(C)
 \xrightarrow{j}H_2(C^+)
 \xrightarrow{p}E
 \xrightarrow{\beta}H_1(C)
 \xrightarrow{\iota_*}H_1(C^+)
 \longrightarrow0.
\]. If $r_R=\operatorname{rank}\beta$ is set, then \[
 \begin{aligned}
 \dim H_1(C^+)&=\dim H_1(C)-r_R,\\
 \dim H_2(C^+)&=\dim H_2(C)+\dim E-r_R,\\
 H_0(C^+)&=H_0(C).
 \end{aligned}
\]. If $\dim E=r_R$ is set, then $j$ is isomorphic. The identity read also induces the well-defined $\bar T:H_1(C)\to Q$ , $\bar T^+:H_1(C^+)\to Q$ , and satisfies $\bar T^+\iota_*=\bar T$ .\end{theorem}
\begin{certificate}
[Connect the map to the original image of each core]Save $BA=BR=TA=TR=0$ and the bases of the three chain groups. Saving $u_e$ for $e\in\ker\beta$ makes $Re=Au_e$ , so $(-u_e,e)\in\ker A^+$ and $p(-u_e,e)=e$ . For $\xi\in\ker\iota_*$, save the representative element $x=Au+Re$, whose class is $\beta(e)$. Also save the basis and rank of $\ker A$, $\ker A^+$ and $\operatorname{im}\beta$.\end{certificate}
\begin{proof}
$BR=0$ ensures that $Re$ is closed, so $\beta$ is defined. $j$ injective and $\ker p$ are exactly all $(u,0)$ , $Au=0$ , and therefore $\operatorname{im}j=\ker p$ . If $(u,e)\in\ker A^+$, then $Au+Re=0$, so $\beta(e)=[Re]=0$. In reverse, $\beta(e)=0$ gives $Re=Au_e$, and from $A^+(-u_e,e)=0$ we get $e\in\operatorname{im}p$. So $\operatorname{im}p=\ker\beta$. Theorem ~ \ref{hgov:thm-extension} gives $\operatorname{im}\beta=\ker\iota_*$ and $\iota_*$ surjective. Each adjacent composite in the exact column is therefore zero, and each kernel has a listed preimage.

From the rank-zero degree formula, \[
 \dim\ker\beta=\dim E-r_R,\qquad
 \dim\ker A^+=\dim\ker A+\dim\ker\beta.
\] and then $\iota_*$ surjective and its core dimension $r_R$ are used to obtain two dimensional equations. The degree-boundary $B$ remains, so the zero-degree homology is equal. When $\dim E=r_R$ is $\ker\beta=0$, then $\ker A^+=\operatorname{im}j$, $j$ are isomorphic. Finally define $\bar T([x])=Tx$ and $\bar T^+([x])=Tx$. Change the representative elements to increase $Au$ or $Au+Re$ respectively, and their $T$ readings are both zero, so the two mappings are well-defined and satisfy the listed commutation relation. The deduction chain is \[
 \begin{gathered}
 BR=0\Rightarrow\beta(e)=[Re]\ \text{well-defined},\\
 Re=Au_e
 \Rightarrow A^+(-u_e,e)=0
 \Rightarrow\ker\beta=\operatorname{im}p,\\
 x=Au+Re
 \Rightarrow[x]_C=\beta(e)
 \Rightarrow\ker\iota_*=\operatorname{im}\beta,\\
 \operatorname{Exact}\land\operatorname{FiniteDimension}
 \Rightarrow\operatorname{RankIdentities},\\
 TA=TR=0\Rightarrow
 \operatorname{WellDefined}(\bar T,\bar T^+)
 \land\bar T^+\iota_*=\bar T.
 \end{gathered}
\]\end{proof}


\section { Specify the existence and minimum of expansion in the allowed direction }\label{exp:hom:restricted-directions}


Take a registered linear subspace $W_{\mathrm{adm}}\subseteq\ker B\cap\ker T$ and require each new operation to satisfy $\operatorname{im}R\subseteq W_{\mathrm{adm}}$. This linear model allows the synthesis of directions with arbitrary domain coefficients; non-negative constraints on the range of motion are specified separately in the definition ~ \ref{hgov:def-executable}. The target is still the fixed subspace $Z\subseteq\ker B\cap\ker T$.

\begin{theorem}
[Image Criterion of Allowed Directions and Restricted Minimum Expansion]\label{exp:hom:restricted-minimum}
There are finite dimensions $E$ and allowed $R:E\to W_{\mathrm{adm}}$ such that each $z\in Z$ has a preimage $z=Au+Re$, if and only if \[
 \pi(Z)\subseteq\pi(W_{\mathrm{adm}})
 \quad\Longleftrightarrow\quad
 Z\subseteq U+W_{\mathrm{adm}}.
\] this condition holds, the smallest $\dim E$ is equal to $r_Z=\dim\pi(Z)$ and enables $\operatorname{rank}(\pi R)=r_Z$. When the condition fails, there exist $z_*\in Z$ and linear functional $\lambda:C_1\to\mathbb k$ such that \[
 \lambda(U)=0,\qquad
 \lambda(W_{\mathrm{adm}})=0,\qquad
 \lambda(z_*)\ne0.
\] If all given matrices and subspace bases are on $\mathbb Q$, the above original image, minimum expansion and veto functional can all be constructed by finite rational elimination.\end{theorem}
\begin{certificate}
[Allow base, target quotient base, and deny readings]Save the basis of $W_{\mathrm{adm}}$ and $Z$ , $U=\operatorname{im}A$ , and column-wise $Bw=Tw=0$ . Successful branch selection $z_1,\ldots,z_{r_Z}$ makes $\pi(z_j)$ the base of $\pi(Z)$; commits $w_j\in W_{\mathrm{adm}}$ and $u_j\in C_2$ for each $j$ makes $z_j=Au_j+w_j$ , and define $Re_j=w_j$. The failure branch saves zero readings on $z_*$, $\lambda$ and all allowed bases.\end{certificate}
\begin{proof}
If $z=Au+Re$ and $\operatorname{im}R\subseteq W_{\mathrm{adm}}$, then $z\in U+W_{\mathrm{adm}}$; use $\pi$ to get the necessary inclusions listed. The two inclusions are equivalent to each other because $\pi(z)=\pi(w)$ is exactly equivalent to $z-w\in U$.

If the inclusion is true, select the above $w_j,u_j$ one by one. For any $z\in Z$, there is a coefficient $\alpha_j$ that satisfies $\pi(z)=\sum_j\alpha_j\pi(z_j)=\sum_j\alpha_j\pi(w_j)$, so there is $u$ such that $z=Au+R\alpha$ exists. Each $w_j$ belongs to the allowed subspace, so $R$ satisfies both allowability and compatibility. Its quotient image consists of the independent $\pi(w_j)$, so $\operatorname{rank}(\pi R)=\dim E=r_Z$. Theorem ~ \ref{hgov:thm-minimal-extension} gives the same lower bound for all candidates. When $r_Z=0$ , take $E=0$ , and $Z\subseteq U$ .

If failure is included, take $z_*\in Z\setminus(U+W_{\mathrm{adm}})$. First take the basis of $U+W_{\mathrm{adm}}$, then add $z_*$, and expand this set of linearly independent vectors into the basis of $C_1$. Let $\lambda$ be zero on the basis of the previous subspace, one on $z_*$, and zero on the remaining complementary basis to obtain the required rejection. Under the rational basis, the same process of selecting the basis and finding the original image is realized by finite elimination. Its predicate chain is \[
 \begin{gathered}
 \operatorname{Allowed}(R)\land
 Z\subseteq U+\operatorname{im}R
 \Rightarrow Z\subseteq U+W_{\mathrm{adm}},\\
 Z\subseteq U+W_{\mathrm{adm}}
 \land\operatorname{QuotientBasis}(\pi(z_j))
 \Rightarrow\exists R:\dim E=r_Z
 \land\forall z\in Z\ \exists(u,e):z=Au+Re,\\
 \lambda(U+W_{\mathrm{adm}})=0\land\lambda(z_*)\ne0
 \Rightarrow
 \forall R\ \bigl(\operatorname{Allowed}(R)
       \Rightarrow z_*\notin U+\operatorname{im}R\bigr).
 \end{gathered}
\]\end{proof}


\section { 3D restoration illustration with orthogonal distances }\label{hgov:three-dimensional}


\begin{proposition} [One new direction resolves one independent obstruction]
\label{hgov:prop-three-dimensional}
Over $\mathbb Q$, take
\[
 B=\begin{pmatrix}1&1&1\end{pmatrix},\quad
 A=p=\begin{pmatrix}1\\-1\\0\end{pmatrix},\quad
 \omega=\begin{pmatrix}1\\0\\-1\end{pmatrix},\quad T=B.
\].
The original complex has $H_1(C^-)\cong\mathbb Q$ and $[\omega]\ne0$. After adding one column $R=\omega$,
\[
 A^+=
 \begin{pmatrix}
 1&1\\-1&0\\0&-1
 \end{pmatrix},\qquad
 H_1(C^+)=0.
\].
Under the standard Euclidean inner product, the minimum residual of the original repair is $\sqrt{3/2}$, while the extended minimum residual is zero. One added direction attains the lower bound of Theorem ~ \ref{hgov:thm-minimal-extension}.
\end{proposition}
\begin{certificate} [Original obstruction, new preimage, and rank]
Retain
\[
 \begin{gathered}
 BA=0,\quad B\omega=0,\quad
 \lambda=\begin{pmatrix}0\\0\\1\end{pmatrix},\quad
 \lambda^{\mathsf T}A=0,\quad
 \lambda^{\mathsf T}\omega=-1,\\
 A^+\begin{pmatrix}0\\1\end{pmatrix}=\omega,\quad
 \operatorname{rank}B=1,\quad
 \operatorname{rank}A=1,\quad
 \operatorname{rank}A^+=2.
 \end{gathered}
\].
The orthogonal projection and residual are
\[
 P\omega=\frac12p,\qquad
 r=\omega-\frac12p
   =\begin{pmatrix}1/2\\1/2\\-1\end{pmatrix},\qquad
 p^{\mathsf T}r=0,\quad r^{\mathsf T}r=\frac32.
\].
\end{certificate}
\begin{proof}
$\lambda^{\mathsf T}A=0$ and $\lambda^{\mathsf T}\omega=-1$ give a rejection certificate for the original equation. The rank-nullity formula gives $\dim\ker B=2$, hence
\[
 \dim H_1(C^-)=2-1=1,\qquad
 \dim H_1(C^+)=2-2=0.
\].
Since $p^{\mathsf T}\omega=1$ and $p^{\mathsf T}p=2$, the projection coefficient is $1/2$; Theorem ~ \ref{hgov:thm-projection} gives the residual value. Take $Z=\mathbb Q\omega$. Its image $\pi(Z)$ is one-dimensional, so one additional column attains the lower bound. The equation-by-equation causal chain is
\[
 \lambda^{\mathsf T}A=0\wedge
 \lambda^{\mathsf T}\omega\ne0
 \Longrightarrow[\omega]_{C^-}\ne0
 \Longrightarrow r_Z=1
 \Longrightarrow
 \bigl(\dim E\ge1\bigr),
\].
\[
 R=\omega\wedge BR=TR=0
 \Longrightarrow A^+(0,1)=\omega
 \Longrightarrow[\omega]_{C^+}=0.
\]
\end{proof}

\section { The combination of forward cone, precondition and finite execution certificate }\label{hgov:executable}


\begin{definition}
[executable repair word]\label{hgov:def-executable}
Let $L=[a_1\ \cdots\ a_m]$ in the real substrate, the state field is $X\subseteq\mathbb R^d$, where $d$ is the row number of $L$. The action of class $j$ has non-negative amplitude $t$, predicate $\operatorname{Pre}_j(x,t,b)$, cost $c_j(x,t)\ge0$ and actual successor \[
 F_j(x,t)=x+ta_j.
\]. The predicate includes $F_j(x,t)\in X$ and $c_j(x,t)\le b$. The repair word $w=((j_1,t_1),\ldots,(j_N,t_N))$ starting from $(x_0,b_0)$ satisfies \[
 \begin{aligned}
 &t_n\ge0,\quad
 \operatorname{Pre}_{j_n}(x_{n-1},t_n,b_{n-1}),\\
 &x_n=F_{j_n}(x_{n-1},t_n),\qquad
 b_n=b_{n-1}-c_{j_n}(x_{n-1},t_n).
 \end{aligned}
\] and the target is changed to $\omega=x_N-x_0$. The algebraically reachable set is $\operatorname{im}L$, the non-negative net action set is $\mathcal K_L=\{Lu:u\in\mathbb R_{\ge0}^m\}$, and the actual reachable set consists of the net actions of the above legal finite words.\end{definition}


\begin{lemma}
[Inclusion of executability, positive cones and linear solvability]\label{hgov:lem-cone}
For each initial state and budget, \[
 \mathcal R_{\mathrm{exec}}(x_0,b_0)
 \subseteq\mathcal K_L\subseteq\operatorname{im}L.
\] if a column vector $\lambda$ satisfies $\lambda^{\mathsf T}L\ge0$ (component-wise) and $\lambda^{\mathsf T}\omega<0$, then $\omega\notin\mathcal K_L$.\end{lemma}
\begin{certificate}
[Amplitude Aggregation and Positive Cone Vetoing]Save $u_j=\sum_{\{n:j_n=j\}}t_n$ for legal words. Save all column checks for $\lambda^{\mathsf T}a_j\ge0$, and $\lambda^{\mathsf T}\omega<0$ against forward cone veto.\end{certificate}
\begin{proof}
According to the definition $u_j\ge0$, and the stepwise addition gives \[
 x_N-x_0=\sum_{n=1}^Nt_na_{j_n}
 =\sum_{j=1}^mu_ja_j=Lu.
\], so two levels of inclusion are obtained. If $\omega=Lu$ and $u\ge0$, then $\lambda^{\mathsf T}\omega
 =\sum_ju_j\lambda^{\mathsf T}a_j\ge0$, contradicting the given strictly negative values. Therefore \[
 \operatorname{Executed}(w)
 \Longrightarrow\exists u\ge0\ (Lu=\omega)
 \Longrightarrow\exists u\in\mathbb R^m\ (Lu=\omega),
\] and $\lambda$ listed refute the first existence proposition.\end{proof}


\begin{theorem}
[Finite prefix preservation of legal repair words]\label{hgov:thm-execution}
If the words defining ~ \ref{hgov:def-executable} gradually pass the pre-check, $x_0\in X$, $b_0\ge0$, and $BL=0$, $TL=0$, then for $0\le n\le N$, \[
 x_n\in X,\quad b_n\ge0,\quad
 Bx_n=Bx_0,\quad Tx_n=Tx_0.
\] If there is another task represented by $q:X\to Q$, and each step actually used satisfies the \[
 qF_{j_n}(x_{n-1},t_n)
 =\bar F_{j_n}(qx_{n-1},t_n),
\] definition $\bar x_0=q(x_0)$, and let $\bar x_n=\bar F_{j_n}(\bar x_{n-1},t_n)$, then the corresponding abstract execution satisfies $\bar x_n=qx_n$. When the target set satisfies $G=q^{-1}(\bar G)$, $\bar x_N\in\bar G$ is equivalent to $x_N\in G$.\end{theorem}
\begin{certificate}
[Prefix induction and task blocks]Step by step save the actual status, amplitude, precondition, cost and budget balance; save the $Ba_j=Ta_j=0$ of each column; when using task representation, also save the exchange equation and target preimage equation of the actual call point. The induction predicate is established in the previous $n$ step for all listed conclusions.\end{certificate}
\begin{proof}
The zero step is established by the input. Assume that the conclusion is established at $n-1$. precondition gives $x_n\in X$ and $b_n=b_{n-1}-c_{j_n}\ge0$ . There is also \[
 Bx_n=Bx_{n-1}+t_nBa_{j_n}=Bx_0,\qquad
 Tx_n=Tx_{n-1}+t_nTa_{j_n}=Tx_0.
\]. If $\bar x_{n-1}=qx_{n-1}$, then \[
 \bar x_n=\bar F_{j_n}(\bar x_{n-1},t_n)
 =\bar F_{j_n}(qx_{n-1},t_n)
 =qF_{j_n}(x_{n-1},t_n)=qx_n.
\]. Therefore, $P(0)\wedge\forall n\le N\,[P(n-1)\Rightarrow P(n)]$ derives all finite prefix conclusions, and the target pre-image formula then gives the final state readback.\end{proof}


\begin{remark}
[verifiable effect of precondition]\label{hgov:rem-preconditions}
Taking the integer state field $X=\mathbb N^2$, the two inseparable action increments are $a_1=(-1,2)^{\mathsf T}$ and $a_2=(1,-1)^{\mathsf T}$. Each action requires the consumed coordinate to be at least one. Starting from $x_0=(0,0)^{\mathsf T}$, $\omega=(0,1)^{\mathsf T}=a_1+a_2$ belongs to the non-negative integer net action set; both initial preconditions are false, and the legal finite words are only empty words. This set of data registers the net effect equation and prefix-by-prefix legality as two independently verifiable fields. For the new direction of proposition ~ \ref{hgov:prop-three-dimensional}, take $X=\mathbb R_{\ge0}^3$ and $x_0=(0,0,1)^{\mathsf T}$, and the action $x\mapsto x+\omega$ with amplitude one requires $x_3\ge1$; Cost one and budget one give a legal one-step final state. $(1,0,0)^{\mathsf T}$.\end{remark}


\section { finite generates the complete dual of a positive cone certificate }\label{exp:hom:cone-alternative}


This section fixes the real matrix $L=[a_1\ \cdots\ a_m]\in\mathbb R^{d\times m}$. The column number $m$ is part of the certificate and the allowed amplitude is a real non-negative value. The finite column closes the positive cone, and accordingly the veto condition of Lemma ~ \ref{hgov:lem-cone} is promoted to a necessary and sufficient alternative result.

\begin{theorem}
[non-negative net-effect preimage or detached certificate]\label{exp:hom:finite-farkas}
For any $\omega\in\mathbb R^d$, exactly one of the following two situations is true: \[
 \begin{array}{ll}
 \mathrm{(P)}&
 \exists u\in\mathbb R_{\ge0}^m,\quad Lu=\omega,\\[2pt]
 \mathrm{(D)}&
 \exists\lambda\in\mathbb R^d,\quad
 L^{\mathsf T}\lambda\ge0,\quad
 \lambda^{\mathsf T}\omega<0.
 \end{array}
\] If $\mathrm{(P)}$ is true, it can be obtained that the columns corresponding to the non-zero support of $u$ are linearly independent; therefore, the required number of non-zero amplitudes does not exceed $\operatorname{rank}L$. If $\mathrm{(D)}$ is established, $p$ can be the closest point from $\omega$ to $\mathcal K_L$, and take $\lambda=p-\omega$. At this time, \[
 \lambda^{\mathsf T}\omega=-\|\omega-p\|^2<0.
\]\end{theorem}
\begin{certificate}
[Non-negative magnitude or column-wise separation]The original image branch saves $u_j\ge0$ and $Lu=\omega$, and can save the non-zero-order subformula of the support sequence as an independence witness. The dual branch saves all $m$ checks for $\lambda^{\mathsf T}a_j\ge0$ and $\lambda^{\mathsf T}\omega<0$. The closest point construction also holds $p\in\mathcal K_L$ , $(\omega-p)^{\mathsf T}p=0$ , and $(\omega-p)^{\mathsf T}a_j\le0$ .\end{certificate}
\begin{proof}
First prove that finite generates a closed positive cone $\mathcal K_L$. Given $x=Lu$ and $u\ge0$, if the columns on the positive support set are linearly related, obtain the non-zero relationship $Lv=0$ in the support set, and replace $v$ with $-v$ if necessary, so that at least one $v_j>0$ . Setting \[
 \tau=\min_{\{j:v_j>0\}}\frac{u_j}{v_j}>0,
 \qquad u'=u-\tau v.
\] For the coordinates of $v_j>0$, the minimum value definition gives $u'_j\ge0$; for the coordinates of $v_j\le0$, there is $u'_j\ge u_j\ge0$. At least one original positive coordinate becomes zero, and $v_j=0$ is outside the support set. So $Lu'=Lu$, the positive support is strictly reduced. By induction on the support size, independent column representation is obtained in finite steps.

For each independent column set $I$ , the map $L_I:\mathbb R^{I}\to\operatorname{im}L_I$ has a continuous linear inverse. If $L_Iu^{(n)}$ converges and $u^{(n)}\ge0$ , its limit belongs to the closed subspace $\operatorname{im}L_I$ , and the inverse mapping gives that the coefficient limit is still non-negative component-wise. So $L_I\mathbb R_{\ge0}^{I}$ is closed. There are only finitely many sets of independent columns, and together with the empty set representing zero, $\mathcal K_L$ is the finite union of these closed sets and is therefore closed. It also immediately satisfies convexity by linearity and nonnegativity.

If $\omega\notin\mathcal K_L$, take the nearest point $p$ on the closed convex set. Its existence restricts the minimization to $\{z\in\mathcal K_L:\|z-\omega\|\le\|\omega\|\}$: the set contains zeros, is closed, and is bounded, so it is compact in finite-dimensional space. Order $r=\omega-p$. For any $z\in\mathcal K_L$, the line segments $p+t(z-p)$ and $0\le t\le1$ remain in the forward cone, and the minimum right derivative gives \[
 r^{\mathsf T}(z-p)\le0.
\]. Taking $z=0$ and $z=2p$ respectively, we get $r^{\mathsf T}p=0$; then take $z=a_j$ to get $r^{\mathsf T}a_j\le0$. Therefore, $\lambda=-r$ satisfies \[
 L^{\mathsf T}\lambda\ge0,\qquad
 \lambda^{\mathsf T}\omega
 =-r^{\mathsf T}(p+r)=-\|r\|^2<0.
\], which constructs $\mathrm{(D)}$. If $\omega\in\mathcal K_L$ , its definition gives $\mathrm{(P)}$ . The mutual exclusivity of the two branches comes from $u\ge0\Rightarrow\lambda^{\mathsf T}Lu\ge0$. The complete deduction chain is \[
 \begin{gathered}
 \operatorname{FiniteColumns}(L)
 \Rightarrow\operatorname{IndependentSupportReduction}
 \Rightarrow\operatorname{ClosedConvex}(\mathcal K_L),\\
 \omega\notin\mathcal K_L
 \Rightarrow\operatorname{NearestPoint}(p)
 \Rightarrow\mathrm{(D)},\\
 \omega\in\mathcal K_L\Rightarrow\mathrm{(P)},\qquad
 \mathrm{(P)}\land\mathrm{(D)}\Rightarrow 0\le\lambda^{\mathsf T}\omega<0.
 \end{gathered}
\]. The last item is the contradiction between non-negativity and strict negativity of the same reading.\end{proof}


\section { Budget slack variable and executable linear cost interface }\label{exp:hom:budget-cone}


Let $c=(c_1,\ldots,c_m)^{\mathsf T}\ge0$ be the fixed unit amplitude cost, and $b\ge0$ be the given budget. The fee stipulation for this section is $c_j(x,t)=c_jt$, so the total fee for a word is determined by the aggregation range. Register the unconsumed budget as slack variable $s\ge0$ to verify net effects and expenses at the same time.

\begin{theorem}
[Complete certificate and global execution implementation of budget constraints]\label{exp:hom:budget-alternative}
Exactly one of the following two situations is true: \[
 \begin{array}{ll}
 \mathrm{(P_b)}&
 \exists u\ge0,\quad Lu=\omega,\quad c^{\mathsf T}u\le b,\\[2pt]
 \mathrm{(D_b)}&
 \exists\lambda\in\mathbb R^d,\ \mu\ge0,\quad
 L^{\mathsf T}\lambda+\mu c\ge0,\quad
 \lambda^{\mathsf T}\omega+\mu b<0.
 \end{array}
\] If the state field is $X=\mathbb R^d$, each action is defined as $F_j(x,t)=x+ta_j$, and its precondition is exactly the current budget of $t\ge0$ and $c_jt\le$, then $\mathrm{(P_b)}$ It is also equivalent to budgeting admissible execution of a finite word $b$ from any given $x_0$ and getting the net effect of $\omega$ . The preimage $u$ can be implemented in column number order, using the amplitude $u_j$ once per column.\end{theorem}
\begin{certificate}
[Augmented matrix and prefix-by-prefix balance]Save \[
 \widetilde L=
 \begin{pmatrix}L&0\\c^{\mathsf T}&1\end{pmatrix},
 \qquad
 \widetilde\omega=\begin{pmatrix}\omega\\b\end{pmatrix}.
\] and save the original image branch $u\ge0$, $s=b-c^{\mathsf T}u\ge0$ and $\widetilde L(u,s)^{\mathsf T}=\widetilde\omega$. The dual branch holds all augmented column checks for $(\lambda,\mu)$. The global implementation branch also saves \[
 x_j=x_0+\sum_{i=1}^{j}u_i a_i,\qquad
 b_j=b-\sum_{i=1}^{j}c_i u_i,\qquad 0\le j\le m.
\]\end{certificate}
\begin{proof}
Block-by-block multiplication gives \[
 \widetilde L\begin{pmatrix}u\\s\end{pmatrix}
 =\widetilde\omega
 \Longleftrightarrow
 Lu=\omega\ \land\ c^{\mathsf T}u+s=b.
\] so the non-negative augmented preimage is exactly equivalent to $\mathrm{(P_b)}$. Applying the theorem ~ \ref{exp:hom:finite-farkas} to $\widetilde L$ , the nonnegative reading of the last column of the dual is $\mu\ge0$ and the readings of the remaining columns are $\lambda^{\mathsf T}a_j+\mu c_j\ge0$; the target reading is $\lambda^{\mathsf T}\omega+\mu b<0$. So we get $\mathrm{(D_b)}$, which is alternative.

Under the listed global implementation conditions, if $u$ is the budget preimage, each prefix satisfies \[
 0\le\sum_{i=1}^{j}c_i u_i
 \le\sum_{i=1}^{m}c_i u_i\le b.
\], so $b_j\ge0$ and $c_ju_j\le b_{j-1}$, and the precondition is passed in step $j$. Mathematical induction gives all $m$ prefix methods, and the net effect of the end point is $Lu=\omega$. In reverse, any legal word is aggregated into $u_j=\sum_{\{n:j_n=j\}}t_n$; the linear reduction of the cost gives the total cost $c^{\mathsf T}u\le b$, so $\mathrm{(P_b)}$ is obtained. The empty word corresponds to $u=0$. The predicate causal chain is \[
 \begin{gathered}
 Lu=\omega\land c^{\mathsf T}u\le b\land u\ge0
 \Longleftrightarrow
 \exists s\ge0:\widetilde L(u,s)^{\mathsf T}
                   =\widetilde\omega,\\
 \operatorname{FiniteColumns}(\widetilde L)
 \Rightarrow\mathrm{(P_b)}\ \text{and}\ \mathrm{(D_b)}
          \ \text{\shortstack[c]{\scriptsize exactly one\\[-1pt]\scriptsize holds}},\\
 \mathrm{(P_b)}\land\operatorname{GlobalLinearActions}
 \Rightarrow\bigwedge_{j=0}^{m}(b_j\ge0)
 \Rightarrow\operatorname{Executed}(w)\land(x_m-x_0=\omega).
 \end{gathered}
\]\end{proof}


\begin{remark}
[3D illustration of directions and budget readings]\label{exp:hom:budget-example}
In the proposition ~ \ref{hgov:prop-three-dimensional}, let $L=[p\ \omega]$ , $c=(0,1)^{\mathsf T}$ . For the target $-\omega$, take $\lambda=(0,0,-1)^{\mathsf T}$, then \[
 \lambda^{\mathsf T}L=(0,1),\qquad
 \lambda^{\mathsf T}(-\omega)=-1.
\]. So the extended signed linear image contains $-\omega$, and the target has dual veto when non-negative amplitudes are allowed.

For target $\omega$ and budget $b<1$, take $\lambda=(0,0,1)^{\mathsf T}$ and $\mu=1$, then when \[
 L^{\mathsf T}\lambda+\mu c=\begin{pmatrix}0\\0\end{pmatrix},
 \qquad
 \lambda^{\mathsf T}\omega+\mu b=b-1<0.
\] budgets $b=1$, $u=(0,1)^{\mathsf T}$ is the accepted certificate. Therefore, under the premise of linear cost specification and global action, the minimum cost of $\omega$ is exactly one. If you use the state field defined ~ \ref{hgov:def-executable} with other preconditions, the preimage and the minimum cost reading can be used as input to perform the check prefix-by-prefix.\end{remark}


\chapter{Formal Parameters, Arity Filtrations, and Order-by-Order Homotopy Obstructions}
\label{exp:ch:tower}


\section { order-by-order obstacles in total differential and high-order operations }\label{hgov:higher-obstruction}


Total homology and homotopy repair tower require that the new direction be compatible with all existing continuation formulas. The order and sign are fixed below, and the formal parameter order and operand number are dealt with respectively. Its law-space interface comes from the literature ~ \cite{S01,S38}. For the coderivation expression of high-order Lie operation, please refer to the literature ~ \cite{ExtLadaMarkl}.

\begin{definition} [The operator cochain complex and formal parameters]
\label{hgov:def-end}
Suppose $\mathbb k$ has characteristic zero, $C^\bullet$ is a graded space, and $D_0$ is a degree-one operator with $D_0^2=0$. On $\operatorname{End}^\bullet(C)$ set
\[
 \delta A=[D_0,A]
 =D_0A-(-1)^{|A|}AD_0.
\].
The parameter $t$ is a central formal variable of degree zero, and every $K_r$ has degree one. The order-$n$ truncation is
\[
 D^{(n)}(t)=D_0+\sum_{r=1}^nt^rK_r.
\].
\end{definition}

\begin{theorem}
[Necessary and sufficient conditions for order-by-order closure of the total differential]\label{hgov:thm-total-tower}
Let $n\ge1$ , and $\bigl(D^{(n-1)}(t)\bigr)^2=0\pmod{t^n}$ . Define the $n$-th order barrier \[
 O_n=\sum_{\substack{i+j=n\\i,j\ge1}}K_iK_j,
 \qquad O_1=0.
\] of degree two, then \[
 \delta O_n=0,
 \qquad
 \bigl(D^{(n)}(t)\bigr)^2=0\pmod{t^{n+1}}
 \Longleftrightarrow
 \delta K_n=-O_n.
\]. Therefore, the $n$-th order expansion exists if and only if $[O_n]=0\in H^2(\operatorname{End}C,\delta)$. Give $K_1,\ldots,K_N$ step by step to obtain the closed certificate up to the level $N$.\end{theorem}
\begin{certificate}
[Coefficients, fractions and order induction]Save $D_0^2=0$, the known operators $K_i$, coefficients $O_n$, and \[
 D_0O_n-O_nD_0=0,\qquad
 D_0K_n+K_nD_0=-O_n.
\] old order coefficients and new order coefficients of each order are retained one by one; the end point records the specific truncation order $N$.\end{certificate}
\begin{proof}
The definition ~ \ref{hgov:def-end} is an instance of Lemma ~ \ref{hgov:lem-hom} when the source and target are the same, so $\delta^2=0$. Set $A(t)=D^{(n-1)}(t)$. The associative law directly gives $A(t)A(t)^2-A(t)^2A(t)=0$. According to the assumption, the coefficients of $A(t)^2$ below $n$ are all zero, and its $n$-th order coefficient is $O_n$, so the $n$-th order of the listed identities is \[
 D_0O_n-O_nD_0=\delta O_n=0.
\]. New After $t^nK_n$, the square coefficient of order $n$ is \[
 D_0K_n+K_nD_0+O_n=\delta K_n+O_n.
\]. All products involving $K_n$ and old terms of positive parameter order have a parameter order greater than $n$, and the old low-order coefficients remain. Therefore, this order of closure is exactly equivalent to the listed filling equation. The closed cochain $O_n$ has the preimage $-K_n$ which is exactly equivalent to its class being zero. Taking $P(n)$ to be squared is zero modulo $t^{n+1}$, then \[
 \begin{gathered}
 P(0):D_0^2=0,\\
 P(n-1)\Longrightarrow\delta O_n=0,\\
 P(n-1)\wedge(\delta K_n=-O_n)
 \Longrightarrow P(n),\\
 P(0)\wedge
 \bigwedge_{n=1}^{N}(\delta K_n=-O_n)
 \Longrightarrow P(N).
 \end{gathered}
\] This gives the order-by-order mathematical induction certificate.\end{proof}


\section { Compatible closure of graded suspension and operands }\label{exp:hom:arity-grading}


The formal parameter order is measured by the multiplication of the central variables, and the input arity is measured by the insertion of the coderivation. The two satisfy $t^it^j=t^{i+j}$ and $a$ meta-operations respectively. After inserting the $b$ meta-operation, there are $a+b-1$ inputs. This section fixes the degree and sign of the antisymmetric representation to the symmetric coalgebraic representation, so that each of the two sets of step-by-order equations has a clear type.

\begin{lemma}
[Hover interface for antisymmetric rendering and symmetric coderivation]\label{hgov:lem-suspension}
Assume that the original grading space is $W$, the degree of $\ell_n:\Lambda^nW\to W$ is $2-n$, and the symbol convention is \ref{hgov:linfty:convention}. Let $V=W[1]$ , where $W[1]^j=W^{j+1}$ , and let $s:W\to V$ be the specified linear isomorphism of degree $-1$. Aligning the input definition \[
 q_n(sx_1\odot\cdots\odot sx_n)
 =(-1)^{\alpha_n(x)}s\ell_n(x_1,\ldots,x_n),
 \qquad
 \alpha_n(x)=\sum_{k=1}^{n}(n-k)|x_k|.
\] then $q_n:S^nV\to V$ is well-defined and the degree is one. Expand $q_n$ to coderivation $\widehat q_n$ by shuffling sum, let $Q=\sum_{n\ge1}\widehat q_n$ , then \[
 \operatorname{pr}_V Q^2(sx_1\odot\cdots\odot sx_n)
 =(-1)^{\alpha_n(x)}s\mathcal R_n(x_1,\ldots,x_n).
\] and therefore $Q^2=0$ if and only if all $\mathcal R_n=0$ .\end{lemma}
\begin{certificate}
[Shift times and shuffle symbols]Save the shuffle expansion of $s,s^{-1}$, input times, index $\alpha_n$, adjacent swap symbols and coderivation. Only finite terms for each fixed input element participate in the calculation.\end{certificate}
\begin{proof}
The number of outputs minus the number of inputs is $-1+(2-n)+n=1$. After swapping adjacent $x,y$, the change in exponent and the antisymmetric swap together give \[
 (-1)^{|x|+|y|}\bigl(-(-1)^{|x||y|}\bigr)
 =(-1)^{(|x|-1)(|y|-1)},
\], which is exactly the swap sign of $sx,sy$ in $S^nV$, so it is well defined. For the first $i$ inputs of the unpermuted inner layer and the outer layer $j=n+1-i$, the index of the two suspensions is \[
 \begin{aligned}
 &\sum_{k=1}^{i}(i-k)|x_k|
 +(j-1)\left(\sum_{k=1}^{i}|x_k|+2-i\right)
 +\sum_{k=i+1}^{n}(n-k)|x_k|\\
 &\hspace{2em}\equiv\alpha_n(x)+i(j-1)\pmod2.
 \end{aligned}
\]. The general shuffling consists of adjacent exchanges; the previous commutation relation will successively change the Koszul symbol of the suspension space to $\chi(\sigma;x)$ on $W$. differ from the corresponding index. Therefore, each term has a common factor $(-1)^{\alpha_n(x)}s$, and the remaining terms are exactly $(-1)^{i(j-1)}\chi(\sigma;x)$ multiplied by the corresponding $\ell_j\ell_i$. Summing gives the equation listed. The square of coderivation of degree one is still coderivation , and is uniquely determined by its projection onto the remainder generator $V$, so the ternary projection to zero is equivalent to $Q^2=0$. Its predicate chain is \[
 \begin{gathered}
 \operatorname{ShiftDegree}\land\operatorname{SwapSign}\land
 \operatorname{ShuffleExpansion}\\
 \Rightarrow\operatorname{CorestrictionIdentity}
 \Rightarrow(Q^2=0\Longleftrightarrow\forall n\ge1,\mathcal R_n=0).
 \end{gathered}
\]\end{proof}


\begin{corollary}
[homotopy algebra imputation in operand filtering]\label{hgov:cor-arity}
Let $V$ be the graded space, $S^c(V)=\bigoplus_{m\ge1}S^m(V)$ be the reduced graded symmetric coalgebra, and use Koszul notation when exchanging inputs. The only extension of degree one mapping $q_m:S^mV\to V$ is coderivation $\widehat q_m$. Suppose $q_1^2=0$ , and the square of $Q_{<n}=\sum_{m=1}^{n-1}\widehat q_m$ is zero when the arity is less than $n$ , $n\ge2$ . Denote $q_1^{(n)}$ as the differential induced by $q_1$ on $S^nV$, let \[
 \begin{aligned}
 \delta_nf&=q_1f-(-1)^{|f|}f q_1^{(n)},\\
 \omega_n&=
 \operatorname{pr}_{V}\,Q_{<n}^{\,2}\big|_{S^nV}.
 \end{aligned}
\] be $\delta_n^2=0$, $\delta_n\omega_n=0$, and $n$ The necessary and sufficient condition for the existence of meta-operation is \[
 \bigl(\exists q_n\in\operatorname{Hom}^1(S^nV,V):
 \delta_nq_n=-\omega_n\bigr)
 \Longleftrightarrow
 [\omega_n]=0
 \in H^2\bigl(\operatorname{Hom}(S^nV,V),\delta_n\bigr).
\]\end{corollary}
\begin{certificate}
[Separate registration of yuan and degree]Save the degree of each $q_m$, input arity $m$, coderivation extended Koszul symbol and closed low arity relationship. This stage saves $\omega_n$ with degree two and $q_n$ with degree one, and verifies $q_1q_n+q_nq_1^{(n)}=-\omega_n$ base by base.\end{certificate}
\begin{proof}
On tensor powers, $q_1^{(n)}$ is the signed sum acting on each input position; each co-position square cancels by $q_1^2=0$, and the two orders at different positions cancel out due to the opposite Koszul sign of degree one. The effect drops to a graded symmetric power, so $(q_1^{(n)})^2=0$. So $(S^nV,q_1^{(n)})$ and $(V,q_1)$ are both cochain complex, and $\delta_n^2=0$ is established by the lemma ~ \ref{hgov:lem-hom}. The two elements are coderivation composites of $a,b$, and the input element of its core component is $a+b-1$. Therefore, after adding $\widehat q_n$, the square core of the element $n$ is exactly \[
 \omega_n+q_1q_n+q_nq_1^{(n)}
 =\omega_n+\delta_nq_n.
\]. The composite containing $q_n$ and $q_j$ and $j\ge2$ has an element of at least $n+1$. , the old low arity relationship is maintained. In the $n$ core of the identity $[Q_{<n},Q_{<n}^{\,2}]=0$, all terms involving the squared core below the $n$ element are zero, leaving $\delta_n\omega_n=0$. So the intermediate equivalence of \[
 \begin{gathered}
 \operatorname{Closed}_{<n}(Q_{<n})
 \Longrightarrow\delta_n\omega_n=0,\\
 \delta_n\omega_n=0
 \Longrightarrow
 \left[
 \bigl(\exists q_n:\delta_nq_n=-\omega_n\bigr)
 \Longleftrightarrow[\omega_n]=0
 \right],\\
 \operatorname{Closed}_{<n}(Q_{<n})
 \wedge(\delta_nq_n=-\omega_n)
 \Longrightarrow
 \operatorname{Closed}_{\le n}(Q_{<n}+\widehat q_n).
 \end{gathered}
\] is established by the kernel-image quotient definition, and the last step uses the constructed padding.\end{proof}


\section { finite Defibration and selection parameters of repair tower }\label{exp:hom:tower-fibres}


Let \[
 \mathcal A^r=\operatorname{End}^r(C),\qquad
 Z^r_\delta=\ker(\delta:\mathcal A^r\to\mathcal A^{r+1}),\qquad
 B^r_\delta=\delta(\mathcal A^{r-1}).
\] be the fixed integer $N\ge0$ , define the finite fix tower \[
 \mathcal T_N=
 \left\{(K_1,\ldots,K_N)\in(\mathcal A^1)^N:
 \delta K_r+\sum_{\substack{i+j=r\\i,j\ge1}}K_iK_j=0
 \text{ for }1\le r\le N\right\}.
\] by $\mathcal T_0=\{()\}$ , and delete the last coefficient with $\tau_n:\mathcal T_n\to\mathcal T_{n-1}$ . These definitions only involve $N$ parameter orders and finite products therein.

\begin{theorem} [Empty fibers, affine fibers, and the first-order homotopy quotient of the truncation map]
\label{exp:hom:tower-fibre}
Fix $\mathbf K_{<n}=(K_1,\ldots,K_{n-1})\in\mathcal T_{n-1}$ and, following Theorem ~ \ref{hgov:thm-total-tower}, compute $O_n$. If $[O_n]\ne0$, then
$\tau_n^{-1}(\mathbf K_{<n})=\varnothing$.
If $K_n^0$ has been obtained with $\delta K_n^0=-O_n$, then
\[
 \tau_n^{-1}(\mathbf K_{<n})
 \cong K_n^0+Z^1_\delta.
\]
where the relation on the same fiber
\[
 K_n\sim K'_n
 \quad\Longleftrightarrow\quad
 K'_n-K_n=\delta S
 \text{ for some }S\in\mathcal A^0
\]
gives a quotient solution set on which $H^1(\mathcal A,\delta)$ acts freely and transitively. After choosing $K_n^0$, the quotient coordinate is $[K_n-K_n^0]$.

If $C$ has only finite multiple non-zero divisions and each division has finite dimension, after fixing the basis, $\mathcal T_N$ is given by a finite group of polynomial equations of degree at most two. For a given rational low-order branch, the existence of its next fiber, a preimage and all linear free parameters can be verified by finite rational elimination.\end{theorem}
\begin{certificate}
[Branches, obstacles, preimages and nuclear bases]Record the low-order concrete coefficients, all closed relations, the matrix of $O_n$, and the matrix of $\delta:\mathcal A^1\to\mathcal A^2$. Non-empty fibers save the basis of $K_n^0$ and $Z^1_\delta$; the quotient solution set additionally saves the basis of $B^1_\delta$ and its coordinates in $Z^1_\delta$. An empty fiber holds the row vectors $\lambda^{\mathsf T}$ , $\lambda^{\mathsf T}\delta=0$ and $\lambda^{\mathsf T}O_n\ne0$ .\end{certificate}
\begin{proof}
After fixing the low-order branch, $O_n$ is a known degree two element, and is closed by the theorem ~ \ref{hgov:thm-total-tower}. Therefore, this order coefficient must satisfy the linear equation $\delta K_n=-O_n$. If $K_n^0$ exists, the difference between the two solutions satisfies $\delta(K_n-K_n^0)=0$, otherwise any $Z\in Z^1_\delta$ gives the solution $K_n^0+Z$. This proves the fiber formula and emptiness for non-zero obstacles. $\delta^2=0$ guarantees that $\delta S$ belongs to $Z^1_\delta$; taking the class of the closed element modulo $B^1_\delta$ is the same as the difference-by-difference proof of the theorem ~ \ref{exp:hom:filler-torsor}, and obtains the listed free transfer effect.

Under the finite graded basis, each coordinate of $\delta K_r$ is a linear equation, and each coordinate of $K_iK_j$ is a quadratic equation; the number of coefficients and the number of orders are both finite, so a finite polynomial presentation is obtained. When the low-order branches are given, all quadratic terms have become definite right-hand ends, and this order is eliminated by the theorem ~ \ref{hgov:prop-elimination} finite. The right end of subsequent steps is recalculated from the selected coefficients. The predicate chain is \[
 \begin{gathered}
 \mathbf K_{<n}\in\mathcal T_{n-1}
 \Rightarrow \delta O_n=0,\\
 \delta O_n=0
 \Rightarrow
 \bigl(\tau_n^{-1}(\mathbf K_{<n})\ne\varnothing
       \Longleftrightarrow[O_n]=0\bigr),\\
 \delta K_n^0=-O_n
 \Rightarrow
 \tau_n^{-1}(\mathbf K_{<n})\cong K_n^0+Z^1_\delta,\\
 \operatorname{GivenRationalBranch}(\mathbf K_{<n})
 \Rightarrow\operatorname{FiniteLinearFibreCheck}.
 \end{gathered}
\]\end{proof}


\section { Covariant transport of form base change and obstacles }\label{exp:hom:obstruction-transport}


In the following, we first deal with the formal base change on the same substrate, and then deal with the structural mapping between the two operator algebras. All base swaps and structure mappings are saved. This way, low-order coefficients, obstacle classes, and primitives can be transported along the same witness.

\begin{theorem} [Finite inverses of formal basis changes, the first-order formula, and obstruction classes]
\label{exp:hom:formal-conjugation}
On $\mathbb k[t]/(t^{N+1})$ take
\[
 G(t)=I+\sum_{j=1}^{N}t^jS_j,\qquad |S_j|=0.
\].
Then $G$ is invertible, and
\[
 D'=GDG^{-1},\qquad (D')^2=GD^2G^{-1}
\]
gives a bijection on $\mathcal T_N$. If $G=I+t^rS+O(t^{r+1})$ and $1\le r\le N$, then
\[
 K'_i=K_i\ (i<r),\qquad K'_r=K_r-\delta S.
\].

Further, take $1\le n\le N$ and the fixed low-order closed branch $A=D_0+\sum_{i=1}^{n-1}t^iK_i$. The only one written as \[
 GAG^{-1}=A'+t^nB_n\pmod{t^{n+1}},
 \qquad
 A'=D_0+\sum_{i=1}^{n-1}t^iK'_i.
\] is $|B_n|=1$, $A'$ is low-order closed, and its $n$-th order obstacle satisfies \[
 O'_n=O_n-\delta B_n,\qquad [O'_n]=[O_n]\in H^2(\mathcal A,\delta).
\]. If $\delta K_n=-O_n$, then $K'_n=K_n+B_n$ satisfies $\delta K'_n=-O'_n$.\end{theorem}
\begin{certificate}
[Truncated inverse and current order basis coefficient]Save all degree zero coefficients of $G$ , let $H=G-I$ , and record \[
 G^{-1}=\sum_{a=0}^{N}(-H)^a\pmod{t^{N+1}}.
\] Save $GG^{-1}=G^{-1}G=I$ , the typed conjugate expansion, and the truncated $A'$ and $n$ order residues $B_n$. Verify $O'_n-O_n+\delta B_n=0$; accept branch save $K'_n=K_n+B_n$ padding.\end{certificate}
\begin{proof}
Each parameter order of $H$ is at least one, so $H^{N+1}=0\pmod{t^{N+1}}$. Finite geometry and multiplying both sides gives the inverses listed. The associative law and $G^{-1}G=I$ give the constant term of \[
 (GDG^{-1})(GDG^{-1})=GD^2G^{-1}.
\] $G$ is identity, so the constant term of $D'$ is still $D_0$; the inverse basis is $D=G^{-1}D'G$, so the repair tower is given Bijective. The first modification order only involves the constant terms $D_0$ and $t^rS$, so the coefficients of \[
 K'_r=K_r+SD_0-D_0S=K_r-\delta S,
\] lower than $r$ remain.

For the second part, the old square has the form
$A^2=t^nO_n+O(t^{n+1})$.
Because the constant coefficient of $G$ is $I$,
\[
 GA^2G^{-1}=t^nO_n+O(t^{n+1}).
\].
On the other side, $GAG^{-1}=A'+t^nB_n$ gives
\[
 (GAG^{-1})^2
 =(A')^2+t^n(D_0B_n+B_nD_0)
 \pmod{t^{n+1}}.
\].
Here $2n\ge n+1$, and products of $B_n$ with all positive-order old coefficients also lie in the discarded orders. Therefore the square coefficients of $A'$ below order $n$ vanish, while at order $n$ we have $O'_n+\delta B_n=O_n$. This proves preservation of the obstruction class and also gives
$\delta(K_n+B_n)=-O_n+\delta B_n=-O'_n$.
The predicate chain is
\[
 \begin{gathered}
 G-I\in t\mathcal A^0[t]
 \Rightarrow\operatorname{FiniteInverse}(G)
 \Rightarrow(D')^2=GD^2G^{-1},\\
 A^2=0\pmod{t^n}
 \Rightarrow(A')^2=0\pmod{t^n}
 \land O'_n=O_n-\delta B_n,\\
 O'_n=O_n-\delta B_n
 \Rightarrow[O'_n]=[O_n],\\
 \delta K_n=-O_n
 \Rightarrow\delta(K_n+B_n)=-O'_n.
 \end{gathered}
\].
\end{proof}

\begin{theorem}
[Obstacle and filling transportation under differential fractional representation]\label{exp:hom:dg-transport}
Suppose $(\mathcal A,\delta)$ and $(\mathcal A',\delta')$ are operator differential fractional algebras obtained from the zero-degree square operator $D_0,D'_0$. A zero linear mapping of given degree $F:\mathcal A\to\mathcal A'$ satisfies \[
 F(XY)=F(X)F(Y),\qquad
 F(\delta X)=\delta'F(X).
\] to define a low-order closed branch $K'_i=F(K_i)$. Then the target branch is also low-order closed, and each source of \[
 O'_n=F(O_n),\qquad
 [O'_n]=H^2(F)[O_n].
\] is filled with $K_n$ and the transport is filled with $F(K_n)$ for the target. If $H^2(F)$ injective , then the source padding of this level exists if and only if the target padding exists. If $F$ is a differential graded algebraic isomorphism, then it also gives the bijective of all this-order filled sets on both sides.

In particular, the co-chain isomorphism $g:(C,D_0)\to(C',D'_0)$ yields $F(X)=gXg^{-1}$, satisfying all conditions.\end{theorem}
\begin{certificate}
[Multiplication, Differentiation and Class Mapping]Save the matrix, multiplication-preserving form, and commutation with differentials form of $F$ in the graded basis used. When saving \[
 \delta'F(K_r)+F(O_r)=F(\delta K_r+O_r)=0.
\] order-by-order using reverse existence, additionally save the basis or rank witness of $H^2(F)$ whose core is zero. When using the bijective of all solutions, save $F^{-1}$ and its inverse on both sides. In the case of co-chain isomorphism, $g^{-1}$ and $D'_0g=gD_0$ are saved.\end{certificate}
\begin{proof}
Multiplication preservation and linearity give \[
 \sum_{\substack{i+j=n\\i,j\ge1}}F(K_i)F(K_j)
 =F\left(\sum_{\substack{i+j=n\\i,j\ge1}}K_iK_j\right)
 =F(O_n).
\] The low-order relation is still zero after applying $F$, so the target low-order branch is closed. The commutation with differentials formula maps the closed element and boundary to the closed element and boundary respectively, so $H^2(F)$ is well-defined, and has the listed formulas. If $\delta K_n=-O_n$, apply $F$ to get $\delta'F(K_n)=-O'_n$. In reverse, if the target padding exists, then $H^2(F)[O_n]=[O'_n]=0$; the injective property gives $[O_n]=0$, and then use the theorem ~ \ref{hgov:thm-total-tower} to construct the source padding. If $F$ is isomorphic, the same derivation is applied to $F^{-1}$ to obtain the bijective of all solution sets.

Finally, for $F(X)=gXg^{-1}$ , the associative law gives a multiplicative hold; aligning times $X$ has \[
 \begin{aligned}
 \delta'F(X)
 &=D'_0gXg^{-1}-(-1)^{|X|}gXg^{-1}D'_0\\
 &=g\bigl(D_0X-(-1)^{|X|}XD_0\bigr)g^{-1}
 =F(\delta X).
 \end{aligned}
\] so the co-chain isomorphism is an instance of the above transport. The effect of each condition is verified by the following chain: \[
 \begin{gathered}
 \operatorname{Multiplicative}(F)\land
 \operatorname{Differential}(F)
 \Rightarrow O'_n=F(O_n)
 \land[O'_n]=H^2(F)[O_n],\\
 \delta K_n=-O_n
 \Rightarrow\delta'F(K_n)=-O'_n,\\
 [O'_n]=0\land\operatorname{Injective}(H^2(F))
 \Rightarrow[O_n]=0
 \Rightarrow\exists K_n:\delta K_n=-O_n.
 \end{gathered}
\]\end{proof}


\section { First-order parameter change and slope of obstacle distance }\label{exp:hom:obstruction-slope}


This section takes the finite-dimensional real graded space $C$. Only finite multiple grades are non-zero, and $D_0$ is fixed. Parameter $s$ is located in the open interval $I\subseteq\mathbb R$. It is stipulated that $K_i(s)\in\mathcal A^1$ and $1\le i<n$ are coordinate-by-coordinate differentiable mappings, and each $s\in I$ satisfies all order relations lower than $n$. These finite matrix equations can thus be differentiated term by term. First-order changes are recorded using the commutative algebra $\mathbb R[\varepsilon]/(\varepsilon^2)$, where $\varepsilon$ has degree zero and is distinguished from the formal variable $t$.

\begin{theorem}
[First-order barrier primitives and observable residual slopes]\label{exp:hom:first-variation}
Fixed $n\ge2$ and $s_0\in I$, remembering that \[
 K_i^0=K_i(s_0),\qquad
 V_i=\frac{dK_i}{ds}(s_0),\qquad
 O_n^0=O_n(s_0).
\] has obtained this level of benchmark filling $\delta K_n^0=-O_n^0$. Define \[
 \dot O_n
 =\sum_{\substack{i+j=n\\i,j\ge1}}
       (V_iK_j^0+K_i^0V_j)\in\mathcal A^2.
\] then $\delta\dot O_n=0$. The following conditions are equivalent: \[
 \begin{gathered}
 \exists V_n\in\mathcal A^1:
 D_0+\sum_{i=1}^{n}t^i(K_i^0+\varepsilon V_i)
 \text{ \shortstack[c]{\scriptsize squared modulus} }(t^{n+1},\varepsilon^2)\text{ vanishes},\\
 \exists V_n\in\mathcal A^1:\delta V_n=-\dot O_n,\\
 [\dot O_n]=0\in H^2(\mathcal A,\delta).
 \end{gathered}
\] Fix the positive definite inner product on $\mathcal A^2$, let $P$ be the orthogonal projection to $\operatorname{im}(\delta:\mathcal A^1\to\mathcal A^2)$, and define \[
 e_n(s)=\min_{X\in\mathcal A^1}\|\delta X+O_n(s)\|,
 \qquad
 a_n=\|(I-P)\dot O_n\|.
\] then \[
 e_n(s)=|s-s_0|a_n+o(|s-s_0|),\qquad
 a_n=0\Longleftrightarrow[\dot O_n]=0.
\] if $a_n>0$ , when $\epsilon_0>0$ exists, make $0<|s-s_0|<\epsilon_0$ when $e_n(s)>0$ .\end{theorem}
\begin{certificate}
[Derivative matrices, dual coefficients and orthogonal components]Save each $K_i^0,V_i$ , the low-order relations and their derivatives, and $\delta K_n^0=-O_n^0$ . Save the finite product sum and closure of $\dot O_n$ item by item. The first-order acceptance branch saves the preimage equation of $V_n$; the rejection branch saves $(I-P)\dot O_n\ne0$ or the corresponding dual row. The distance field holds $P^2=P$, orthogonality, and $(I-P)O_n^0=0$, and records the inner product used and the specification of the parameter $s$.\end{certificate}
\begin{proof}
The coordinate-by-coordinate product rule of finite summation and matrix multiplication gives $\frac{d}{ds}O_n(s_0)=\dot O_n$ . Since $\delta$ is fixed and $\delta O_n(s)=0$ holds for all $s$, derivation gives $\delta\dot O_n=0$. Derivating each low-order coefficient relationship also yields that the first-order coefficient of its dual expansion is zero. The even number coefficient of this order is exactly \[
 \begin{aligned}
 \delta(K_n^0+\varepsilon V_n)
  +O_n^0+\varepsilon\dot O_n
 &=(\delta K_n^0+O_n^0)
   +\varepsilon(\delta V_n+\dot O_n)\\
 &=\varepsilon(\delta V_n+\dot O_n).
 \end{aligned}
\], so the even number closure is equivalent to $\delta V_n=-\dot O_n$. The closure property then gives it the equivalence to the cohomology class of zero.

Theorem ~ \ref{hgov:thm-projection} is applied to the fixed matrix $\delta$ to get $e_n(s)=\|(I-P)O_n(s)\|$ . Baseline padding guarantee $(I-P)O_n^0=0$. Differentiability then gives \[
 (I-P)O_n(s)
 =(s-s_0)(I-P)\dot O_n+o(|s-s_0|).
\] Using the inverse triangle inequality of the norm, we get the distance expansion listed. The image of $P$ is exactly the boundary space, so $a_n=0\Longleftrightarrow\dot O_n\in\operatorname{im}\delta$; combined with $\delta\dot O_n=0$, the class criterion is obtained. When $a_n>0$, select the neighborhood so that the residual norm is smaller than $\frac12 a_n|s-s_0|$, and get $e_n(s)\ge\frac12a_n|s-s_0|>0$. All conclusions are composed of the following predicate chain: \[
 \begin{gathered}
 \operatorname{Differentiable}(K_i)\land
 \operatorname{Closed}_{<n}(K_i(s))
 \Rightarrow\delta\dot O_n=0,\\
 \delta K_n^0=-O_n^0
 \Rightarrow
 \bigl(\operatorname{FirstOrderFill}
       \Longleftrightarrow\exists V_n:\delta V_n=-\dot O_n\bigr),\\
 \delta\dot O_n=0
 \Rightarrow
 \bigl(\operatorname{FirstOrderFill}
       \Longleftrightarrow[\dot O_n]=0
       \Longleftrightarrow a_n=0\bigr),\\
 a_n>0\land\operatorname{RemainderBound}
 \Rightarrow
 0<|s-s_0|<\epsilon_0\Rightarrow e_n(s)>0.
 \end{gathered}
\]\end{proof}


\begin{remark}
[Computable finite parameter branches]\label{exp:hom:parameter-matrix}
Take $C^0=\mathbb R v_0$, $C^1=\mathbb R v_1$, $C^2=\mathbb R v_2$, $D_0=0$, and sort by $(v_0,v_1,v_2)$. Let \[
 K_1(s)=
 \begin{pmatrix}0&0&0\\1&0&0\\0&s&0\end{pmatrix},
 \qquad
 O_2(s)=K_1(s)^2=
 \begin{pmatrix}0&0&0\\0&0&0\\s&0&0\end{pmatrix}.
\] and each $K_1(s)$ be of degree one, and all first-order relations $\delta K_1(s)=0$ are established. Here $\delta=0$ , so the second-order preimage exists if and only if $s=0$ ; where $s=0$ is preferable $K_2^0=0$ , the remaining choices are all $Z^1_\delta=\mathcal A^1$ . When the matrix unit $E_{20}:v_0\mapsto v_2$ is taken as unit length, \[
 \dot O_2(0)=E_{20},\qquad
 e_2(s)=|s|,\qquad a_2=1.
\] therefore this set of finite matrices simultaneously gives non-empty fibers, empty fibers, first-order vetoes and accurate obstacle distance slopes.\end{remark}


\begin{remark}
[Combination Interface for Structural Repair]\label{hgov:rem-interface}
The finite representation in this section registers the governance chain as \[
 \begin{gathered}
 \text{\shortstack[c]{\scriptsize Typed\\[-1pt]\scriptsize continuation}}
 \longrightarrow J\in\ker\delta
 \longrightarrow[J]\ \text{\shortstack[c]{\scriptsize and dual\\[-1pt]\scriptsize certificates}}\\
 \longrightarrow
 \text{\shortstack[c]{\scriptsize Compatible new\\[-1pt]\scriptsize directions and\\[-1pt]\scriptsize higher-order\\[-1pt]\scriptsize fillings}}
 \longrightarrow
 \text{\shortstack[c]{\scriptsize Admissible\\[-1pt]\scriptsize finite execution}}
 \longrightarrow
 \text{\shortstack[c]{\scriptsize Task readback\\[-1pt]\scriptsize and preservation\\[-1pt]\scriptsize certificates}}.
 \end{gathered}
\] homology class to describe the direction that the existing boundary image has not yet covered; the compatible extension is measured by the dimensionality of the task obstacle subspace in the original boundary quotient; the higher-order equation verifies the compatibility of the new operation with the existing structure; the prefix verification is performed to connect the algebraic solution to the actual path. Together, these fields give a describeable, repairable, and verifiable structure for heterogeneous evolution.\end{remark}


\chapter{A Complete Joint Construction of Negative Curvature Inputs and Finite Repair}
\label{exp:ch:joint}


\label{hgov:sec:joint}

This section configures standard hyperbolic geometry and finite repair complexes in the same task family. All geometric coefficients, task observations, old and new actions, obstacle readings and terminals are given by the following definitions; the required proofs are completed in this section.

\section { Geometric input and non-zero homological obstruction }

\begin{definition}
[Hyperbolic parameters and repair complex]\label{hgov:def:joint}
Fix $a>0$, let $H_a$ be the hyperbolic plane with constant cross-section curvature $-a^2$, take the unit speed geodesic $\gamma$ and the unit transverse parallel field $E$. Define \[
J(t)=\frac{\sinh(at)}aE(t),\qquad
\lambda(t)=\langle J(t),E(t)\rangle=\frac{\sinh(at)}a .
\] to take the repair change space $V=\mathbb R^3$, and fix the old and new chain complex of \[
\mathbf p=\begin{pmatrix}1\\-1\\0\end{pmatrix},\qquad
\mathbf r=\begin{pmatrix}1\\0\\-1\end{pmatrix},\qquad
B=\begin{pmatrix}1&1&1\end{pmatrix}.
\] as \[
C^-:\mathbb R\xrightarrow{A^-=\mathbf p}\mathbb R^3
\xrightarrow{B}\mathbb R,\qquad
C^+:\mathbb R^2\xrightarrow{A^+=[\,\mathbf p\ \mathbf r\,]}\mathbb R^3
\xrightarrow{B}\mathbb R.
\] respectively. The corresponding chain mapping $i:C^-\to C^+$ is $v\mapsto(v,0)$ in degree two, and the degree one and zero are identical. The bridge between the geometry read and the changes to be fixed is defined as $\omega(t)=\lambda(t)\mathbf r$ .\end{definition}


\begin{lemma}
[The same geometric input produces verifiable non-zero barriers]\label{hgov:lem:jointobstruction}
For all $t>0$ , $J$ satisfies the Jacobi equation of the given hyperbolic model, and \[
B\omega(t)=0,\qquad
[\omega(t)]\ne0\in H_1(C^-),\qquad
i_*[\omega(t)]=0\in H_1(C^+).
\]\end{lemma}
\begin{certificate}
[Atom-by-atom bridging of geometric to linear barriers]Save $a,\gamma,E,t$, $\lambda(t)$, $A^\pm,B$ and $w=(0,0,1)^{\mathsf T}$. Verify item by item \[
D_t^2J=a^2J,\quad R(J,\dot\gamma)\dot\gamma=-a^2J,\quad
w^{\mathsf T}A^-=0,\quad
w^{\mathsf T}\omega(t)=-\lambda(t).
\] and add the new filling record as $A^+(0,\lambda(t))^{\mathsf T}=\omega(t)$.\end{certificate}
\begin{proof}
The parallel field satisfies $D_tE=0$, and the Jacobi equation and initial value can be obtained by derivation twice. $B\mathbf p=B\mathbf r=0$ proves the linkage on both sides; $A^+i_2=A^-$ proves that $i$ is a chain mapping. Because $t>0$, $a>0$, there is $\lambda(t)>0$. If $\omega(t)=A^-v$, then $w^{\mathsf T}\omega(t)=w^{\mathsf T}A^-v=0$ is inconsistent with $-\lambda(t)<0$. The new column directly gives the boundary witness. So the semantics of \[
\begin{gathered}
t>0\Rightarrow\lambda(t)>0
\Rightarrow w^{\mathsf T}\omega(t)\ne0
\Rightarrow[\omega(t)]_{C^-}\ne0,\\
B A^+=0\land A^+(0,\lambda(t))^{\mathsf T}=\omega(t)
\Rightarrow i_*[\omega(t)]_{C^+}=0 .
\end{gathered}
\] here $\omega(t)$ are determined by the explicit bridge definition and its task input.\end{proof}


\section { Actual action, task quotient and smooth homotopy }

\begin{definition}
[Parameterized task status and unit filling]\label{hgov:def:jointtask}
For the integer $m\ge1$, take the fixed integer \[
t_m=a^{-1}\operatorname{arsinh}(am),\qquad p_m=\gamma(t_m),
\qquad \lambda(t_m)=m .
\] $h\ge0$. On the product substrate $M_m=H_a\times\mathbb R_s\times\mathbb R_u\times\mathbb R_f^h$, the task status field is fixed $p=p_m$ and $u\in\{0,\ldots,m\}$, and $m$, coefficient $s$ and redundancy $f$ are recorded; record this status field as $X_m$. The associated actual change \[
x(s,u)=s\mathbf p+u\mathbf r,\qquad
\omega_m=m\mathbf r,\qquad
\zeta_m(s,u)=\omega_m-x(s,u).
\] task observation includes $(m,p,s,u,\zeta_m)$ , the identity includes $(m,p,Bx)$ , and the target is $G_m\iff \zeta_m=0$ . For each real number label $c$, define the global old action \[
 F_c(p,s,u,f)=(p,s+c,u,f),\qquad
 F_{\mathrm{nf}}(p,s,u,f)=(p,s,u,0).
\]. The predicates of both are always true on $X_m$, and the execution cost is zero. The new unit is filled as part of the action \[
 \begin{aligned}
 \operatorname{Pre}_+(p,s,u,f)&\iff u<m,\\
 F_+(p,s,u,f)&=(p,s,u+1,f),\qquad c_+=1.
 \end{aligned}
\]. The initial state is $(p_m,0,0,f_0)$, where $f_0$ is a fixed input. Each action tag, precondition, and cost has an explicit coordinate-by-coordinate definition as described above.\end{definition}


\begin{proposition}
[Joint tasks of dynamic quotient and smooth homotopy]\label{hgov:prop:jointquotient}
Mapping \[
q(p,s,u,f)=(p,s,u),\qquad i_q(p,s,u)=(p,s,u,0)
\] retains all of the above task observations, identities, targets, action availability, successors and execution costs. Here, the state domain of $q$ is $X_m$ and the quotient domain is $Q_m=q(X_m)$; the same coordinate formula also defines the smooth projection $M_m\to H_a\times\mathbb R_s\times\mathbb R_u$ on the environment substrate. On $M_m$, \[
H_\theta(p,s,u,f)=(p,s,u,(1-\theta)f),\quad 0\le\theta\le1,
\] satisfies $qi_q=\mathrm{id}$ and $qH_\theta=q$, and gives a homotopy equal to $i_qq$. Its smooth singular chain realizes the existence of the prism operator $P$, which satisfies \[
(i_qq)_\#-\mathrm{id}=\partial P+P\partial.
\]\end{proposition}
\begin{certificate}
[quotient and prism interface]Save $O=\bar Oq$, $I=\bar Iq$, $\mathbf1_{G_m}=\mathbf1_{\bar G_m}q$, and \[
 \begin{aligned}
 \bar F_c(p,s,u)&=(p,s+c,u),&
 \bar F_{\mathrm{nf}}&=\mathrm{id}_{Q_m},\\
 \bar F_+(p,s,u)&=(p,s,u+1),&
 \overline{\operatorname{Pre}}_+(p,s,u)&\iff u<m.
 \end{aligned}
\] item by item. The quotient-side predicate of the old action is always true, and the cost is zero; the cost of the new action is one. Save $qF_a=\bar F_aq$, $\operatorname{Pre}_a=\overline{\operatorname{Pre}}_a q$, $c_a=\bar c_aq$ for each label, and save the homotopy endpoints and directed decomposition of the prism.\end{certificate}
\begin{proof}
All task fields are read by $(m,p,s,u)$, so observations, identities and targets are factored. The three types of actions respectively satisfy \[
 \begin{aligned}
 qF_c(p,s,u,f)&=(p,s+c,u)=\bar F_cq(p,s,u,f),\\
 qF_{\mathrm{nf}}(p,s,u,f)&=(p,s,u)=\bar F_{\mathrm{nf}}q(p,s,u,f),\\
 qF_+(p,s,u,f)&=(p,s,u+1)=\bar F_+q(p,s,u,f).
 \end{aligned}
\]. The predicate of the old action is always true, and the new predicate $u<m$ is read from the quotient coordinates; the constant costs of zero and one are also factored respectively. $H_\theta$ fixes all task coordinates and satisfies the listed endpoints. For each smooth simplex, the product of it and the interval is decomposed into a directed prism simplex. The internal faces cancel each other, leaving only the upper and lower ends and side surfaces, and the prism identity \cite{ExtHatcher,S38} is obtained. Therefore \[
\begin{gathered}
\operatorname{CoordinateFactorization}
\land\operatorname{DomainMatch}\\
\land\operatorname{CostFactorization}
\land\operatorname{ActionCommute}
\Rightarrow\operatorname{DynamicQuotient},\\
\operatorname{PrismBoundary}\Rightarrow\operatorname{ChainHomotopy}.
\end{gathered}
\]\end{proof}


\section { Edge-by-edge implementation and exchange of three types of complexes }

\begin{definition}
[Finite execution subgraphs of joint tasks and three types of complexes]\label{hgov:def:jointgraph}
Fixed $m,f_0$, registering actual execution vertices $v_k=(p_m,0,k,f_0)$, $0\le k\le m$, and unit action edges $e_k:v_k\to v_{k+1}$, $0\le k<m$. The chain complex of the execution subgraph is \[
 K_{m,0}=\mathbb R\{v_0,\ldots,v_m\},\quad
 K_{m,1}=\mathbb R\{e_0,\ldots,e_{m-1}\},\quad
 K_{m,2}=0,\qquad
 \partial_me_k=v_{k+1}-v_k.
\]. Its quotient graph uses $\bar v_k=(p_m,0,k)$ and $\bar e_k:\bar v_k\to\bar v_{k+1}$, and defines \[
 q_{m,0}v_k=\bar v_k,\qquad
 q_{m,1}e_k=\bar e_k,\qquad q_{m,2}=0.
\]. The three types of complexes are registered as \[
 C_\bullet^{\mathrm{sm}}(M_m;\mathbb R),\qquad
 K_{m,\bullet},\qquad C^\pm.
\] respectively. They successively save the smooth simplex of the environment substrate, the finite directed record of the actual execution, As well as changing constraints and fixing generators. The task status mapping $v_k\mapsto(p_m,0,k,f_0)$ is registered item by item with each observation value.\end{definition}


\begin{proposition}
[Explicit exchange of smooth paths, task graphs and fix changes]\label{hgov:prop:jointgraph}
For the definition ~ \ref{hgov:def:jointgraph} , there is $q_{m,0}\partial_m=\bar\partial_mq_{m,1}$ . Define smooth implementation $\mathcal R_m:K_{m,\bullet}\to
C_\bullet^{\mathrm{sm}}(M_m;\mathbb R)$: \[
 \mathcal R_m(v_k)=[(p_m,0,k,f_0)],\qquad
 \mathcal R_m(e_k)=\sigma_k,\quad
 \sigma_k(\tau)=(p_m,0,k+\tau,f_0),\quad 0\le\tau\le1.
\] then $\partial_{\mathrm{sm}}\mathcal R_m=\mathcal R_m\partial_m$. The quotient side takes the $q$ projection of the same path and gets $q_\#\mathcal R_m=\bar{\mathcal R}_mq_m$. Further define \[
 \rho_m(v_k)=k\mathbf r,\qquad
 \chi_m(e_k)=(0,1)^{\mathsf T},\qquad
 \xi_m(v_k)=(m-k)\mathbf r.
\] These mappings satisfy the bridge type \[
 \rho_m\partial_m=A^+\chi_m,\qquad
 B\rho_m=0,\qquad
 \xi_m\partial_m=-A^+\chi_m.
\] Treat $x(s,u)=s\mathbf p+u\mathbf r$ by coordinate as three smooth observations on the environment substrate, then \[
 \int_{\sigma_k}dx=\mathbf r=A^+\chi_m(e_k),\qquad
 \int_{\sigma_k}d\zeta_m=-\mathbf r,
\] and $\omega(t_m)=\rho_m(v_m)=\xi_m(v_0)$ .\end{proposition}
\begin{certificate}
[Edge-by-edge realization, boundaries and integrals]Save all vertex and edge tables, the vertex binding of $q_{m,0}$, the path $\sigma_k$ and its two endpoints, and save \[
 \partial_{\mathrm{sm}}\sigma_k
 =[(p_m,0,k+1,f_0)]-[(p_m,0,k,f_0)].
\] edge by edge. Change the bridge to save $A^+(0,1)^{\mathsf T}=\mathbf r$, $B\mathbf r=0$, $dx=\mathbf p\,ds+\mathbf r\,du$, and $\sigma_k^*ds=0$, $\sigma_k^*du=d\tau$. Geometry input save $\lambda(t_m)=m$.\end{certificate}
\begin{proof}
For each edge generator, the boundary of the smooth simplex \[
 q_{m,0}\partial_me_k
 =\bar v_{k+1}-\bar v_k
 =\bar\partial_mq_{m,1}e_k.
\] is the end point minus the starting point, so $\partial_{\mathrm{sm}}\mathcal R_me_k
=\mathcal R_m\partial_me_k$ is given edge by edge; the boundaries on both sides of the zero degree are zero, and the other degrees are empty, and linear continuation results in a chain map. Projection $q$ acts on each parameter value of $\sigma_k$, directly obtaining the quotient side path, so $q_\#\mathcal R_m=\bar{\mathcal R}_mq_m$. And \[
 \begin{aligned}
 \rho_m\partial_me_k
 &=(k+1)\mathbf r-k\mathbf r
 =\mathbf r=A^+\chi_m(e_k),\\
 \xi_m\partial_me_k
 &=(m-k-1)\mathbf r-(m-k)\mathbf r
 =-\mathbf r,\\
 \int_{\sigma_k}dx
 &=\int_0^1\bigl(\mathbf p\,0+\mathbf r\,d\tau\bigr)
 =\mathbf r.
 \end{aligned}
\]. Since $\zeta_m=\omega_m-x$ and $\omega_m$ are fixed in this execution, the second integral formula is immediately established. $B\mathbf r=0$ gives $B\rho_m=0$; a geometry read of $\lambda(t_m)=m$ gives the last endpoint bridged. The entire predicate chain is \[
 \begin{gathered}
 \operatorname{EdgeTable}
 \Longrightarrow q_{m,0}\partial_m=\bar\partial_mq_{m,1},\\
 \operatorname{SmoothEndpoints}(\sigma_k)
 \Longrightarrow
 \partial_{\mathrm{sm}}\mathcal R_m=\mathcal R_m\partial_m,\\
 \sigma_k^*dx=\mathbf r\,d\tau
 \Longrightarrow
 \int_{\sigma_k}dx=\rho_m\partial_me_k=A^+\chi_m(e_k).
 \end{gathered}
\] and the smooth edge internal parameters are located in the environment substrate $M_m$; the execution accepts stepwise verification on integer vertices and registered unit actions.\end{proof}


\section { finite hit, least direction and complete cost }

\begin{theorem}
[Strictly reachable expansion, finite hit and action lower bounds for joint tasks]\label{hgov:thm:joint}
For each $m\ge1$, the old action family cannot reach $G_m$ from the initial state; the new action word is composed of $m$ sub-unit padding and accurately reaches $G_m$. The word maintains task identity, dynamic quotient, and geometric input throughout, and each hit word is padded with at least $m$ subunits. Along the construction path, the potential function \[
\Phi(s,u)=\tfrac12\|\zeta_m(s,u)\|^2
\] satisfies \[
\Phi(0,k)=(m-k)^2,\qquad
\Phi(0,k)-\Phi(0,k+1)=2(m-k)-1>0
\] and holds for $0\le k<m$.\end{theorem}
\begin{certificate}
[finite induction locus for all parameters]Retain the parameter $m$, analytic geometric input $t_m,p_m$, the actual action sequence, and row $k$:
\[
(s_k,u_k,x_k,\zeta_k,\Phi_k,r_k)
=(0,k,k\mathbf r,(m-k)\mathbf r,(m-k)^2,m-k).
\].
Each row records the action precondition, addition in three coordinates, $Bx_k=0$, the task-quotient image, and unit cost, and binds $v_k,e_k$ in the finite execution subgraph to the path realization of Proposition ~ \ref{hgov:prop:jointgraph}. The final row records $\zeta_m=0$ and total execution cost $m$.\end{certificate}
\begin{proof}
The old action remains $u=0$. The third change coordinate of any old successor is zero, and the third coordinate of the target $m\mathbf r$ is $-m$ , so the old family is unreachable. The new family takes $s_0=u_0=0$ as the base example. If $k<m$, then $u_k=k<m$, the unit action is legal and $u_{k+1}=k+1$ is generated, and the other fields remain. All listed trajectories are mathematically establish by induction, and the goal of step $m$ is established. The identity preservation is given by $B\mathbf p=B\mathbf r=0$ with the geometric coordinates unchanged. Since $\mathbf p,\mathbf r$ is linearly independent, the target is equivalent to $s=0,u=m$. The old action does not change $u$, each unit fill only increases by one, so any hit requires at least $m$ units to fill. Finally, for $\|\mathbf r\|^2=2$, the potential function can be directly expanded to obtain the amount of decline. The atomic predicate causal chain is \[
\begin{gathered}
m\ge1\land u_0=0\land
\operatorname{OldPreserves}(u)
\Rightarrow u_{\rm old}=0<m
\Rightarrow\neg G_m,\\
u_k=k<m
\Rightarrow\operatorname{Enabled}(a_k)
\Rightarrow u_{k+1}=k+1
\Rightarrow r_{k+1}=r_k-1,\\
r_0=m\land r_m=0
\Rightarrow x_m=m\mathbf r
\Rightarrow\zeta_m=0
\Rightarrow G_m .
\end{gathered}
\]. The geometric bridge, chain mapping, homotopy quotient and actual repair are established under the same parameters and task specification.\end{proof}


\begin{corollary}
[Minimality and complete cost of a new independent direction]\label{hgov:cor:jointcost}
$\dim H_1(C^-)=1$ and $H_1(C^+)=0$ in this task family. Make $\operatorname{span}\{\mathbf r\}$ all become boundaries, add at least one independent direction, and define \ref{hgov:def:joint} to reach the lower bound. The complete delivery cost is registered as \[
\begin{aligned}
C_{\rm total}={}&C_{\rm geometry}+C_{\rm input}+C_{\rm search}
+C_{\rm execution}\\
&+C_{\rm verification}+C_{\rm encoding}+C_{\rm readback}
\end{aligned}
\]. The lower bound of the actual execution fee is $m$, and the above structure reaches this lower bound.\end{corollary}
\begin{certificate}
[Rank and itemized costs]Record $\operatorname{rank}B=1$, $\operatorname{rank}A^-=1$, $\operatorname{rank}A^+=2$; record the measurement unit and actual number of calls of each cost primitive. Symbolic preservation of analytical expressions and numerical evaluation to specified precision use separate cost records.\end{certificate}
\begin{proof}
$\dim\ker B=2$, so the homology dimensions on both sides are $2-1$ and $2-2$ respectively. The target class is non-zero in the old homology, so at least one new quotient direction is required; column $\mathbf r$ gives the equal sign construction. Any hit must have $m$ unit actions, each costing one, and other costs are non-negative, so the execution cost is at least $m$; the constructed word only has these $m$ actions, so it reaches the lower bound. The whole process costs are accumulated from the required stages according to the same measurement protocol. \begin{gather*}
\operatorname{RankWitness}
\Rightarrow\operatorname{OneMissingDirection}
\Rightarrow\operatorname{MinimalExtension},\\
\operatorname{UnitProgressBound}\Rightarrow C_{\rm execution}\ge m .\qedhere
\end{gather*}\end{proof}


\begin{remark}
[The structure maintained by the union construct]This family maintains the metric and negative curvature, task identity and geometric parameters of the hyperbolic plane throughout the process; the progress coordinate changes the actual residual, and the certification expansion changes the available repair direction. The task quotient eliminates redundant fibers, and the homology expansion makes the specified target class a boundary. The above-mentioned maintained relationships are verified by coordinate formula, chain mapping, prism homotopy and execution induction respectively.\end{remark}


\section { Derived structure fields with a conservative extension of the same task }

\begin{lemma} [Lifting execution to the graph of structural fields]
\label{exp:joint:derived-fields}
Let the state domain of the task system be $X$, with a fully defined structural assignment $H:X\to Z$. Set
\[
 \widehat X=\{(x,H(x)):x\in X\},\qquad
 j_H(x)=(x,H(x)),\qquad p_H(x,H(x))=x.
\]
and lift each action to
$\widehat F_a(x,H(x))=(F_a x,H(F_a x))$.
Preconditions, observations, identity, goals, and costs are all read through $p_H$. Then $j_H,p_H$ are inverse maps, and all finite admissible executions, their acceptance results, and their costs correspond exactly.
\end{lemma}
\begin{certificate} [Recomputation of derived values and fieldwise squares]
Record the source, target, actual evaluation, and version of $H$. At every step retain the recomputation record for $H(F_ax)$, agreement of preconditions, and
$p_H\widehat F_a=F_ap_H$, $\widehat F_aj_H=j_HF_a$.
\end{certificate}
\begin{proof}
By definition, $p_Hj_H=\mathrm{id}_X$ and $j_Hp_H=\mathrm{id}_{\widehat X}$. The action commutation relation follows by substitution of coordinates; all preconditions and task fields are read by projection and therefore agree value by value. Induct on the length of the action word. The empty word is the identity on both sides. Appending one commuting step to a verified prefix gives corresponding successors and equal partial cost sums. Thus the graph of structural fields preserves all registered executions of the original task.
\end{proof}

\section { homology fix common residuals with higher-order compatibility }

\begin{definition}
[Derived higher-order operations in union instances]\label{exp:joint:algebra-lift}
This section takes $\mathbb K=\mathbb R$. Let $W$ be the original graded space of the theorem ~ \ref{hgov:linfty:model}, and use all the symbols that define ~ \ref{exp:higher:parameter-family}. For the joint task state $z=(p_m,s,u,f)$ , register the derived structure \[
 H_m(z)=\ell_\bullet^{m,mu},\qquad
 \ell_1(z)=\ell_1,\quad
 \ell_2(z)=m\ell_2,\quad
 \ell_3(z)=mu\ell_3.
\] and the remaining higher-order operations are zero. Binary coefficients are fixed by geometry reading $\lambda(t_m)=m$ and ternary coefficients are read by actual progress $u$. Let $\eta_3:C_1=\mathbb R^3\to\mathbb R$ be the third physical change coordinate, define \[
 \mathcal B_m:C_1\longrightarrow
 \operatorname{Hom}^{0}(\Lambda^3W,W),\qquad
 \mathcal B_m(v)=-m\eta_3(v)\Omega.
\] where $\Omega$ is the non-zero ternary composite operator given in Chapter ~ \ref{exp:ch:higher} ~, whose source target and physical change vector are located $C_1$ is registered separately.\end{definition}


\begin{theorem} [Linear repair and higher-order closure along the same admissible path]
\label{exp:joint:simultaneous-repair}
Under Definition ~ \ref{exp:joint:algebra-lift}, every state satisfies
\[
 \mathcal R_3(z)=m(m-u)\Omega
 =\mathcal B_m\zeta_m(s,u),\qquad
 \mathcal R_n(z)=0\quad(n\ne3).
\]
and the task goal has the necessary and sufficient criterion
\[
 \zeta_m(s,u)=0
 \quad\Longleftrightarrow\quad
 s=0\ \land\ \operatorname{HigherClosed}(W,H_m(z)).
\].
Along the execution specified in Theorem ~ \ref{hgov:thm:joint}, consisting of $m$ steps, the geometric input, binary operation, and task identity are preserved, while the linear residual and the full higher-order compatibility residual both vanish at step $m$.
\end{theorem}
\begin{certificate} [A linear bridge between the two residuals]
Register step by step
\[
 u_k=k,\quad s_k=0,\quad
 \zeta_k=(m-k)\mathbf r,\quad
 \beta_k=mk,\quad
 \mathcal R_{3,k}=m(m-k)\Omega.
\]
and retain $\eta_3(\mathbf p)=0$, $\eta_3(\mathbf r)=-1$, $\Omega(e_1,e_2,e_3)=e_3$, and the multiplication record showing that every unit action increases $\beta$ by $m$. All operation fields are lifted according to Lemma ~ \ref{exp:joint:derived-fields}.
\end{certificate}
\begin{proof}
The parameter-compatibility theorem ~ \ref{exp:higher:scaled-closure} gives
\[
 \mathcal R_3(z)=(m^2-mu)\Omega=m(m-u)\Omega,
\]
and the remaining arity relations vanish. On the other hand,
\[
 \eta_3\zeta_m(s,u)
 =\eta_3\bigl((m-u)\mathbf r-s\mathbf p\bigr)=-(m-u),
\],
which gives the stated bridging equality. The vectors $\mathbf p,\mathbf r$ are linearly independent, so $\zeta_m=0$ if and only if $s=0,u=m$. By $m\ge1$ and $\Omega\ne0$, higher-order closure holds if and only if $u=m$. Combining these two equivalences yields the criterion for the goal.

On the construction path, the base case $u_0=0,s_0=0$ gives $\zeta_0=m\mathbf r$ and $\mathcal R_{3,0}=m^2\Omega$. If $k<m$, the unit action makes $u_{k+1}=k+1$, while $m,p_m,s,f$ remains, so \[
 \zeta_{k+1}=\zeta_k-\mathbf r,\qquad
 \mathcal R_{3,k+1}=\mathcal R_{3,k}-m\Omega.
\] can be summarized to get all the prefixes, and both residuals are zero at $k=m$. The coefficient of $\ell_2^{m,mu}=m\ell_2$ is fixed throughout, and the identity is maintained by $B\mathbf p=B\mathbf r=0$. So \[
 \begin{gathered}
 \lambda(t_m)=m\land\operatorname{UnitExecution}(k)
 \Rightarrow \zeta_k=(m-k)\mathbf r,\\
 \mathcal R_{3,k}=\mathcal B_m\zeta_k
 \Rightarrow
 \bigl(k=m\Longleftrightarrow\zeta_k=0
 \Longleftrightarrow\operatorname{HigherClosed}(H_m(z_k))\bigr).
 \end{gathered}
\]\end{proof}


\begin{corollary} [Exact readout of the total square of the suspended coderivation]
\label{exp:joint:coderivation-square}
Using the suspension of Lemma ~ \ref{hgov:lem-suspension}, namely $s:W\to W[1]$, convert $H_m(z)$ into the degree-one coderivation $Q_{m,u}$. Define the symmetric ternary map of degree two
\[
 \Omega_s(sx\odot sy\odot sz)
 =(-1)^{2|x|+|y|}s\Omega(x,y,z),
\]
and let $\widehat\Omega_s$ be its coderivation extension. Then
\[
 Q_{m,u}^{2}=m(m-u)\widehat\Omega_s,
 \qquad Q_{m,u}^{2}=0\Longleftrightarrow u=m.
\].
\end{corollary}
\begin{certificate} [Aritywise projections and the unique coderivation extension]
Retain the suspension signs, the ternary projection, certificates that every other arity projection vanishes, and the nonzero readout of $\Omega_s$ on tuples of suspended basis elements.
\end{certificate}
\begin{proof}
The suspension interface takes $Q_{m,u}^2$, projects it to cogenerators at arity $n$, and converts that projection into the signed suspension of the corresponding $\mathcal R_n$. Theorem ~ \ref{exp:joint:simultaneous-repair} makes all nonternary projections zero, while the ternary projection equals $m(m-u)\Omega_s$. The square of a degree-one coderivation is a coderivation, uniquely determined on the reduced symmetric coalgebra by its projection to cogenerators. Hence the total-square equality follows. $\Omega_s\ne0$ and $m\ne0$ give the final necessary and sufficient condition.
\end{proof}

\section { The complete value of parameter two certificate }

\begin{proposition}
[All acceptance records for two-step instances]\label{exp:joint:m-two}
Take $m=2$ and the geometric input is $t_2=a^{-1}\operatorname{arsinh}(2a)$, then $\lambda(t_2)=2$. From $s=u=0$, two consecutive unit actions are executed to obtain the following table; where $g_k$ is defined as $\mathcal R_{3,k}=g_k\Omega$.\begin{center}
\begin{tabular}
{ccccccc}\toprule
$k$&
$u_k$&
$x_k$&
$\zeta_k$&
$\Phi_k$&
$\beta_k$&
$g_k$\\
\midrule
0&
0&
$(0,0,0)$&
$(2,0,-2)$&
4&
0&
4\\
1&
1&
$(1,0,-1)$&
$(1,0,-1)$&
1&
2&
2\\
2&
2&
$(2,0,-2)$&
$(0,0,0)$&
0&
4&
0\\
\bottomrule
\end{tabular}
\end{center}
Each bank meets $Bx_k=0$, and the total cost is $k$. In the second step, the original task goal and complete higher-order compatibility are achieved at the same time, and the actual number of unit actions reaches the lower bound of two.\end{proposition}
\begin{certificate}
[Coordinate-wise addition, potential difference, and scalar multiplication]Save input $m=2$, vector $\mathbf r=(1,0,-1)^{\mathsf T}$, antecedents $0<2$ and $1<2$, two additions, potential difference $4-1=3$ and $1-0=1$, coefficient $\beta_k=2k$ and $g_k=4-2k$.\end{certificate}
\begin{proof}
Use $x_k=k\mathbf r$, $\zeta_k=(2-k)\mathbf r$, $\|\mathbf r\|^2=2$ row by row to get the vectors and potential values in the table. The parameter formulas give $\beta_k=2k$ and $g_k=2(2-k)$. The old action remains $u=0$ and cannot reach $u=2$; each new action only increases by one, so it takes at least two hits to reach this limit in the table.\end{proof}


\begin{remark}
[Multiple layers of reads on the same task]The geometry read $\lambda(t_m)$, the boundary preimage of the inpainted complex, the coderivation square of the higher-order operation, the execution ranking and the normal form encoding are each associated to the same registration state by explicit mappings. The target acceptance reads the complete predicate $s=0\land u=m$, and each intermediate certificate participates in joint verification according to its source target, domain and version.\end{remark}


\chapter{Task Normal Forms, Finite Sufficient Quotients, and Capability Reuse}
\label{exp:ch:normalform}


\label{hgov:sec:normalform}

This chapter uses fixed tasks, admissible actions, observations, and cost primitives to gradually construct a representation that can preserve all specified future actions. First, the complete identification of observation readback and dynamic equivalence classes is distinguished, and then the finite table algorithm, extractable distinguishing words and all calculation examples are given. Finally, it is explained how to encapsulate the certified strategy on quotient into capabilities, and count the one-time construction, actual call and re-authentication into costs respectively. This development follows the evolutionary saturation and formal structure of the open source manuscript \cite{S19,S38,GFrameworkSynthesis}.

\section { Observe the uniqueness of decoding and normal form }\label{exp:gov:normalform-section}


\begin{definition}
[task normal form and executable capability package]\label{hgov:def:normalform}
Fixed task fields $X$, prescribed observations $O:X\to Y$ and admissible actions. When the task encoding $N:X\to Z$ together with the decoder $D:N(X)\to Y$ satisfies $DN=O$, it is called a readback representation of the observation; if $N$ also accurately identifies the declared dynamic task equivalence class, it is called the canonical type of the task domain. The capability package is registered as \[
\mathcal K=(N,D,\gamma,\mathsf{Dom},\mathsf{Fail},\mathcal C,
\mathsf{Cost},\mathsf{Version}),
\], where $\gamma$ is the verified action path, $\mathsf{Dom}$ is the applicable domain, $\mathsf{Fail}$ is the certified invalid domain, and the remaining components store the certification, cost and version. normal form and path encoding continuation \cite{S19,S38,GFrameworkSynthesis}.\end{definition}


\begin{theorem}
[Fiber criterion for normal form decoding]\label{hgov:thm:nf}
There is a unique $D:N(X)\to Y$ that satisfies $DN=O$, if and only if \[
\forall x,y\in X,\quad N(x)=N(y)\Rightarrow O(x)=O(y).
\] requires readback of all original objects, then the condition is $N$ injective.\end{theorem}
\begin{certificate}
[Coding fibers and specific decoding]Save the code, observations, proof of consistency for each relevant fiber, and program or finite table defining $D(N(x))=O(x)$. Requires bit-by-bit recovery while saving the actual encoded version.\end{certificate}
\begin{proof}
If $DN=O$, then the two inputs with the same encoding will get the same observation through $D$. In reverse, it is known from fiber consistency that $D(N(x))=O(x)$ has nothing to do with the representative element; every value in $N(X)$ has a representative, so $D$ exists and is unique. When taking $O=\mathrm{id}_X$, the condition is equivalent to the same object with the same encoding. Therefore \[
\operatorname{FiberConstant}(O,N)
\Rightarrow\operatorname{WellDefined}(D)
\Rightarrow DN=O,
\quad
DN=O\Rightarrow\operatorname{FiberConstant}(O,N).
\]\end{proof}


\begin{theorem}
[normal form unique to unique bijective of the same task equivalence relationship]\label{exp:gov:normalform-unique}
Assume $E$ is the equivalent relationship on $X$, and the two codes $N_i:X\to Z_i$ satisfy \[
 N_i(x)=N_i(y)\quad\Longleftrightarrow\quad xEy
 \qquad(i=1,2).
\] such that $Q_i=N_i(X)$. The only bijective $\varphi:Q_1\to Q_2$ exists makes $\varphi N_1=N_2$ . If $O=D_iN_i$ , then \[
 D_2\varphi=D_1.
\] If the domain of a certain partial action $F$ is the union of $E$ classes, and $xEy$ and $x,y\in\operatorname{Dom}(F)$ contain $F(x)EF(y)$ , then the induced partial actions on the two normal types satisfy \[
 \varphi\bar F_1=\bar F_2\varphi
\] and have a domain corresponding to $\varphi$. Also if $\nu:X\to X$ satisfies $\nu(x)Ex$ and is constant on every $E$ class, then \[
 \nu^2=\nu,\qquad
 \nu(x)=\nu(y)\Longleftrightarrow xEy.
\] therefore represents an idempotent normalization given by the meta-choice while maintaining the same task quotient between different legal choices.\end{theorem}
\begin{certificate}
[Two encoded bidirectional tables conjugated with successor]Register the precise fiber criteria for $E$ and $N_1,N_2$, as well as $\varphi(N_1x)=N_2x$ and $\psi(N_2x)=N_1x$. Save the observation decoding formula, the saturation of part of the action domain, and the successor compatibility formula. When using representative elements, register that the selected object belongs to this category on a category-by-category basis, and save the version of the representative element selection table.\end{certificate}
\begin{proof}
If $N_1x=N_1y$ , then $xEy$ , and thus $N_2x=N_2y$ ; therefore $\varphi$ is well-defined. In the same way, $\psi$ is well-defined, and the two codes \[
 \psi\varphi(N_1x)=N_1x,\qquad
 \varphi\psi(N_2x)=N_2x.
\] are surjective on its image, so these two equations prove to be mutually inverse; the value of any mapping that satisfies $\varphi N_1=N_2$ on each image value has been determined, and is therefore unique. For $z=N_1x$ , the domain saturation of \[
 D_2\varphi(z)=D_2N_2x=O(x)=D_1z.
\] makes $\bar F_i(N_ix)=N_iF(x)$ a well-defined partial function. So for its legal input, \[
 \varphi\bar F_1(N_1x)=N_2F(x)
 =\bar F_2(N_2x)=\bar F_2\varphi(N_1x).
\] and finally, $\nu(x)Ex$ with a constant value on the class gives $\nu(\nu(x))=\nu(x)$ . If $\nu(x)=\nu(y)$, use $xE\nu(x)=\nu(y)Ey$ to get $xEy$; the reverse is a class constant. Therefore \[
\begin{gathered}
 \operatorname{SameEquivalenceFibers}(N_1,N_2)
 \Rightarrow\operatorname{UniqueBijection}(\varphi),\\
 \operatorname{Readback}\land\operatorname{DomainSaturated}
 \land\operatorname{SuccessorCompatible}
 \Rightarrow\operatorname{DecoderAndActionConjugacy},\\
 \operatorname{ClassRepresentative}(\nu)\land\operatorname{ClassConstant}(\nu)
 \Rightarrow\nu^2=\nu\land\operatorname{CompleteClassCode}(\nu).
\end{gathered}
\]\end{proof}


\section { Future observation equivalence and finite partition refinement }\label{exp:gov:refinement-section}


\begin{definition}
[successor saturation of finite task testing]\label{hgov:def:saturation}
Fixed finite state set $X$ with explicit enumerations, finite label set $L$ , and partial successor $F_\ell:\operatorname{Dom}(F_\ell)\to X$ given by explicit finite tables. Each specified observation has an evaluable finite table, and the equivalence of the table items is decidable; each action cost has a finite table that can be accurately summed, and the equivalence is decidable. Nonnegative rational costs give a direct realization of this input condition. Add undefined mark $\bot$ , let $X_\bot=X\sqcup\{\bot\}$ , and define \[
 \widehat F_\ell(x)=
 \begin{cases}
 F_\ell(x),&x\in\operatorname{Dom}(F_\ell),\\
 \bot,&x\notin\operatorname{Dom}(F_\ell).
 \end{cases}
 \qquad \widehat F_\ell(\bot)=\bot.
\] . Each original observation $f:X\to Y_f$ is expanded to $\widehat f:X_\bot\to Y_f\sqcup\{*_f\}$ , where $\widehat f(\bot)=*_f$ , and the remaining values are read according to the original table. Target and availability take Boolean extensions: $\mathbf1_G(\bot)=0$, $e_\ell=\mathbf1_{\operatorname{Dom}(F_\ell)}$, $e_\ell(\bot)=0$. The full list of action charges with type \[
 \widehat c_\ell(x)=
 \begin{cases}
 c_\ell(x),&x\in\operatorname{Dom}(F_\ell),\\
 *_\ell^{\mathrm{cost}},&x\notin\operatorname{Dom}(F_\ell)
 \ \text{or}\ x=\bot.
 \end{cases}
\] The initial finite test family $\mathcal P_0$ contains all specified observations, task identities, targets, each $e_\ell$ , each $\widehat c_\ell$ , and $\mathbf1_{\{\bot\}}$ . All members are defined in $X_\bot$. Recursively define \[
\mathcal P_{n+1}
=\mathcal P_n\cup\{f\circ\widehat F_\ell:
f\in\mathcal P_n,\ell\in L\},
\quad
x\sim_ny\iff\forall f\in\mathcal P_n,\ f(x)=f(y).
\] Let $b_n(x)$ be the current block number, and the actual refinement is grouped by the finite signature \[
 s_{n+1}(x)=
 \left(b_n(x),\bigl(b_n(\widehat F_\ell x)\bigr)_{\ell\in L}\right)
\]. The initial numbering is determined by the total table value of $\mathcal P_0$, and the block numbering adopts fixed rules under the given enumeration.\end{definition}


\begin{theorem}
[The construction, sufficiency and roughest properties of finite stable quotient]\label{hgov:thm:saturation}
If the initial partition has $b_0$ blocks, the above structure will be stable after at most $|X_\bot|-b_0$ strict refinements. The stable quotient $q_*:X_\bot\to Q_*$ preserves the initial observations after all finite words, and is the coarsest quotient that preserves these observations and the label successor. Its action domain, cost per step, and cumulative cost of any finite legal word can be read back from quotient. Under a given enumeration, the minimum number of each stable block is taken as the code, and the computable normal form of the dynamic task equivalence relationship is obtained.\end{theorem}
\begin{certificate}
[Partition refinement and word length induction]Save the complete successor table, extended observation table, marked expense table and judgment procedures for table items. Each split saves the same block's different value pairs and their observations, labels or action words; the final state saves each block's observation constants, each label successor same block, legal actions and cost constants. Save the initial number of blocks, the number of blocks increased each time, and the total number of states. The accumulated cost certificate stores the actual cost and partial sum by prefix; the normal form stores the minimum number and representative element reading table of the stable block.\end{certificate}
\begin{proof}
Finite transition tables can be evaluated, observation tables can be evaluated, etc., so the initial grouping and each finite signature comparison are completed in finite steps. By definition, \[
 x\sim_{n+1}y
 \Longleftrightarrow
 x\sim_ny\ \land\
 \bigwedge_{\ell\in L}
 \bigl(\widehat F_\ell x\sim_n\widehat F_\ell y\bigr),
\] so the signature algorithm exactly implements the defined test refinement. Strict refinement adds at least one non-empty block, and the number of blocks is at most $|X_\bot|$, so the number of strict refinements is bounded. If the partitions in two adjacent rounds are the same, each label successor input in the same block is still in the same block, otherwise a certain $f\circ\widehat F_\ell$ will be split. Summed up by word length, the successors of all finite words remain in the same block. The initial observation is constant over the block, so all finite successor observations fall to a stable quotient. Among them, the constancy of $e_\ell$ gives the domain matching, and the constancy of $\widehat c_\ell$ gives the cost function $\bar c_\ell$ on the legal domain, satisfying $c_\ell=\bar c_\ell q_*$. If $w=(\ell_1,\ldots,\ell_N)$ is legal in $x_0$, let $x_k=F_{\ell_k}(x_{k-1})$, then the word length is summarized to give $\bar x_k=q_*x_k$, and then \[
 \begin{aligned}
 C_0&=\bar C_0=0,\\
 C_{k+1}
 &=C_k+c_{\ell_{k+1}}(x_k)
 =\bar C_k+\bar c_{\ell_{k+1}}(\bar x_k)
 =\bar C_{k+1}.
 \end{aligned}
\]. This also proves that all cost prefixes and the final total cost are maintained. In the reverse direction, any equivalence relationship between the initial observation and the label successor is maintained, and each $f\circ\widehat F_w$ is inductively maintained by word length, thus being finer than the stable relationship. Therefore the \[
\begin{gathered}
\operatorname{FiniteTables}\land\operatorname{DecidableEntries}
\Rightarrow\operatorname{EffectiveRefinement},\\
\operatorname{EffectiveRefinement}\land\operatorname{BoundedSplits}
\Rightarrow\operatorname{StablePartition},\\
\operatorname{StablePartition}\Rightarrow\operatorname{SuccessorCongruence}
\Rightarrow\operatorname{AllWordObservations},\\
\operatorname{AnySufficientCongruence}
\Rightarrow\operatorname{Refines}(q_*)
\Rightarrow\operatorname{Coarsest}(q_*).
\end{gathered}
\] undefined flag ensures that the action domain is preserved with the same rules. All specified tests in the stable block take the same value, so decoding can be obtained by reading from the representative element corresponding to the smallest number. The minimum number of two states is the same if and only if they belong to the same stable block, thus obtaining the normal form.\end{proof}


\begin{theorem}
[Finite deep recognition, distinguishing words and stabilizing the universal properties of quotient]\label{exp:gov:future-equivalence}
Follow the definition ~ \ref{hgov:def:saturation}. The word $w=(\ell_1,\ldots,\ell_k)\in L^*$ stipulates that \[
 \widehat F_w=\widehat F_{\ell_k}\cdots\widehat F_{\ell_1},
 \qquad \widehat F_{\varepsilon}=\mathrm{id}_{X_\bot}.
\] defines all future observations equivalent to \[
 x\equiv_\infty y
 \quad\Longleftrightarrow\quad
 \forall w\in L^*\ \forall f\in\mathcal P_0,
 \quad f(\widehat F_wx)=f(\widehat F_wy).
\], then there is an accurate identification formula for each $n\ge0$ \[
 x\sim_ny
 \quad\Longleftrightarrow\quad
 \forall |w|\le n\ \forall f\in\mathcal P_0,
 \quad f(\widehat F_wx)=f(\widehat F_wy).
\]. Taking $r=\min\{n\ge0:\sim_{n+1}=\sim_n\}$, then $\sim_r=\equiv_\infty$ . For any two non-equivalent states there exists a distinguishing pair $(f,w)$ that can be extracted from the split record, and $|w|\le r\le |X_\bot|-b_0$ .

If a surjection $p:X_\bot\to R$ satisfies
$f=\widetilde f p$ for every $f\in\mathcal P_0$, and
$p\widehat F_\ell=\widetilde F_\ell p$,
then there exists a unique surjection
$u:R\to Q_*$ such that
\[
 q_*=up,\qquad
 u\widetilde F_\ell=\bar F_\ell u.
\].
In particular, every dynamic representation that exactly preserves all future observations factors through $Q_*$.
\end{theorem}
\begin{certificate} [The splitting tree and the unique factorization table]
An initial split records different values of an observation $f\in\mathcal P_0$, with corresponding word $\varepsilon$. A successor-induced split records the label $\ell$ and an existing distinction record for a pair of successors; during backtracking, prepend $\ell$ to that distinguishing word. Record the first stable round, the number of blocks at every round, and the tables showing commutation of observations and successors for other candidate representations. Verify the factor map using the finite table $u(p(x))=q_*(x)$.
\end{certificate}
\begin{proof}
Induct on depth $n$. When $n=0$, only the empty word occurs, so the conclusion is precisely the initial observation partition. Assume the result at depth $n$. Two states belong to the same block in the next round if and only if they currently share a block and every $\ell$-successor shares a block at round $n$. The induction hypothesis converts the latter condition into, for all $|w|\le n$,
\[
 f(\widehat F_w\widehat F_\ell x)
 =f(\widehat F_w\widehat F_\ell y).
\].
These are exactly the tests with first label $\ell$ and total length at most $n+1$. Adding the empty word covers every required word and completes the induction.

The stable block successor is still the same block. The zero-step identity and one-step continuation for word length show that $\sim_r$ implies the equality of all future observations; the reverse is obtained by testing that the length does not exceed $r$. If $x\not\sim_ry$ , the above finite depth equivalent provides a differentiator whose length does not exceed $r$. The backtracking of the split record adds at most one first label each round, so the executable structure of the word is also given. The bounds on the number of strict splits are provided by the theorem ~ \ref{hgov:thm:saturation}; there is a strict split at each round before the first invariance, so $r\le |X_\bot|-b_0$ .

Finally, if $p(x)=p(y)$, the word length induction is $p\widehat F_wx=p\widehat F_wy$: the empty word is the antecedent, and the induction step uses successor exchange. Observing factorization then gives $x\equiv_\infty y$ , thus $u(p(x))=q_*(x)$ is independent of the representative element. $p$ surjective gives uniqueness, $q_*$ surjective gives $u$ surjective; substitute any $p(x)$ to get \[
 u\widetilde F_\ell(p(x))
 =u(p\widehat F_\ell x)
 =q_*\widehat F_\ell x
 =\bar F_\ell q_*x
 =\bar F_\ell u(p(x)).
\]. In summary, \[
\begin{gathered}
 \operatorname{Depth}(n)\land\operatorname{OneStepRefinement}
 \Rightarrow\operatorname{Depth}(n+1),\\
 \operatorname{StableAt}(r)
 \Rightarrow (\sim_r=\equiv_\infty)
 \Rightarrow\operatorname{FiniteSeparatingWitness},\\
 \operatorname{CandidateCongruence}(p)
 \Rightarrow\ker p\subseteq\ker q_*
 \Rightarrow\operatorname{UniqueFactor}(u).
\end{gathered}
\] here $\ker p$ represents the iso-image equivalence relationship rather than the kernel of the linear mapping.\end{proof}


\section { All refinement and readback calculations of the five-state model }\label{exp:gov:finite-example-section}


\begin{proposition}
[Five-state example of two rounds of refinement and preservation of non-trivial fibers]\label{exp:gov:five-state}
Take $X=\{a,b,c,d,e\}$, unique action label $\lambda$, and given the following table. The action is defined in all five real states, the legal step cost is $1$; the task identity is a constant, the target set is $G=\{e\}$, and the specified observation is $U=\mathbf1_G$.\begin{center}
\begin{tabular}
{ccccc}\toprule
Status $x$&
$F_\lambda(x)$&
$U(x)$&
$e_\lambda(x)$&
$c_\lambda(x)$\\
\midrule
$a$&
$c$&
$0$&
$1$&
$1$\\
$b$&
$d$&
$0$&
$1$&
$1$\\
$c$&
$e$&
$0$&
$1$&
$1$\\
$d$&
$d$&
$0$&
$1$&
$1$\\
$e$&
$e$&
$1$&
$1$&
$1$\\
\bottomrule
\end{tabular}
\end{center}
After adding the independent $\bot$ as defined, the whole process of partitioning is\begin{center}
\begin{tabular}
{cl}\toprule
rounds&
All chunks\\
\midrule
$0$&
$\{a,b,c,d\},\ \{e\},\ \{\bot\}$\\
$1$&
$\{a,b,d\},\ \{c\},\ \{e\},\ \{\bot\}$\\
$2$&
$\{a\},\ \{b,d\},\ \{c\},\ \{e\},\ \{\bot\}$\\
$3$&
Same as round $2$\\

\bottomrule
\end{tabular}
\end{center}
The stable quotient has four real task states, and the non-trivial class is $\{b,d\}$. Taking $a<b<c<d<e<\bot$ as the enumeration, the normal representative element satisfies \[
 N(a)=a,\quad N(b)=N(d)=b,\quad N(c)=c,\quad N(e)=e,
 \quad N(\bot)=\bot.
\]. In the real state, the shortest word length to reach the target is \[
 2,\quad\infty,\quad1,\quad\infty,\quad0.
\], where $\infty$ means that there is no finite hit word.\end{proposition}
\begin{certificate}
[Full list, distinguishing words and absorption categories]Save the above two finite tables, $\widehat F_\lambda(\bot)=\bot$, independent tags without definition and fee, and stable quotient transfer \[
 [a]\mapsto[c],\qquad [b,d]\mapsto[b,d],\qquad
 [c]\mapsto[e],\qquad[e]\mapsto[e],\qquad[\bot]\mapsto[\bot].
\] to distinguish $c$ from $a$ using $\lambda$; to distinguish $a$ from $b$ uses $\lambda\lambda$. The indistinguishable certificates of the same type $b,d$ are the same as the current observation and the successors are all $d$. The normal form is decoded as $D(a)=D(b)=D(c)=0$ and $D(e)=1$, and the $\bot$ mark is read according to the convention.\end{certificate}
\begin{proof}
In round zero, the four non-target states are merged, and the target and undefined marks are independent. In the first round, only the successor of $c$ enters the target block, so $c$ is separated from the quad block. In the second round, the successors of $a$ enter $\{c\}$, and the successors of $b,d$ all enter the block containing $d$, so $a$ is branched out again. In the third round, all members of each existing block have the same successor block, and the partition is stable. For $b,d$, the empty word observations are all zero; the positive length words are all in $d$ after the first step, and are always fixed thereafter. All observations are summarized by word length, and the legality and cost tests are the same.

For $a,b$, the empty word and one-step observation are both zero, and they are at $e,d$ in two-step, so the two-step word distinction is not only effective, but also has the smallest length. The hit paths $a\to c\to e$ and $c\to e$ give the upper bound; the actual table values of zero step and one step give the corresponding lower bound. All positive-length successors of $b,d$ are $d$, so they cannot be finitely hit; $e$ is initially on the target. The cost of any legal $k$ step is $k$, which is obtained by induction from $C_0=0$ and $C_{k+1}=C_k+1$. It is also true on the stable quotient. \[
\begin{gathered}
 \operatorname{FullTransitionTable}
 \Rightarrow\mathcal B_0\Rightarrow\mathcal B_1\Rightarrow\mathcal B_2=\mathcal B_3,\\
 \operatorname{SameAbsorbingSuccessor}(b,d)
 \Rightarrow b\equiv_\infty d,\\
 U(F_\lambda^2a)=1\ne0=U(F_\lambda^2b)
 \Rightarrow a\not\equiv_\infty b,\\
 \operatorname{StableClasses}\land\operatorname{ChosenRepresentatives}
 \Rightarrow DN=\widehat U\land N^2=N.
\end{gathered}
\] Here $\mathcal B_k$ represents the $k$ round of partitions.\end{proof}


\section { Necessary refinement of new observation and dynamic tasks }\label{exp:gov:new-observation-section}


\begin{proposition}
[Add the necessary refinement and rigorous information increment of new observations]\label{hgov:prop:refine}
If there is already a quotient $q:X\to Q$ and the new target observation $h$ takes different values on a certain $q$ fiber, then there is no accurate reader that only relies on $q(x)$ to recover all $h(x)$. If $h$ can be obtained from actual input or verifiable new evidence, then \[
q'(x)=(q(x),h(x)),\qquad r(q(x),h(x))=q(x)
\] satisfies $rq'=q$ and the new observation is retained. In the finite dynamic task domain, the dynamics can be fully refined by subsequent saturation.\end{proposition}
\begin{certificate}
[Identical images and different values and sources of evidence]Save $x_1,x_2$, $q(x_1)=q(x_2)$, $h(x_1)\ne h(x_2)$, as well as the input sources, evaluation rules, versions and actual results of the newly added observations.\end{certificate}
\begin{proof}
If there is $h=Dq$, the identical input will give $h(x_1)=h(x_2)$, which is inconsistent with the certificate. The joint observation directly gives $rq'=q$ and the second coordinate readback; the dynamic conclusion is obtained by the theorem \ref{hgov:thm:saturation}. Thus \[
\begin{gathered}
\operatorname{SameOldCode}\land\operatorname{DifferentTarget}
\Rightarrow\operatorname{NoOldDecoder},\\
\operatorname{Available}(h)\Rightarrow\operatorname{Refinement}(q')
\Rightarrow\operatorname{NewReadback}.
\end{gathered}
\] all old quotient algorithms can be run on the new representation via $r$ with new observations having independent readback, so solvability is strictly increased in this family of tasks.\end{proof}


\section { quotient policy fallback, capability reuse and complete cost }\label{exp:gov:reuse-section}


\begin{theorem}
[Certified growth of finite capability library]\label{hgov:thm:library}
Assume that every package in the initial capability library passes type, scope, task readback, and path checks, and that these checks are performed according to proven validation rules. Only new packages that pass the same checks are added each time, and existing packages and their versions are retained. Then the library in any finite growth stage still consists of certified packages; the reuse of packages is only called when the corresponding applicable conditions are met.\end{theorem}
\begin{certificate}
[Package number induction and calling preconditions]Register the initial library, parent dependencies of each newly added package, verification results, applicable domains, and prerequisite witnesses for each call.\end{certificate}
\begin{proof}
establish by induction the number of increases. The initial conclusion is established by the antecedent; in the induction step, the old package has not changed and the new package has passed the same check, so the union still satisfies packet-by-packet authentication. When calling, replace the actual input into the conditional conclusion of the package, and obtain the corresponding result by reliable instantiation of theorem \ref{hgov:thm:atomsound}. \[
\operatorname{Certified}(\mathcal L_n)\land\operatorname{Checked}(K_{n+1})
\Rightarrow\operatorname{Certified}(\mathcal L_n\cup\{K_{n+1}\}),
\quad
\operatorname{Input}\in\mathsf{Dom}(K)\Rightarrow\operatorname{ValidReuse}.
\]\end{proof}


\begin{theorem}
[Ranking strategy fallback and repeated calls on formal types]\label{exp:gov:policy-reuse}
Take the exact stable quotient of theorem ~ \ref{hgov:thm:saturation}, restricted to the real state and its image. Let $R\subseteq Q_*\setminus\{q_*(\bot)\}$ be the registered applicable domain, and $\bar G$ be its target observation quotient state set equal to $1$. Given the strategy $\pi:R\setminus\bar G\to L$ and the ranking $\rho:R\to\mathbb N$, if \[
\begin{aligned}
 &\rho(z)=0\Longleftrightarrow z\in R\cap\bar G,\\
 &\bar e_{\pi(z)}(z)=1,
 \qquad \bar F_{\pi(z)}(z)\in R,\\
 &\rho(\bar F_{\pi(z)}(z))<\rho(z)
 \qquad(z\in R\setminus\bar G).
\end{aligned}
\] is required, then for any $x$ that satisfies $q_*x\in R$, press $a_k=\pi(q_*x_k)$ to select the specific action, at most through $\rho(q_*x)$ Step to reach $G$. Each step, each cost prefix, and final state specifies that observations are identical to the quotient level run. Strategy and ranking can be used as a public certificate for a capability package; each call must still verify the version, enter the applicable domain, and open the package's prerequisites.

If the one-time construction authentication cost is $B\ge0$ and the actual cost of the $t$ complete call is $c_t\ge0$, and is accounted for with non-overlapping primitives, then the total workload of $K$ calls is $W_K=B+\sum_{t=1}^Kc_t$. Compared with the benchmark of completing the same output obligation and the actual cost is $a_t$, the necessary and sufficient condition for strict savings is $\sum_{t=1}^K(a_t-c_t)>B$.\end{theorem}
\begin{certificate}
[Strategy table, drop table and sequential call ledger]Save the entire exchange table for $q_*$, $R,\bar G,\pi,\rho$, and the drop witness for each non-target quotient state. The $t$ call saves the proof, current version, actual action and fee prefix of the input $q_*^{-1}(R)$. $c_t$ includes input reading, encoding or maintenance, policy table lookup, real actions, result readback and current verification; re-authentication caused by migration is independently charged to the one-time account or the corresponding call account, and is only counted once. The cost comparison saves the same tasks, acceptance and output obligations as the baseline and this plan.\end{certificate}
\begin{proof}
Order $z_0=q_*x_0$. At step zero, the actual state and the quotient state satisfy the binding relationship, and the costs are both zero. If $z_k\notin\bar G$ and $q_*x_k=z_k$, the quotient availability is $1$ with domain factorization. The launch $a_k=\pi(z_k)$ is legal in $x_k$; the successor is exchanged to \[
 q_*x_{k+1}=q_*F_{a_k}x_k
 =\bar F_{a_k}q_*x_k=z_{k+1}\in R.
\] and the cost formula is given to $c_{a_k}(x_k)=\bar c_{a_k}(z_k)$, together with the equal parts and launches the next fee prefix equal. This concludes the execution length. Each non-target step decreases the non-negative integer ranking by at least one, so there must be rank zero before more than $\rho(z_0)$ non-target steps; the target test has been accurately preserved, so the specific state at this time belongs to $G$. When $\rho(z_0)=0$, the zero step is successful and the initial base case is covered.

The number of calls is also establish by induction: $W_0=B$, and the non-duplicate primitives are accounted for $W_{K+1}=W_K+c_{K+1}$, so the total workload formula is obtained. The workload difference between the two plans is $\sum_{t=1}^Ka_t-W_K=\sum_{t=1}^K(a_t-c_t)-B$, and its positivity is the necessary and sufficient condition. Therefore \begin{gather*}
 \operatorname{ExactQuotient}\land\operatorname{LegalPolicy}
 \land\operatorname{RankingDecrease}\land(q_*x\in R)\\
 \Rightarrow\operatorname{LegalLiftedPrefixes}
 \Rightarrow\operatorname{FiniteGoalHitAndCostEquality},\\
 \operatorname{VersionAndPremisesCheckedEachCall}
 \Rightarrow\operatorname{CertifiedReuse},\\
 \operatorname{DisjointCostLedger}
 \Rightarrow W_K=B+\sum_{t=1}^Kc_t
 \Rightarrow\left(W_K<\sum_{t=1}^Ka_t
 \Longleftrightarrow\sum_{t=1}^K(a_t-c_t)>B\right).\qedhere
\end{gather*}\end{proof}


\begin{remark}
[Calculation obligations for finite tables]The termination bound of partition refinement measures strictly the number of splitting rounds; reading the entire input table, comparing table entries, and generating certificates all have costs. If the table entry is an arbitrary object that has not yet been evaluated, or the real action table has not yet been obtained, the finite state number itself cannot replace these calculation antecedents. The benefit of capability reuse is determined by the actual cost listed in the above formula, and the existence of the formal type itself does not induce any implementation acceleration.\end{remark}


\chapter{Law-Object Migration, Reflexive Governance, and Joint Preservation}
\label{exp:ch:governance}


\label{hgov:sec:governance}

Law object updates involve states, actions, observations, complexes, and proof rules simultaneously. This chapter registers each of these updates as typed mappings, and demonstrates what conclusions they hold together in finite composition, certificate transport, and actual execution. The old version of the mathematical assertion continues to be interpreted according to the original object; the new version of the assertion is established with the new migration witness. This approach is consistent with the public document's task-sufficient quotient, representation exchange, and atomic certificate closure continuation \cite{S38,GFrameworkSynthesis}.

\section { Type migration, observation decoding and composite readback }\label{exp:gov:migration-section}


\begin{definition}
[law object update and task migration certificate]\label{hgov:def:migration}
The law object $\mathfrak S_n$ in stage $n$ has task complex $C^{(n)}$, state domain $X_n$, task identity $I_n$, target set $G_n$, action, observation, verification core and capability library. The update consists of the following data: state migration $T_n:X_n\to X_{n+1}$, chain mapping $j_n:C^{(n)}\to C^{(n+1)}$, action correspondence $\tau_n$, observation decoding $D_n$, version mapping and rule reliability proof. Fixed identity field $\mathcal I$, requiring \[
I_{n+1}T_n=I_n,\qquad
\partial_{n+1}j_n=j_n\partial_n,\qquad
T_n(G_n)\subseteq G_{n+1}.
\] to specify that old actions, observations and fees meet \[
\begin{aligned}
 \operatorname{Pre}_{n,a}(x)
 &\Longleftrightarrow
 \operatorname{Pre}_{n+1,\tau_n(a)}(T_nx),\\
 T_nF_{n,a}x&=F_{n+1,\tau_n(a)}T_nx
 &&\text{\shortstack[c]{\scriptsize on the\\[-1pt]\scriptsize corresponding\\[-1pt]\scriptsize admissible\\[-1pt]\scriptsize domain}},\\
 O_n&=D_nO_{n+1}T_n,\\
 c_{n,a}(x)&=c_{n+1,\tau_n(a)}(T_nx)
 &&\text{\shortstack[c]{\scriptsize on the\\[-1pt]\scriptsize corresponding\\[-1pt]\scriptsize admissible\\[-1pt]\scriptsize domain}}.
\end{aligned}
\] target forward to remain separately registered as $\mathbf1_{G_n}(x)=1\Rightarrow\mathbf1_{G_{n+1}}(T_nx)=1$. When $C^{(n)}$ is a chain implementation of the task action graph, $j_{n,0}[x]=[T_nx]$ is also required on the registration vertex; the repair change complex uses its declared change bridge. If a new readback is required to accept output, a set $\mathcal A_{n+1}^{\mathrm{read}}\subseteq G_{n+1}$ and \[
 r_n:\mathcal A_{n+1}^{\mathrm{read}}\longrightarrow G_n,\qquad
 O_nr_n=D_nO_{n+1},\qquad I_nr_n=I_{n+1}.
\] covering the actual delivered output is given. Each readback actually verifies that its input belongs to the definition domain; the multi-stage readback verifies the continuation domain one by one. Rule updates are verified against a previously authenticated verifier or against an external mathematical proof that has been given, with the open premise in the proof passed along with the certificate. If the current law selection and verification protocols are encoded as objects in the next phase, the same migration obligations continue to apply.\end{definition}


\begin{theorem}
[Migrate certificate compounding and accept the exact continuation field of the readback]\label{exp:gov:migration-composition}
Assume that stages $0\to1$ and $1\to2$ respectively satisfy the definitions ~ and \ref{hgov:def:migration}. Clearly write out the observation type $O_i:X_i\to Y_i$ and decoding types $D_0:Y_1\to Y_0$ and $D_1:Y_2\to Y_1$. Define \[
\begin{aligned}
 T_{02}&=T_1T_0,& j_{02}&=j_1j_0,\\
 \tau_{02}&=\tau_1\tau_0,& D_{02}&=D_0D_1.
\end{aligned}
\], then the above data meets the identity, chainability, old action legality, successor, observation, target and old action cost obligations of compound migration. $j_{02,0}[x]=[T_{02}x]$ is also satisfied when the vertex binding of the original task action graph is applicable.

If the single-step accepted readbacks are $r_0:\mathcal A_1^{\mathrm{read}}\to G_0$ and $r_1:\mathcal A_2^{\mathrm{read}}\to G_1$, then the registration domain of the composite readback is \[
 \mathcal A_{02}^{\mathrm{read}}
 =\{z\in\mathcal A_2^{\mathrm{read}}:
             r_1z\in\mathcal A_1^{\mathrm{read}}\}.
\]. In this domain, $r_{02}=r_0r_1$ satisfies \[
 r_{02}:\mathcal A_{02}^{\mathrm{read}}\to G_0,
 \qquad O_0r_{02}=D_{02}O_2,
 \qquad I_0r_{02}=I_2.
\]. The identity migration and the above composite have unit law and associative law; partial readback The associativity of also requires that the fields of the written continuation be equal.\end{theorem}
\begin{certificate}
[Field-by-field composition and readback domain membership proof]Save two single-step migration certificates, all mapped source targets, and the compounding sequence. For the actual new output $z$, save the evidence of $z\in\mathcal A_2^{\mathrm{read}}$ and $r_1z\in\mathcal A_1^{\mathrm{read}}$ respectively, and then call $r_0$. Observation decoding is fully defined in the registered observation value domain; if the implementation uses a partial decoder, it is necessary to successively save the proof that the actually called observation value belongs to its definition domain.\end{certificate}
\begin{proof}
Identity, linkage and observation are calculated according to the following substitutions: \[
\begin{aligned}
 I_2T_{02}&=I_2T_1T_0=I_1T_0=I_0,\\
 \partial_2j_{02}
 &=\partial_2j_1j_0=j_1\partial_1j_0=j_1j_0\partial_0,\\
 O_0&=D_0O_1T_0=D_0D_1O_2T_1T_0.
\end{aligned}
\] For each old label $a$, two preconditions are equivalent to \[
 \operatorname{Pre}_{0,a}(x)
 \Longleftrightarrow
 \operatorname{Pre}_{1,\tau_0(a)}(T_0x)
 \Longleftrightarrow
 \operatorname{Pre}_{2,\tau_{02}(a)}(T_{02}x).
\]. On this legal domain, two successor exchanges are to \[
 T_{02}F_{0,a}x
 =T_1F_{1,\tau_0(a)}T_0x
 =F_{2,\tau_{02}(a)}T_{02}x.
\]. The cost equation is concatenated in the same way. The goal includes rolling out $T_0x\in G_1$ and $T_{02}x\in G_2$ sequentially from $x\in G_0$. Vertex binding obtained by $j_{1,0}j_{0,0}[x]=j_{1,0}[T_0x]=[T_1T_0x]$.

For $z\in\mathcal A_{02}^{\mathrm{read}}$, both readback function calls are legal, and the \[
 O_0r_0r_1z=D_0O_1r_1z=D_0D_1O_2z,
 \qquad I_0r_0r_1z=I_1r_1z=I_2z.
\] readback value belongs to the $G_0$ value range declaration from $r_0$. The identity law and associative law of mapping give the overall mapping conclusion. For three consecutive readbacks, no matter what kind of brackets are used, the domain is exactly the set that satisfies \[
 z\in\mathcal A_3^{\mathrm{read}},\quad
 r_2z\in\mathcal A_2^{\mathrm{read}},\quad
 r_1r_2z\in\mathcal A_1^{\mathrm{read}}
\], so some readbacks also have the same value and domain. Zero-order migrations take their respective identity maps; by the above two-step composite induction, any finite migration chain can be constructed. \[
\begin{gathered}
 \operatorname{Migration}(0,1)\land\operatorname{Migration}(1,2)
 \Rightarrow\operatorname{Migration}(0,2),\\
 z\in\mathcal A_{02}^{\mathrm{read}}
 \Rightarrow\operatorname{BothReadbacksDefined}
 \Rightarrow\operatorname{AcceptedOldReadback},\\
 \operatorname{TypedComposition}\land\operatorname{MatchedPartialDomains}
 \Rightarrow\operatorname{AssociativeFiniteMigration}.
\end{gathered}
\]\end{proof}


\section { Version freeze and logical scope of atomic transport }\label{exp:gov:version-atoms-section}


\begin{definition}
[Version freezing, atom-by-atom transportation and prerequisite ledger]\label{hgov:def:certificate-transport}
Interpretation library $\mathfrak M_{\le n}$ stores registered objects as of stage $n$. Every object, function, quantifier domain, input, and certificate has an irreversible version identifier. The new stage expands the interpretation library, and the existing version fields continue to be interpreted according to their original values. All source objects and quantifier fields of proposition $P^{[j]}$ are bound to the registered version $j$.

For a finite certificate to be migrated, expand the local open premise into leaf nodes. Note that its node predicate is $P_v$ and its parent collection is $\operatorname{Par}(v)$. The domain of partial translation $\Theta_n$ covers all atoms, compound formulas and rule instances actually rewritten this time; each of its outputs has a clear source target, frequency and version. Keep the node to directly read the original version $P_v$, and rewrite the node to output the registered new proposition $P'_v=\Theta_nP_v$. Each rewrite atom $A$ holds the specific transport witness and the conditional truth-preserving formula \[
 \mathfrak M_{\le n+1}\models
 \left(A\land\bigwedge\Delta_{n,A}\right)\Rightarrow\Theta_nA.
\] where $\Delta_{n,A}$ is the set of finite open premises actually used by the transport. Each non-leaf node also saves the verified new rule instance \[
 \mathfrak M_{\le n+1}\models
 \left(
 \bigwedge_{u\in\operatorname{Par}(v)}P'_u
 \land\bigwedge\Delta_{n,v}
 \right)\Rightarrow P'_v.
\]. The reserved branch can use the original rule instance, and the consumed objects and premises continue to carry the original version. Let $\Gamma_n$ be all the unresolved premises of the certificate at this stage, let $\Delta_n$ be the union of premises introduced by this object migration, atomic transportation and new rules, then the output premise ledger is defined as \[
 \Gamma_{n+1}=\Gamma_n\cup\Delta_n.
\]. The proven premises are also accompanied by elimination witnesses; the atomic truth value, rule reliability and premise elimination are saved separately.\end{definition}


\begin{proposition}
[version maintains specific transport with generated atoms]\label{hgov:prop:atomic-transport}
Version Freeze maintains the interpretation of all original version formulas. The definition ~ \ref{hgov:def:migration} gives the following truth-preserving rewrite atom: \[
 \begin{aligned}
 y=F_{n,a}x
 &\Longrightarrow T_ny=F_{n+1,\tau_n(a)}T_nx,\\
 x\in G_n&\Longrightarrow T_nx\in G_{n+1},\\
 O_nx=o&\Longrightarrow D_nO_{n+1}T_nx=o,\\
 \partial_nu=z
 &\Longrightarrow \partial_{n+1}(j_nu)=j_nz,\\
 I_nx=i&\Longrightarrow I_{n+1}T_nx=i .
 \end{aligned}
\] An action atom transports both its legality and its cost. Extra atoms are transported in a truth-preserving manner via the corresponding truth-preserving registrations in the definition ~ \ref{hgov:def:certificate-transport}.\end{proposition}
\begin{certificate}
[Frozen domain and atomic commutation witness]Save version identifiers of old objects and domains, actual values of source atoms, state and chain mappings, target inclusions, observation decodes, and action commutation relations. Each output atom saves the substituted equation and the open premises used. The obstacle certificate of the old complex continues to bind the old complex; the filling conclusion of the new complex saves the new preimage and boundary calculation.\end{certificate}
\begin{proof}
establish by induction the structure of the original version of the formula. Objects, functions and relational interpretations in atoms retain their original values; propositional connections are computed by the same truth functions; typed quantifiers traverse the same frozen domain. Therefore the truth value of each original version of the formula remains. For the action atom, substitute $y=F_{n,a}x$ to get \[
 T_ny=T_nF_{n,a}x=F_{n+1,\tau_n(a)}T_nx.
\]. The precondition equivalence and cost formula respectively give the legality and cost atom transportation. The target atom was obtained from $T_n(G_n)\subseteq G_{n+1}$. The observation atoms and boundary atoms are obtained from \[
 D_nO_{n+1}T_nx=O_nx=o,\qquad
 \partial_{n+1}j_nu=j_n\partial_nu=j_nz
\] respectively, and their identities are the same. The predicate causal chain item by item is \[
 \begin{gathered}
 \operatorname{FrozenObjects}\land\operatorname{FrozenDomains}
 \Longrightarrow\operatorname{HistoricalTruthPreserved},\\
 \operatorname{SourceAtom}\land\operatorname{CheckedTransport}
 \Longrightarrow\operatorname{TargetAtom},\\
 \operatorname{GoalInclusion}\land(x\in G_n)
 \Longrightarrow T_nx\in G_{n+1}.
 \end{gathered}
\]. Each transportation only consumes its listed structural formula and premises.\end{proof}


\begin{proposition}
[Positive formula closure and full formula preserving conditions for one-way atomic transport]\label{exp:gov:formula-transport}
For each object type $s$ actually used this time, the domain mapping $T_s:X_s\to X'_s$ is given. Assume that the translation $\Theta$ maintains the logical structure of the formula, and the transport formula of each registered atom is uniformly established in all specified assignments: \[
 \bigwedge\Gamma\Rightarrow
 \bigl(A(\xi)\Rightarrow(\Theta A)(T\xi)\bigr).
\] Here $\xi$ is a typed assignment, $T\xi$ applies the $T_s$ of its type variable by variable; $\Gamma$ contains all unresolved prerequisites for the above unified transportation certificate. Then for the formula $P$ constructed from the registered atoms finite times $\land,\lor,\exists$, there is \[
 \bigwedge\Gamma\Rightarrow
 \bigl(P(\xi)\Rightarrow(\Theta P)(T\xi)\bigr).
\]. If further each $T_s$ is surjective, and all atomic transports are two-way truth equivalent, then $\neg,\forall$ is allowed. Any registered first-order formula of also has bidirectional equivalence. When there is only one-way atomic transport, negation and the universal quantifier traversing the entire newly added domain do not belong to the above general preservation conclusion.\end{proposition}
\begin{certificate}
[Formula construction trees, quantifier domains, and unified assignment checks]Save the finite construction tree of the formula, the type and version of each variable, the assignment domain covered by the atomic transport, and the prerequisites required for the unified proof of each atom. The existence quantifier saves the source witness and its $T_s$ image; the full formula two-way transport also saves the surjective witness of each domain mapping and the two-way equivalence of the atom. The finite domain implementation replaces uncovered sample checks with a full assignment table.\end{certificate}
\begin{proof}
Fix an explanation that makes $\Gamma$ true, and Induct on the formula construction tree. The atomic base example is the antecedent; the conjunction is obtained by transporting two sub-formulas at the same time, and the disjunction is obtained by transporting the branches that are actually established. If $\exists x\in X_s\,P(x,\xi)$ is established, take its witness $x$; induction hypothesis gives $(\Theta P)(T_sx,T\xi)$, so $T_sx\in X'_s$ is the witness of the target existence quantifier. This completes the structural induction of the positive formula.

Under the conditions of two-way atomic equivalence and surjective, all formulas are establish by induction. Negation is maintained directly by truth equivalence of subformulas. For any $x'\in X'_s$ in the target domain, the surjective property $x$ satisfies $T_sx=x'$; therefore any instance of the source full name conclusion can be transported to this $x'$. The reverse universal conclusion holds for every image like $T_sx$, and the subformula is used for bidirectional equivalence back to the source domain. The reverse of the existential quantifier also uses the preimage of the target witness.

Two finite models illustrate precisely the purpose of additional conditions. First take the zero-element atom $A$ which is false in the source interpretation and true in the target interpretation; $A\Rightarrow\Theta A$ holds, but $\neg A$ cannot be transported as $\neg\Theta A$ . Secondly, take $X=\{0\}$, $X'=\{0,1\}$, $T(0)=0$, the source predicate $P(0)$ is true, the target predicate $P'(0)$ is true and $P'(1)$ is false. The atom is held across all source assignments, $\forall x\in X\,P(x)$ is true and $\forall x'\in X'\,P'(x')$ is false. Therefore, the old full name formula under the \[
\begin{gathered}
 \operatorname{UniformForwardAtoms}\land\operatorname{PositiveFormula}
 \Rightarrow\operatorname{ForwardFormulaTransport},\\
 \operatorname{UniformEquivalentAtoms}\land\operatorname{SurjectiveDomains}
 \Rightarrow\operatorname{FirstOrderTruthEquivalence}.
\end{gathered}
\] version freeze is still quantified in the old domain, so the proposition ~ \ref{hgov:prop:atomic-transport} maintains the original interpretation; this proposition deals with the formula after being actively rewritten to the target version.\end{proof}


\section { certificate dependency graph, openness premise and secure elimination }\label{exp:gov:proof-ledger-section}


\begin{lemma}
[Finite prefix reliability for atom-by-atom certificate transport]\label{hgov:lem:transport-prefix}
If the source finite certificate is established under $\Gamma_n$ and has the version and transportation witness defining ~ \ref{hgov:def:certificate-transport}, then each node will be established under $\Gamma_{n+1}$ after migration. Keep certificate to continue to express the original version of the proposition, and rewrite certificate to express its registered target version of the proposition.\end{lemma}
\begin{certificate}
[Topology prefix and new prerequisites]Save the source certificate topological order, the retention or rewriting identifier of each node, the parent index, the atomic transport formula, the new rule instance, and the local premise set. Let $\Delta_{n,v}^{\mathrm{all}}$ be the new premise union of node $v$ and its transportation certificate. Accumulated set of records for the first $k$ nodes \[
 \Lambda_k=\Gamma_n\cup
 \bigcup_{v\le k}\Delta_{n,v}^{\mathrm{all}},
 \qquad \Lambda_k\subseteq\Gamma_{n+1}.
\]\end{certificate}
\begin{proof}
Assume $\Gamma_{n+1}$ is established. All the premises of the source certificate $\Gamma_n$ are established, and the version is frozen so that the conclusion can continue to be read. Summarize according to the topological order of the new certificate: the retained leaf is established by the original version proof; the rewritten leaf is established by the transport formula of the source atom and the included $\Delta_{n,A}$. If the first $k$ nodes have been established, the parent conclusions of the $k+1$ non-leaf nodes are all before, and its $\Delta_{n,v}$ is also included in $\Gamma_{n+1}$, so the corresponding new rule is introduced $P'_{k+1}$. The same argument gives $\bigwedge\Lambda_k\Rightarrow\bigwedge_{v\le k}P'_v$ when using only the prefix premise. So \[
 \begin{gathered}
 \operatorname{SourceValid}(\Gamma_n)
 \land\operatorname{LeafTransport}
 \land\operatorname{SoundTargetRules}\\
 \Longrightarrow
 \left[
 \bigwedge\Gamma_{n+1}
 \Rightarrow\bigwedge_{v}P'_v
 \right].
 \end{gathered}
\] completes the finite prefix induction, and all new premises are kept in the output ledger.\end{proof}


\begin{theorem}
[Node local premise support and premise elimination with evidence]\label{exp:gov:premise-support}
Take a finite acyclic certificate graph with all predicates, types, and object versions fixed. Each leaf node $v$ saves the verified conditional conclusion $\bigwedge B_v\Rightarrow P_v$; the open premise leaf $P_v=\gamma$ can be obtained by $B_v=\{\gamma\}$. Each non-leaf node saves the rule instance \[
 \left(\bigwedge_{u\in\operatorname{Par}(v)}P_u
       \land\bigwedge\Delta_v\right)\Rightarrow P_v,
\] and recursively registers \[
 B_v=\Delta_v\cup\bigcup_{u\in\operatorname{Par}(v)}B_u.
\], then each node satisfies $\bigwedge B_v\Rightarrow P_v$. For the output set $V_{\mathrm{out}}$, let $\Gamma_{\mathrm{out}}=\bigcup_{v\in V_{\mathrm{out}}}B_v$, then \[
 \bigwedge\Gamma_{\mathrm{out}}
 \Rightarrow\bigwedge_{v\in V_{\mathrm{out}}}P_v.
\]. Therefore, only the ancestor closure of the desired output and all its premises can be retained.

If each $\gamma$ in $E\subseteq\Gamma_{\mathrm{out}}$ has another finite reliable certificate, and only $H=\Gamma_{\mathrm{out}}\setminus E$ is used as the unresolved premise to prove $\bigwedge H\Rightarrow\gamma$, then the output premise can be reduced from $\Gamma_{\mathrm{out}}$ to $H$. The same rules apply to the premises introduced for atomic transport.\end{theorem}
\begin{certificate}
[Support set recursion and independent elimination proof]Save topological order, local premises of each rule, node-by-node set union operation and output ancestor index. Save an independent certificate for each prerequisite to be eliminated, and its open leaves all belong to $H$, and verify that the rules used are reliable, the object versions match, and the dependency graph is loop-free. Leaves from older versions retain the original interpretation, and rewritten leaves retain the corresponding atomic transport and new premises.\end{certificate}
\begin{proof}
Sort by topological order. The conditional conclusion of the leaf node is established by the input certificate. If the conclusions of all parent nodes are true under their supporting premises, the definition of $B_v$ ensures that each parent supporting and rule local premise is included in $B_v$. Therefore, under $B_v$, all the parent conclusions are established, and then $P_v$ is obtained from the registered rules. This gives the induction step and overrides the initial base case with no parent. The union of output premises simultaneously covers all output supports, thus leading to the output conjunction. The ancestor closure contains all parent nodes of each retained node, so deleting the remaining nodes does not change the above derivation.

Assume $H$ holds. Each independent elimination certificate gives a $\gamma\in E$, so all members of $H\cup E=\Gamma_{\mathrm{out}}$ are established; all outputs are obtained by substituting the proven output condition conclusions. This process uses actual proof to eliminate the premise without changing the version or meaning of any original proposition. \[
\begin{gathered}
 \operatorname{AcyclicRules}\land\operatorname{LeafSequents}
 \Rightarrow\bigwedge_v\left(\bigwedge B_v\Rightarrow P_v\right),\\
 \operatorname{OutputAncestorClosure}
 \Rightarrow\left(\bigwedge\Gamma_{\mathrm{out}}
 \Rightarrow\bigwedge_{v\in V_{\mathrm{out}}}P_v\right),\\
 \bigwedge_{\gamma\in E}\operatorname{Certified}(H\Rightarrow\gamma)
 \Rightarrow\left(\bigwedge H\Rightarrow
             \bigwedge_{v\in V_{\mathrm{out}}}P_v\right).
\end{gathered}
\]\end{proof}


\section { Compatibility of complex version, non-zero class and new padding }\label{exp:gov:homology-version-section}


\begin{theorem}
[Obstacle maintenance, class disappearance and retraction criteria under version freezing]\label{exp:gov:homology-version}
Let $C_\bullet,C'_\bullet$ be the chain complex on domain $\mathbb K$, and $j:C_\bullet\to C'_\bullet$ be the chain mapping. Given the old closed chain $\omega\in C_k$ and the linear functional $\lambda:C_k\to\mathbb K$, if \[
 \partial\omega=0,\qquad
 \lambda\partial_{k+1}=0,\qquad\lambda(\omega)\ne0.
\] is satisfied, then the old version asserts that $[\omega]_{H_k(C)}\ne0$ has an explicit certificate. This assertion can be true simultaneously with the fact that the new version $h\in C'_{k+1}$ satisfies $\partial'h=j\omega$.

If there are also a chain map $r:C'\to C$ and an operator of degree $+1$,
$H:C_k\to C_{k+1}$,
satisfying in every degree
\[
 rj-\mathrm{id}_C=\partial H+H\partial,
\],
then $j_*:H_k(C)\to H_k(C')$ is injective, and the image of an old nonzero class remains nonzero. Thus transporting nonzero classes requires a proved injectivity condition of this kind, or a direct nonvanishing certificate in the target complex.
\end{theorem}
\begin{certificate} [Old separating functional, new filling, and retraction homotopy]
The old certificate records the complex version, coefficient field, bases, boundary matrix, $\omega$, and $\lambda$, and checks $\partial\omega=0$, $\lambda\partial=0$, and $\lambda\omega\ne0$ item by item. The new filling separately records the new version, $j,h$, and $\partial'h=j\omega$. If transport preserving nonvanishing is used, verify the chain-map property of $r$ on each basis element, together with $rj-\mathrm{id}=\partial H+H\partial$. The type conventions for chain complexes, induced homology, and chain homotopies are given in \cite{ExtHatcher,ExtWeibel,S38}.
\end{certificate}
\begin{proof}
If the old chain were a boundary $\omega=\partial u$, then $\lambda(\omega)=\lambda\partial u=0$, contradicting the premise; hence the old class is nonzero. The chain-map property gives only $\partial'j\omega=j\partial\omega=0$, so the induced homology image is defined. If there is additionally a new filling $h$, its image class is zero by definition, whereas the old boundary matrix is unchanged and the old separation certificate remains valid.

The following finite model directly checks this coexistence. Take \[
 C_1=\mathbb K e,\qquad C_0=\mathbb K v,
 \qquad\partial_1=0,
 \qquad C_k=0\ (k\notin\{0,1\}).
\] and specify $C'_0=C_0,C'_1=C_1,C'_2=\mathbb K f$, and the remaining times are zero, specifying $\partial'_2f=e$ and $\partial'_1=0$. $j$ is identical in degrees zero and one, and is a mapping starting from the zero space in other degrees. The bounding square is zero, because the only possible non-zero composition is $\partial'_1\partial'_2=0$; the chainability of $j$ holds directly from degree to degree. So \[
 H_1(C)=\mathbb K e,\qquad
 H_1(C')=\mathbb K e/\mathbb K e=0.
\] takes $\lambda(e)=1$ to get the old non-zero certificate, and takes $h=f$ to get the new padded certificate.

Now assume that the retraction homotopy condition holds. If the image of a closed chain $z\in C_k$ is the boundary $jz=\partial'h$, then \[
 rjz=\partial(rh),\qquad
 rjz-z=\partial(Hz)+H\partial z=\partial(Hz),
\] and therefore $z=\partial(rh-Hz)$, the original class is zero. This proves $\ker j_*=0$; applied to $\omega$ that excludes new padding. \[
\begin{gathered}
 \operatorname{OldCycleAndSeparator}
 \Rightarrow [\omega]_{H_k(C)}\ne0,\\
 \operatorname{ChainMap}(j)\land(\partial'h=j\omega)
 \Rightarrow [j\omega]_{H_k(C')}=0,\\
 \operatorname{HomotopyRetraction}(r,j,H)
 \Rightarrow\operatorname{Injective}(j_*)
 \Rightarrow\operatorname{NonzeroClassPreserved}.
\end{gathered}
\] The superscript complexes in the first two lines are different, so they describe compatible conclusions in different versions.\end{proof}


\begin{remark}
[Separate verification of version keys and target constraints]An obstacle conclusion binds at least the coefficient field, the complex, the boundary operator, the class representative and the allowed filling space. When the operation family or relationship surface is allowed to change, a new version and a new nuclear image proof should be generated for the new complex. The old non-fillable and the new fillable can jointly form a capability expansion certificate; if a homology class also undertakes task information that must be protected, corresponding observation retention or retraction conditions need to be registered in task migration.\end{remark}


\section { finite Inductive invariant of interleaved migration and execution }\label{exp:gov:finite-governance-section}


\begin{theorem}
[The task of finite law object migration is to maintain the certificate]\label{hgov:thm:migration}
If the rules of the initial law object are reliable and each step satisfies the definition ~ \ref{hgov:def:migration} and the definition ~ \ref{hgov:def:certificate-transport}, then each finite compound $T_{0:n}=T_{n-1}\cdots T_0$ and $j_{0:n}=j_{n-1}\cdots j_0$ Maintain identity, chainability, provision of old runs and their fees. After each old successfully runs the migration, it will still enter the new target, and the certificate after the migration specifically meets \[
 I_nT_{0:n}=I_0,\qquad
 \partial_nj_{0:n}=j_{0:n}\partial_0,\qquad
 T_{0:n}(G_0)\subseteq G_n.
\] is valid under the complete prerequisite account book \[
 \Gamma_n=\Gamma_0\cup\bigcup_{0\le k<n}\Delta_k
\]; the original version certificate and the certified target version certificate are saved separately.\end{theorem}
\begin{certificate}
[Summary of update times]Save each stage's object type, matrix, state map, action precondition, task identity, target inclusion, observation readback, cost, rule proof, and version parent. Save the composite invariant \[
I_nT_{0:n}=I_0,\qquad
\partial_nj_{0:n}=j_{0:n}\partial_0,\qquad
T_{0:n}(G_0)\subseteq G_n.
\] and observe the decoding to save $D_{0:n}=D_0D_1\cdots D_{n-1}$ and $O_0=D_{0:n}O_nT_{0:n}$. Each certificate saves the atomic transport with $\Gamma_{n+1}=\Gamma_n\cup\Delta_n$; the accepted readback saves the domain proof for each actual call.\end{certificate}
\begin{proof}
Zero-degree composition is identity, and the base case is established. Assuming that the $n$ stage is established, then \[
I_{n+1}T_nT_{0:n}=I_nT_{0:n}=I_0,
\] \[
\partial_{n+1}j_nj_{0:n}
=j_n\partial_nj_{0:n}=j_nj_{0:n}\partial_0.
\] if $x\in G_0$, the induction hypothesis gives $T_{0:n}x\in G_n$, and the target contains further gives $T_nT_{0:n}x\in G_{n+1}$. The one-step exchange of actions is equivalent to the precondition by induction by word length, keeping all the specified old runs, and the cost of each step is the same so that the sum of all cost parts is the same. The observation corresponds to \[
 O_0
 =D_{0:n}O_nT_{0:n}
 =D_{0:n}D_nO_{n+1}T_nT_{0:n}.
\] so the final state of the old execution via $T_{0:n+1}$ is still accepted by the new target, which specifies the observation to be read back via the above decoding. When readback is required to accept new output, the corresponding admissible readback is given by successive domain checks and the target value domain of $r_n$.

For certificate , the induction hypothesis gives that it is valid under $\Gamma_n$. Lemma ~ \ref{hgov:lem:transport-prefix} transports the actually registered atoms, leaves and rules to the new stage, and obtains the conclusion under $\Gamma_{n+1}=\Gamma_n\cup\Delta_n$. Version freezing simultaneously preserves the source of the old proposition and the original object interpretation. So the induction invariant is \[
\begin{gathered}
\operatorname{ValidStage}(n)\land\operatorname{MigrationCertificate}(n)
\land\operatorname{AtomTransport}(n)\\
\Rightarrow
\operatorname{IdentityPreserved}(n+1)
\land\operatorname{GoalPreserved}(n+1)\\
\land\operatorname{ChainPreserved}(n+1)
\land\operatorname{OldRunAndCostPreserved}(n+1)\\
\land
\left[\bigwedge\Gamma_{n+1}
\Rightarrow\operatorname{TargetCertificateValid}(n+1)\right].
\end{gathered}
\] mathematical induction gives all finite stages.\end{proof}


\begin{theorem}
[Union invariants when old ability runs are interleaved with law object updates]\label{exp:gov:interleaved-governance}
Assume that finite stage migration satisfies the definition ~ \ref{hgov:def:migration} and the definition ~ \ref{hgov:def:certificate-transport}. Initial reference state $x_0\in X_0$, actual state $y_0=x_0$, current stage $n=0$. Events in a finite event sequence are allowed to be of the following two categories:\begin{enumerate}
[label=(\roman*)]\item
In stage $n$, the current compilation label $\tau_{0:n}(a)$ of a primitive action $a$ is executed; the reference state is updated according to $F_{0,a}$, and the actual state is updated according to the corresponding current action. Check the legality of the original action before execution.\item
Migrate to stage $n+1$, the actual status is updated according to $T_n$, and the reference status is maintained; at the same time, the certificate is transported and the prerequisite ledger is expanded.\end{enumerate}
Here, all zero-order composites take identity mapping. For any finite event prefix, if the current reference state is $x$ and the actual state is $y$, then the partial sum of the \[
\begin{gathered}
 y=T_{0:n}x,\qquad I_ny=I_0x,\qquad
 O_0x=D_{0:n}O_ny,\\
 \partial_nj_{0:n}=j_{0:n}\partial_0.
\end{gathered}
\] reference and actual task action costs is equal; if the reference state belongs to $G_0$, the actual state belongs to $G_n$. All registered certificates are valid as long as the accumulation has not been cancelled. The actual workload required for updating and verification is included in the additional cost based on non-duplicate primitives.\end{theorem}
\begin{certificate}
[Event tables, compilation labels, and dual-state prefixes]Save the category, current phase, reference state, actual state, action compilation chain, precondition, cost or migration matrix for each event. Maintain the prefix table of $(x,y,n,T_{0:n},j_{0:n},D_{0:n},\Gamma)$. An additional cost of $\kappa_i\ge0$ is saved for migration and revalidation; it does not double count against the saved task action charges. The applicable domain, target and open premise of the original capability package will be verified item by item during actual invocation.\end{certificate}
\begin{proof}
establish by induction the number of events. In the empty event prefix $n=0$, all composites are identical, $y_0=x_0$, the sum of the fee parts is zero, and the certificate condition is given by the initial registration.

Assume that the current prefix satisfies the conclusion. If the next event is an action, the legality equivalence of compound migration ensures that the current compilation action is legal in $y$, and the single-step equality of the cost of the successor exchange to \[
\begin{aligned}
 y^+&=F_{n,\tau_{0:n}(a)}y
 =F_{n,\tau_{0:n}(a)}T_{0:n}x\\
 &=T_{0:n}F_{0,a}x=T_{0:n}x^+.
\end{aligned}
\] is added to the already equal partial sum, keeping the next prefix cost equal. The identity and observation correspondence are obtained from the new binding state and the existing composite migration; the certificate instantiates the registered rules according to legal actions.

If the next event is a migration, $x^+=x$ and \[
 y^+=T_ny=T_nT_{0:n}x=T_{0:n+1}x^+.
\] Theorem ~ \ref{exp:gov:migration-composition} gives the updated identity, observation and linkage equations. This event does not add reference or actual task action charges, but its implementation cost is credited to $\kappa_i$. The certificate is transported by lemma ~ \ref{hgov:lem:transport-prefix}, and the premise is saved by $\Gamma^+=\Gamma\cup\Delta_n$. This completes the induction step for both types of events. Finally, $x\in G_0$ is included with the composite target given $y\in G_n$ . If the original action cost of the entire event sequence is $C_{\mathrm{task}}$, the total cost of accounting based on non-repeating primitives is $C_{\mathrm{task}}+\sum_i\kappa_i$. \[
\begin{gathered}
 \operatorname{PrefixInvariant}(k)\land
 \bigl(\operatorname{CompiledAction}\lor\operatorname{CertifiedMigration}\bigr)\\
 \Rightarrow\operatorname{PrefixInvariant}(k+1),\\
 \operatorname{PrefixInvariant}(0)
 \Rightarrow\forall k\le K\ \operatorname{PrefixInvariant}(k),\\
 \operatorname{FinalReferenceGoal}\land\operatorname{PrefixInvariant}(K)
 \Rightarrow\operatorname{ActualGoalAndDecodedObservation}.
\end{gathered}
\]\end{proof}


\section { Joint acceptance and strict capacity expansion }\label{exp:gov:joint-section}


\begin{theorem}
[Joint maintenance and strict expansion of heterogeneous evolutionary governance]\label{hgov:thm:main}
In the same task domain, given the following verified data: typed law object and atomic derivation kernel; task action subgraph and relational surface registered in this finite run; maintained task state quotient and chain mapping of respective types; geometric observation and task bridging; closed defects and allowed repairs; executable actions, preconditions and costs; actual termination certificate; formal decoding and version migration with atomic transport. Register the unresolved prerequisite union of all interfaces and migration certificates with $\Gamma_*$. Then under $\Gamma_*$, all finite operations compounded by these interfaces simultaneously satisfy:\begin{enumerate}
[label=(\roman*)]\item
Semantic differential is mapped and exchanged with the specified chain, and the closed cochain points on the authentication boundary are maintained;\item
Homological obstruction is determined by clear nuclear images or computational complex shapes, and the actual repair has recalculated filling and execution witnesses;\item
task identity, prescribed observations, old target acceptance, and old execution costs are maintained in legal migrations;\item
The reachable set is non-decreasing when the old and new executable control sets are included; it is strictly expanded when there is a verified difference set witness;\item
The results that pass the acceptance can be maintained authenticated by task normal form readback, finite capability library and subsequent legal reuse.\end{enumerate}
The task family of theorem \ref{hgov:thm:joint} satisfies the above conditions at the same time and has strictly reachable expansion.\end{theorem}
\begin{certificate}
[Atomic Acceptance Form for Joint Conclusions]Each entry continues to expand into the book's registration primitives:\begin{center}
\small\begin{tabularx}
{\textwidth}{P{.10\textwidth}YY}\toprule
entrance&
math obligations&
Actual witness\\
\midrule
U01&
Type, object domain and task identity&
input, version, domain\\
U02&
Various types of complex $\partial^2=0$, $d=\partial^*$&
finite execution subgraphs, boundary entries, and basis-wise products\\
U03&
State and chain exchange, domain and cost factorization&
Vertex bindings, action tables, and finite prefixes\\
U04&
Geometric variations, smooth implementations and observed responses&
Metrics, initial values, parameterized paths, and boundary integrals\\
U05&
closed defects and obstacles&
Nuclear basis, projection or separation functional\\
U06&
Allow padding or compatible expansion&
Actual operators, matrices, permissions and antecedents\\
U07&
Performance, Ranking and Cost&
Each step status, action, descent amount\\
U08&
$DN=O$ and applicable domains&
Encoding, decoding, success and failure records\\
U09&
Target inclusion, versioning, and atom-by-atom shipping&
Truth-preserving witness, readback domain and premise union\\
\bottomrule
\end{tabularx}

\end{center}
The task keys, source data, and valid fields of all portals are compatible; each portal retains its own complex type and state, changes, or geometry bridges. The union acceptance is the conjunction of the acceptance predicates and continues to save $\Gamma_*$.\end{certificate}
\begin{proof}
U01 and rule reliability give the validity of atomic judgment. U02--U03 are given by theorem \ref{hgov:thm:complex}, lemma \ref{hgov:lem:semanticdiff} and theorem \ref{hgov:thm:quotient} on their respective types, giving the differential, task execution and cost preservation. The geometric conclusion of U04 has been verified in response to entering the task data, and the smooth implementation executes the subgraph with actual parameterized paths and boundary exchange connections. U05--U06 uses the book's Hom filling, nuclear image criteria and compatible expansion structures, and U07 implements repairs with specific admissible actions and termination rankings. Control set inclusion makes every old admissible path still belong to the new family, so the reachable set contains; the actual hit paths in the difference set give strict inclusion. U08 is given readback and capability library retention by theorem \ref{hgov:thm:nf} and theorem \ref{hgov:thm:library}; U09 is given target forward retention, old cost retention and conditional certificate migration by theorem \ref{hgov:thm:migration}. These results share the same task interface, and the premise is passed through the union, so it can be compounded according to the finite certificate graph: \[
\begin{gathered}
\bigwedge_{i=1}^{9}\operatorname{Verified}(U_i)
\land\operatorname{CompatibleInterfaces}\\
\Rightarrow
\left[
\bigwedge\Gamma_*\Rightarrow
\operatorname{TaskPreserved}
\land\operatorname{RepairExecuted}
\land\operatorname{ReadbackValid}
\right].
\end{gathered}
\] The joint instance gives the specific matrix, bridge, $m$ step path and normal form $N_m(p,s,u,f)=(m,p,s,u)$, and the decoding is calculated according to $x=s\mathbf p+u\mathbf r$ and the residual formula. The task observation already includes $(m,p,s,u)$, so the two states have the same normal form, which is equivalent to the fact that they only differ in $f$; the proposition ~ \ref{hgov:prop:jointquotient} proves that this relationship is compatible with all registered actions and costs. Proposition ~ \ref{hgov:prop:jointgraph} gives the finite execution graph, smooth implementation, quotient chain mapping and change bridging, completing U02--U04 respectively. This example uses an explicit coordinate normal form; the finite test saturation theorem applies to task subdomains with explicit finite tables.

The state field of the old and new law objects is $X_m$. Migration takes $T=\mathrm{id}_{X_m}$, the old action tag corresponds to include, observation decoding and receiving readback are identical, and the chain mapping of the repair complex takes $i$ that defines ~ \ref{hgov:def:joint}. The identity and target of the two phases are the same, so $T(G_m)=G_m$; the predecessor, successor and zero cost of the old action are the same one by one. The transportation of generated atoms is verified step by step by the proposition ~ \ref{hgov:prop:atomic-transport}. The old unreachable certificate retains the old action family and the old complex version, and the new filled certificate saves $A^+(0,m)^{\mathsf T}=m\mathbf r$ and the actual $m$ step path. These interfaces are all explicitly proved under registered input conditions, and the set of additional open premises of U09 is feasible with the empty set. Its old family remains $u=0$ and the new family reaches $u=m$, so strict inclusion is also clearly witnessed.\end{proof}


\begin{corollary}
[Strict ability extension of fixed operation model to mutable law objects]\label{hgov:cor:value}
Fixed a task, input, acceptance predicate and allowed extended specification. If the actual control set of the old model is included in the new model, and there is a target with an unreachable certificate in the old model and a legal finite hit certificate in the new model, the task reachability of the new model is strictly greater than that of the old model. This gain can be achieved while maintaining the specified task identity and necessary constraints.\end{corollary}
\begin{certificate}
[Common tasks and difference witnessing]Saves the same task keys, old unreachable invariants or detached functionals, new hit paths, constraint preservation matrices, and acceptance records for encoding and readback.\end{certificate}
\begin{proof}
The old path is still executable giving the reachable set inclusion; the old unreachability and the new hit give the difference set non-empty: \[
\mathcal R_{\rm old}\subseteq\mathcal R_{\rm new},
\qquad
y_*\notin\mathcal R_{\rm old}\land y_*\in\mathcal R_{\rm new}
\Rightarrow
\mathcal R_{\rm old}\subsetneq\mathcal R_{\rm new}.
\] task identity and constraints are passed by proven holds on the same path. Theorem \ref{hgov:thm:joint} provides a construction witness for this conditional proposition on all integer parameters.\end{proof}


\begin{remark}
[Mathematical Objects and Application Interfaces for Conclusions]This book compares against explicit substrates, representations, and families of actually allowed operations. Geometric transport, chain and cochain, homotopy filling, task quotient and normal form together form the management program through typed mapping. Practical applications instantiate their antecedents with data, executable operators, measurements and resources; new capabilities are measured by reachable difference sets and acceptance certificates under the same task. Theoretical constructs, computational implementations, and real-life readbacks therefore maintain their own verifiable interfaces.\end{remark}


\clearpage \appendix \titleformat {\chapter} [display] {\normalfont\huge\bfseries} { Appendix \  \thechapter } {12pt} {\Huge} \addtocontents {toc} { \protect \renewcommand {\protect\cftchappresnum} { Appendix \space } } \addtocontents {toc} {\protect\setlength{\protect\cftchapnumwidth}{6.6em}}

\chapter{Unified Terminology and Symbol Index}
\label{app:hgov:notation}


The following table collects the main terms and notations in common mathematical substrate, geometry, repair, task quotient and governance in the order of argument in the text. When the same letter is reused in different local models, its source target, frequency, domain and version shall be subject to the corresponding definition; the brief description in the index retains the conditions required by the text, and the specific derivation can be found in the listed position.

\begingroup
\small
\setstretch{1.0}\begin{longtable}
{P{0.24\textwidth}P{0.54\textwidth}P{0.14\textwidth}}\toprule
term or notation&
Responsibilities and applicable boundaries&
Text positioning\\
\midrule
\endfirsthead
\toprule
term or notation&
Responsibilities and Boundaries of Applicability (continued)&
Text positioning\\
\midrule
\endhead
\midrule
\multicolumn {3} {r} { Continued on next page }\\
\endfoot
\bottomrule
\endlastfoot
$\Gamma,\mathcal C$&
Prerequisite environment and finite certificate; save input, conclusion, witness, parent dependency, valid domain and version, and the conditional reliability of the rules is inductively maintained by the finite certificate graph.&
Definition ~ \ref{hgov:def:language}, Theorem ~ \ref{hgov:thm:atomsound}\\
$\alpha_0,\ldots,\alpha_6$&
The registration atom of semantic differential; registers type, action antecedent, successor, observation at both ends, difference and cochain pairing in sequence.&
Definition ~ \ref{hgov:def:atom}\\
$L,\mathfrak S_L,\mathcal E$&
With task law object and common mathematical substrate; states, actions, partial successors, observations, identities, goals and costs are registered with the law object.&
Definition ~ \ref{hgov:def:law}\\
GZ-GMS / $\mathcal G$&
Collectively organize objects, laws, admissible flows, property layers, triple-layer connections and law curvature; ordinary smooth bundles are configured in designated structural sectors.&
Chapter ~ \ref{exp:ch:law}, \ref{exp:ch:geometry}\\
$\ell_n,\mathcal R_n,J_b$&
Jacobiator for higher-order operations of degree $2-n$, ternary-wise consistent residuals, and binary brackets; complete closure requires all $\mathcal R_n=0$ .&
Definition ~ \ref{hgov:linfty:convention}, Theorem ~ \ref{hgov:linfty:model}\\
$C_\bullet,\partial$&
The chain complex generated by task vertices, action edges and authenticated relationship surfaces; the common start and end points of the relationship surfaces are given as $\partial_1\partial_2=0$.&
Definition ~ \ref{hgov:def:taskcomplex}, Theorem ~ \ref{hgov:thm:complex}\\
$C^\bullet,d,\int_c\omega$&
Dual cochains, semantic differentials, pair with finite paths; closed cochains maintain integrals between paths connected by authenticated boundaries.&
Definition ~ \ref{hgov:def:semanticdiff}, Lemma ~ \ref{hgov:lem:semanticdiff}\\

$q,q_\bullet,q^*$&
Status quotient, separately registered chain implementations, and cochain pullbacks; observation factorization, action availability, successor exchange, and fee requirements are verified separately.&
Definition ~ \ref{hgov:def:quotient}, Theorem ~ \ref{hgov:thm:quotient}\\
$g_\alpha,\nabla^{g_\alpha},J$&
Sector metric, Levi--Civita connection and Jacobi variation field; curvature sign, initial value and comparison conditions jointly determine the growth conclusion.&
Definition ~ \ref{hgov:geo:def:sector}, \ref{hgov:geo:def:jacobi}\\
finite observation difference and error budget&
Differentiable conditions are obtained from observed responses, variational amplitudes, and unified residual bounds, and the common valid domain and actual reading accuracy are preserved.&
Definition ~ \ref{hgov:geo:def:observation}, proposition ~ \ref{exp:geo:prop:observation-robustness}\\
$d_\nabla,F_\nabla$&
Covariant external derivatives and curvature; register $d_\nabla^2\xi=F_\nabla\wedge\xi$ and Bianchi identities, and the complex condition of the flat subbundle will be verified separately.&
Definition ~ \ref{hgov:geo:def:exterior}, Theorem ~ \ref{hgov:geo:thm:curvature}\\
Heterogeneous plexus mapping and parallel transport&
Base and bundle maps are jointly registered; the commutation relation of connection, curvature and transport along the path holds on compatible types and valid domains.&
Definition ~ \ref{hgov:geo:def:heterogeneous}, Theorem ~ \ref{hgov:geo:thm:transport}\\
$I_k,\delta$&
Integration map from scalar form to smooth singular cochain; here $\delta$ is the singular cochain differential, satisfying $\delta I_k=I_{k+1}d$.&
Definition ~ \ref{hgov:geo:def:singular}, Theorem ~ \ref{hgov:geo:thm:integration}\\
Curvature characteristics cochain and transfer form&
Obtain the closed form and integral cochain from the invariant curvature polynomial, and compare the corresponding cohomology classes using the transfer formula of the connection deformation.&
Proposition ~ \ref{hgov:geo:prop:curvature-cocycle}, \ref{exp:geo:prop:transgression}\\
$\operatorname{Hom}^r,\delta_{ij}$&
The graded mapping space and Hom differential with source target; $\delta_{ij}f=d_jf-(-1)^rfd_i$, the corresponding graded symbol is compositely saved.&
Definition ~ \ref{hgov:def-hom}, Lemma ~ \ref{hgov:lem-hom}\\
continuation defects and homotopy filling&
The continuation is mapped into the common mathematical substrate to form a closed defect; the filling equation and its obstacle class are determined in the declared Hom complex.&
Definition ~ \ref{hgov:def-defect}, Theorem ~ \ref{hgov:thm-filler}\\
$U\xrightarrow{L}V\xrightarrow{D}W$&
finite linear repair rendering;
$DL=0$, $D\omega=0$, the actual filling solution is $Lu=\omega$, the homology quotient is $\ker D/\operatorname{im}L$.&
Definition ~ \ref{hgov:def-linear}\\
$P_G,r,\lambda$&
Orthogonal projection, residuals and separation functionals; under the positive definite Gram matrix, the residual being zero is equivalent to the disappearance of the obstacle class, and the dual certificate of the unsolvable branch is given.&
Theorem ~ \ref{hgov:thm-projection}\\
Allowed direction and minimum compatible expansion&
The new operation must preserve the old interface; calculate a lower bound for the new independent directions under given compatibility constraints, and register the actual expansion to reach the lower bound.&
Definition ~ \ref{hgov:def-extension}, Theorem ~ \ref{hgov:thm-minimal-extension}\\
Executable fixes and discontinued rankings&
Implement linear filling as an action with precondition, resources and cost; finite termination conclusion binds admissible execution and ranking conditions.&
Definition ~ \ref{hgov:def-executable}, Theorem ~ \ref{hgov:thm-execution}\\
$D^{(n)}(t),K_n,O_n$&
Formal parameter truncation, order-by-order repair operators and closing obstacles; the expansion requirement of the $n$ step is $\delta K_n=-O_n$, and the closing conditions of previous steps are saved.&
Definition ~ \ref{hgov:def-end}, Theorem ~ \ref{hgov:thm-total-tower}\\
$Q,Q^2$&
The coderivation of higher-order operations after suspension and its total square; the element-wise projection continuation of the disappearance of $Q^2=0$ and all higher-order compatibility residuals.&
Chapter ~ \ref{exp:ch:higher}, \ref{exp:ch:tower}\\
$C^-,C^+,\mathbf p,\mathbf r$&
The old and new repair complexes in the joint instance and the specified direction; the same hyperbolic geometry reads into the non-zero old barrier and the new filled direction.&
Definition ~ \ref{hgov:def:joint}, Lemma ~ \ref{hgov:lem:jointobstruction}\\
$m,\zeta_m,\Phi$&
Integer parameters, residuals, and potential functions of the joint task; legal unit padding hits in $m$ steps and reaches the lower bound on the number of actions for the task.&
Theorem ~ \ref{hgov:thm:joint}\\
$C_{\rm total}$&
Complete delivery cost; register geometry, input, search, execution, validation, encoding and readback separately, and the execution cost lower bound and other sub-items are measured separately.&
Inference ~ \ref{hgov:cor:jointcost}\\
$N,D,\mathcal K$&
Task encoding, decoder and capability package; $DN=O$ gives a readback specifying the observation, and the normal form also requires precise identification of the declared dynamic equivalence class.
&
Definition ~ \ref{hgov:def:normalform}, Theorem ~ \ref{hgov:thm:nf}\\
$\widehat F_\ell,\bot$&
Tagged generalization and undefined status of partial successors; observations, availability, and costs retain typed tags and decidable tables respectively.&
Definition ~ \ref{hgov:def:saturation}\\
$\mathcal P_n,\sim_n,q_*$&
Finite test families, successive partitioning, and stable quotients; under explicit finite table conditions, stable quotients preserve the initial observations and legal charges after all finite words.&
Definition ~ \ref{hgov:def:saturation}, Theorem ~ \ref{hgov:thm:saturation}\\
$T_n,j_n,\tau_n$&
State migration, chain mapping and action correspondence in law object update; identity, chainability, action antecedent, old target and running cost are saved separately.&
Definition ~ \ref{hgov:def:migration}\\
$D_n,r_n$&
Decoding of migration observations and readback of newly accepted output; readback is only called on registered and verified continuation domains.&
Definition ~ \ref{hgov:def:migration}, Theorem ~ \ref{exp:gov:migration-composition}\\
$\Gamma_n,\Delta_n,\Gamma_*$&
Stage premise ledger, new unresolved premise and joint premise set; finite migration maintains condition conclusion and version source through premise union.&
Theorem ~ \ref{hgov:thm:migration}, \ref{hgov:thm:main}\\
$\mathcal R_{\rm old},\mathcal R_{\rm new}$&
The reachable set under the same task and constraints; the operation inclusion gives the reachability set inclusion, and strict expansion requires the common witness of the old unreachability and the new hit.&
Inference ~ \ref{hgov:cor:value}\\
U01--U09&
Joint acceptance entry; covering type, complex, exchange, geometry, obstacle, repair, execution, readback and version transport, the source data of each entry is compatible with the valid domain.&
Theorem ~ \ref{hgov:thm:main} and its certificate\\
\end{longtable}
\endgroup

\chapter{Index of Literature Use and Textual Interfaces}
\label{app:hgov:coverage}


The following table lists the basic background, directly used structures and their main text interfaces according to the numbering order of the thirteen existing references. Various assumptions, specific constructions and proofs are subject to the corresponding chapters; complete bibliographic information, open source and local source manuscript location can be found in the references.

\small\begin{longtable}
{P{0.07\textwidth}P{0.21\textwidth}P{0.22\textwidth}P{0.41\textwidth}}\toprule
Ref.&
Topic&
book location&
Scope of use and text interface\\
\midrule
\endfirsthead
\toprule
Ref.&
Topic&
book location&
Scope of use and text interface (continued)\\
\midrule
\endhead
\midrule
\multicolumn {4} {r} { Continued on next page }\\
\endfoot
\bottomrule
\endlastfoot
\cite{S10}&
Homotopy identity, property sets and certificate hierarchies&
Read and verify conventions&
As the basic background registration for identity, property layers and certificated objects; the object types and atomic rules of this book are specifically given in Chapters ~ \ref{exp:ch:logic} and \ref{exp:ch:law}.\\
\cite{GFrameworkSynthesis}&
G-Framework unified structure and task interface&
Chapter ~ \ref{exp:ch:logic}, \ref{exp:ch:law}, \ref{exp:ch:task}, \ref{exp:ch:geometry}, \ref{exp:ch:normalform}, \ref{exp:ch:governance}&
Comprehensive background of law-space, triple-layer connection, task-sufficient quotient, normal form and complete cost; respectively access the semantic differential, observation retention and version migration of this book.\\
\cite{S01}&
Pure mathematics, law-space and GZ-GMS&
Chapter ~ \ref{exp:ch:law}, \ref{exp:ch:geometry}, \ref{exp:ch:hom}, \ref{exp:ch:tower}&
Common mathematical substrate, triple-layer connection and homotopy expansion background of objects and laws; specifically implemented by defining ~ \ref{hgov:def:law} with task fields and subsequent parting geometry.\\
\cite{S38}&
Rates and Quotients, strict sectors and task preservation&
Chapter ~ \ref{exp:ch:higher}, \ref{exp:ch:task}, \ref{exp:ch:hom}, \ref{exp:ch:tower}, \ref{exp:ch:joint} -- \ref{exp:ch:governance}&
Strict structure and complete high-order rendering, dynamic task quotient, commutation with differentials, prismatic witness, regular readback, cost and certificate migration.\\
\cite{ExtHatcher}&
Algebraic topology, singular chains and chain homotopy&
Chapter ~ \ref{exp:ch:task}, \ref{exp:ch:transport}, \ref{exp:ch:joint}, \ref{exp:ch:governance}&
Chain boundaries, smooth singular chains, prismatic identities, and induced homology; preserve explicit chain mappings and homotopy across joint instances and version migrations.\\

\cite{ExtWeibel}&
Homology algebra and Hom complex&
Chapter ~ \ref{exp:ch:hom}, \ref{exp:ch:extension}, \ref{exp:ch:governance}&
Graded Hom differential, kernel-image quotient, relative complex and exact sequence; access to continuation defects, finite fixes and homology version shipping.\\
\cite{ExtStacks}&
chain complex, cochain complex and homotopy&
Chapter ~ \ref{exp:ch:hom}, \ref{exp:ch:extension}&
Type conventions for complexes, chain mapping, and homotopy; supports Hom continuation, repair-change complexes, and structural verification of relative complexes.\\
\cite{ExtLadaMarkl}&
Strong homotopy Lie algebra&
Chapter ~ \ref{exp:ch:higher}, \ref{exp:ch:tower}&
The order of higher-order operations, shuffling, and Koszul notation; establishing an aryontical correspondence with the total square of the suspension coderivation.\\
\cite{ExtKobayashiNomizu}&
principal bundle, connection and curvature&
Chapter ~ \ref{exp:ch:geometry}, \ref{exp:ch:transport}&
Partially smooth substrate, covariant external derivatives, connection pullback, curvature exchange and parallel transport; used for geometric compatibility between heterogeneous bundles.\\
\cite{ExtLeeRiemann}&
Riemannian geometry compared with Jacobi&
Geometry input for Chapter ~ \ref{exp:ch:geometry}; Chapter ~ \ref{exp:ch:joint}&
Curvature convention, geodesic variation and hyperbolic space model; based on this, the text derives error conditions for growth, finite angle separation and task observation.\\
\cite{ExtBottTu}&
Differential Forms and Algebraic Topology Interface&
Chapter ~ \ref{exp:ch:transport}&
Transfer formulas for Stokes exchanges, integral cochains, flat subbundles, characteristic forms of curvature, and connection deformations.\\
\cite{S37}&
Basic disorder spectrum and certificate-based repair&
Chapter ~ \ref{exp:ch:hom} and subsequent repair interfaces&
The background of obstacle classes, homotopy filling and finite certificate; this book expands allowed expansion and order-by-order closure in chapters ~ \ref{exp:ch:extension} and \ref{exp:ch:tower} respectively.\\
\cite{S19}&
connection transmission, path encoding and normal form&
Chapter ~ \ref{exp:ch:normalform}&
Background on continuations of task encoding, normalization, and path representation; the book gives decoding fiber criteria, finite successor saturation, and executable capability reuse.\\
\end{longtable}
\normalsize

\clearpage
\begin{thebibliography} {99}
\addcontentsline {toc} {chapter} {\bibname}
\small
\bibitem {S10} GaoZheng. Homotopy Identity, Property Sets, and Monad Set Theory: Certified Foundations for Layered Discrete and Continuous Continua (Working Notes: Formal-System Version). 2026-01-12.
\newblock Included in Chinese Working Notes v5.2; collection DOI: \url{https://doi.org/10.5281/zenodo.22739417}.
\newblock Local source manuscript: \path{docs/project_docs/notebook_tex_v5.2/20260112_gframework_homotopy_identity_monadic_set_theory_foundation_formal_system.tex}, lines 1--1132.
\bibitem {GFrameworkSynthesis} GaoZheng. Unified Mathematical Structures, Computational Paradigms, and Philosophical Synthesis of the G-Framework: From Pan-Logic Analysis, Law-Space, and Homotopy Repair to Task-Sufficient Quotients and Rate--Quotient Unification. Basic version, 2026-11-01.
\newblock Local source manuscript: \path{docs/G_Framework/tex1/G_Framework_Unified_Synthesis_ZH.tex}; relevant chapters concern law-space, the triple-layer connection, macro-micro total homology, task-sufficient quotients, Rates and Quotients, and full costs.
\bibitem {S01} GaoZheng. Meta-Mathematical Theory based on Pan-Logic Analysis and Pan-Iterative Analysis (GaoZheng G-Framework) and the Principal-Bundle-Based Generalized Noncommutative Lie algebra (GaoZheng G-Algebra): An Integrated Construction. 2025.
\newblock Zenodo, cited version v1.0; the local source file is marked Version 1.1, 2025.
\newblock DOI: \url{https://doi.org/10.5281/zenodo.17651584}.
\newblock Local source manuscript: \path{docs/project_docs/Pub_GFramework_PureMath.tex}, lines 1--15562.
\bibitem {S38} GaoZheng. Rates and Quotients: A Formal System of Homotopy Mappings between G-Framework Semantic Calculus and Traditional Calculus. 2026-09-12.
\newblock Included in Chinese Working Notes v5.2; collection DOI: \url{https://doi.org/10.5281/zenodo.22739417}.
\newblock Local source manuscript: \path{docs/project_docs/notebook_tex_v5.2/20260912_gframework_rate_quotient_semantic_calculus_traditional_calculus_homotopy_mapping_formal_system.tex}, lines 1--3934.
\bibitem {ExtHatcher} Allen Hatcher. \emph {Algebraic Topology}. Cambridge University Press, 2002.
\newblock Author's publicly available full text: \url{https://pi.math.cornell.edu/~hatcher/AT/AT.pdf}.
\newblock Material used: singular chains, prism operators, chain homotopies, and induced homology.
\bibitem {ExtWeibel} Charles A. Weibel. \emph {An Introduction to Homological Algebra}. Cambridge Studies in Advanced Mathematics 38. Cambridge University Press, 1994.
\newblock DOI: \url{https://doi.org/10.1017/CBO9781139644136}.
\newblock Material used: complexes, Hom differentials, kernel-image quotients, homotopies, and graded compatibility.

\bibitem {ExtStacks} The Stacks Project Authors. \emph {The Stacks Project}, Section 12.13: Complexes, Tag 010V.
\newblock \url{https://stacks.math.columbia.edu/tag/010V}.
\newblock Material used: typed definitions of chain complexes, cochain complexes, homotopies, and homotopy equivalences.
\bibitem {ExtLadaMarkl} Tom Lada and Martin Markl. Strongly homotopy Lie algebras.
\emph {Communications in Algebra}, 23(6):2147--2161, 1995.
\newblock DOI: \url{https://doi.org/10.1080/00927879508825335}.
\newblock Public preprint: \url{https://arxiv.org/abs/hep-th/9406095}.
\bibitem {ExtKobayashiNomizu} Shoshichi Kobayashi and Katsumi Nomizu. \emph {Foundations of Differential Geometry, Volume 1}. Wiley Classics Library edition, 1996.
\newblock ISBN: 978-0-471-15733-5.
\newblock Publisher: \url{https://www.wiley.com/en-us/foundations-of-differential-geometry-volume-1-p-9780471157335}.
\newblock Material used: principal bundles, vector bundles, connections, curvature, and parallel transport.
\bibitem {ExtLeeRiemann} John M. Lee. \emph {Introduction to Riemannian Manifolds}. Second edition, Graduate Texts in Mathematics 176. Springer, 2018.
\newblock DOI: \url{https://doi.org/10.1007/978-3-319-91755-9}.
\newblock Material used: the Levi--Civita connection, curvature, Jacobi fields, and comparison.
\bibitem {ExtBottTu} Raoul Bott and Loring W. Tu. \emph {Differential Forms in Algebraic Topology}. Graduate Texts in Mathematics 82. Springer, 1982.
\newblock DOI: \url{https://doi.org/10.1007/978-1-4757-3951-0}.
\newblock Material used: differential forms, integration cochains, de Rham interfaces, and characteristic curvature forms.
\bibitem {S37} GaoZheng. A Genealogy of Foundational Obstructions in Mathematics, Physics, Computation, and Cognition within the G-Framework: Unified Law-Space, Total-Homology Semantic Fields, and Certified Resolution (Working Notes: Formal-System Version). 2026-07-08.
\newblock Included in Chinese Working Notes v5.2; collection DOI: \url{https://doi.org/10.5281/zenodo.22739417}.
\newblock Local source manuscript: \path{docs/project_docs/notebook_tex_v5.2/20260708_gframework_foundational_obstruction_spectrum_formal_system.tex}, lines 1--10126.
\bibitem {S19} GaoZheng. Set Equivalence of Bundle Spaces with Principal-Bundle Connections, Support Classes of Path Integrals, and Homotopy-Twisted Normal Forms (Working Notes: Formal-System Version). 2026-03-09.
\newblock Included in Chinese Working Notes v5.2; collection DOI: \url{https://doi.org/10.5281/zenodo.22739417}.
\newblock Local source manuscript: \path{docs/project_docs/notebook_tex_v5.2/20260309_gframework_principal_connection_fiber_space_path_integral_homotopy_twist_equivalence_formal_system.tex}, lines 1--5791.
\end{thebibliography}
\clearpage

\end{CJK*}
\end{document}
